\documentclass[11pt]{article}

\usepackage{geometry}
\geometry{margin=1in}

\usepackage{amsmath,amssymb,amsfonts,amsthm}
\usepackage{mathtools}
\usepackage{hyperref}
\usepackage{graphicx}
\usepackage{float}
\usepackage{booktabs}
\usepackage{tikz-cd}
\usepackage{cleveref}
\usepackage{subcaption}
\usepackage{placeins}
\usepackage{enumitem}
\usepackage{array}
\usepackage{xr}

\hypersetup{colorlinks=true, linkcolor=blue, citecolor=blue, urlcolor=blue}

%---------------------------
% Theorem environments (renumbered with S prefix-like scheme)
%---------------------------
\numberwithin{equation}{section}

\newtheorem{theorem}{Theorem}[section]
\newtheorem{lemma}[theorem]{Lemma}
\newtheorem{proposition}[theorem]{Proposition}
\newtheorem{corollary}[theorem]{Corollary}
\theoremstyle{definition}
\newtheorem{definition}[theorem]{Definition}
\newtheorem{assumption}[theorem]{Assumption}
\theoremstyle{remark}
\newtheorem{remark}[theorem]{Remark}

\newcommand{\proofref}[1]{%
  \par\noindent\emph{Proof.} See Section~\ref{#1}.\par\medskip
}
% In the supplement, proofs of body results are housed in this document, not in the
% main paper appendix; keep the \proofappendixref name for compatibility, but route
% it to a local \ref since the supplement contains the appendix itself.
\newcommand{\proofappendixref}[1]{\proofref{#1}}


%---------------------------
% Notation shortcuts
%---------------------------
\newcommand{\K}{\mathcal{K}}
\newcommand{\Hil}{\mathcal{H}}
\newcommand{\bd}{\partial}
\newcommand{\Id}{\mathbf{1}}
\newcommand{\Proj}{\mathbf{P}}
\newcommand{\spec}{\mathrm{spec}}
\newcommand{\dist}{\mathrm{dist}}
\newcommand{\rank}{\mathrm{rank}}
\newcommand{\tr}{\mathrm{tr}}
\newcommand{\Ran}{\mathrm{Ran}}
\newcommand{\opnorm}[1]{\left\lVert#1\right\rVert}
\newcommand{\e}{\varepsilon}
\newcommand{\sgn}{\mathrm{sgn}}
\newcommand{\Indint}{\operatorname{Ind}_{\mathrm{int}}}
\newcommand{\Ch}{\operatorname{Ch}}
\newcommand{\hexagon}{\mathrm{hex}}

%---------------------------
% Example name macros
%---------------------------
\newcommand{\ExOneName}{Disk-like patch (1 boundary)}
\newcommand{\ExTwoName}{Annulus-like band (2 boundaries)}
\newcommand{\ExThreeName}{Three-boundary patch (disk with two holes)}

%---------------------------
% Figure safety helper
%---------------------------
\newcommand{\safeincludegraphics}[2][]{%
  \IfFileExists{#2}{\includegraphics[#1]{#2}}{%
    \fbox{\begin{minipage}[c][0.18\textheight][c]{0.92\linewidth}
      \centering Missing figure file\\[0.3em]\texttt{\detokenize{#2}}
    \end{minipage}}%
  }%
}

\externaldocument[M-]{paper_main}

\title{Supplementary Material for: Boundary and Interface Classes for Universal Chern Hamiltonians\\
via Doubling and Jamming}

\author{
Yitzchak Shmalo\thanks{Corresponding author: \texttt{yitzchak.shmalo@gmail.com}.}\\[0.6em]
Einstein Institute of Mathematics\\
Hebrew University of Jerusalem
}
\date{\today}

\begin{document}
\maketitle

\begin{abstract}
This document collects the secondary results, full appendix proofs, all robustness variants,
hexagonal and square benchmark symbols, the symbolic verification appendix, and the full numerical
record for the main paper ``Boundary and Interface Classes for Universal Chern Hamiltonians via
Doubling and Jamming''. The main results, abstract, introduction, and headline theorems are stated
in the main paper; cross-references to the main paper use the \texttt{M-} prefix. The structure of
the supplement is: full proofs of the finite-\(M\) Schur-complement estimates
(\Cref{sec:schur_proofs}); the sign ledger (\Cref{app:sign_ledger}); two-band hexagonal and square
benchmark symbols and their proofs (\Cref{app:benchmark_hex_square}); the robustness variants for
finite defects, finite corners, tail-split and annular families
(\Cref{app:robustness_finite_defects,app:robustness_finite_corners,app:robustness_tail_annular});
the Section~5 variants (\Cref{app:section5_variants}); the Roe-algebra interface construction
variants (\Cref{sec:bulkboundary,app:roe_interface_variants}); the symbolic verification appendix
(\Cref{app:symbolic_verification}); the detailed appendix proofs of every body result
(\Cref{app:proofs,app:pathb_proofs}); the full numerical experiments
(\Cref{sec:numerics}); and the reproducibility appendix
(\Cref{app:reproducibility_appendix,app:minimal_reproduction_script}). This supplement and the
companion code are distributed through the project GitHub repository
(\href{https://github.com/yspennstate/Doubling-and-Jamming-A-Universal-Chern-Hamiltonian-on-Triangulated-Surfaces-with-Boundary}{\texttt{github.com/yspennstate/Doubling-and-Jamming-A-Universal-Chern-Hamiltonian-on-Triangulated-Surfaces-with-Boundary}}).
\end{abstract}


%====================================================================
\section{Full proofs of finite-dimensional Schur-complement estimates}\label{sec:schur_proofs}
%====================================================================

This section collects the full proofs of the finite-dimensional Schur-complement estimates whose
statements are recorded in the main paper, Section~\ref{M-sec:decoupling}. The proofs themselves
are housed in the detailed-proofs section, \Cref{app:proofs,app:proof:lem:schur}, and are reached
through the cross-references below. Each result of the main paper has a dedicated proof subsection
in \Cref{app:proofs} matching the appendix-label convention used throughout the supplement.
\begin{itemize}[leftmargin=2em]
\item Schur-complement resolvent formula (\Cref{M-lem:schur}): \Cref{app:proof:lem:schur}.
\item Compressed resolvent convergence (\Cref{M-thm:resolvent}): \Cref{app:proof:thm:resolvent}.
\item Boundary Hamiltonian as the large-jam limit (\Cref{M-cor:boundary_hamiltonian_limit}):
\Cref{app:proof:cor:boundary_hamiltonian_limit}.
\item Degenerate full block-resolvent limit (\Cref{M-cor:full_resolvent}):
\Cref{app:proof:cor:full_resolvent}.
\item First-order Schur expansion (\Cref{M-thm:first_order_eff}):
\Cref{app:proof:thm:first_order_eff}.
\item General contour-projection comparison (\Cref{M-thm:proj_general}):
\Cref{app:proof:thm:proj_general}.
\item Bounded-energy spectral localization (\Cref{M-thm:spectral_localization}):
\Cref{app:proof:thm:spectral_localization}; comparison with the relative boundary Hamiltonian
(\Cref{M-cor:relative_spectral_localization}): \Cref{app:proof:cor:relative_spectral_localization}.
\item First-order eigenvalue localization (\Cref{M-thm:first_order_spectral_localization}):
\Cref{app:proof:thm:first_order_spectral_localization}.
\item Copy-side leakage bound (\Cref{M-thm:eigvec_leak}): \Cref{app:proof:thm:eigvec_leak}.
\item Inertia, monotonicity, and no-copy-state bounds
(\Cref{M-thm:inertia_reduction,M-cor:scalar_jam_monotonicity,M-cor:no_copy_spectral_subspace}):
\Cref{app:proof:thm:inertia_reduction,app:proof:cor:scalar_jam_monotonicity,app:proof:cor:no_copy_spectral_subspace}.
\item Inherited gap and direct-sum jam estimates
(\Cref{M-prop:inherited_gap_A,M-cor:uniform_direct_sum_jam}):
\Cref{app:proof:prop:inherited_gap_A,app:proof:cor:uniform_direct_sum_jam}.
\item Internal bulk spectral equivalence (\Cref{M-thm:internal_bulk_spectral_equivalence}):
\Cref{app:proof:thm:internal_bulk_spectral_equivalence}.
\item Routine Schur-complement corollaries: \Cref{app:proof:schur_routine_corollaries}.
\end{itemize}
The proofs themselves now follow the appendix layout used throughout the rest of the supplement.

\section{Setup remarks}\label{sec:setup_remarks}

The fullness hypothesis on $\bd\K$ used throughout the main paper rules out two doubling
collisions:
\begin{enumerate}[label=(\roman*),leftmargin=2em]
\item \emph{Face collision.} If a face $\{v_0,v_1,v_2\}$ has all three vertices on $\bd\K$, its
  reflected copy has the same vertex set, so $D(\K)$ is not an abstract simplicial complex.
\item \emph{Edge collision and valence failure.} If an interior edge $[v_0,v_1]$ has both
  endpoints on $\bd\K$, the reflected copy has the same vertex set and cannot be distinguished
  from the original, and each face incident to $[v_0,v_1]$ in $\K$ produces a reflected face
  also incident to $[v_0,v_1]$, so the edge ends up incident to four or more faces in $D(\K)$.
  The simplex-count formula $|D(\K)_1|=2|\K_1|-|\bd\K_1|$ also fails.
\end{enumerate}
The fullness condition rules both out and holds automatically after one barycentric subdivision.
All refinement families used in the main paper satisfy the condition.

We also record here the closed-case cap-product convention used in \Cref{M-sec:setup}. For the
fundamental cycle $|X\rangle=\sum_{\tau\in\K_2}|\tau\rangle$, and the dual cochains
\(\delta_\xi\), the convention is
\(\delta_\xi\frown[v_0,\dots,v_2]=\delta_\xi([v_{2-p},\dots,v_2])\,[v_0,\dots,v_{2-p}]\),
which gives the Higson--Prodan operator $P$. The relation $BS=-SB^\dagger$ does \emph{not} by
itself imply $(B+B^\dagger)S+S(B+B^\dagger)=0$; this is why the body anticommutation statements
appear either as hypotheses or as measured finite-matrix defects.

% The closed-case gap criteria, doubling identities, and relative cap-product
% lemmas now have full statements (with labels) in the main paper. Only auxiliary
% companions and the bulk-spectral-equivalence corollaries remain here.

\begin{corollary}[Dimension of the doubled simplicial Hilbert space]\label{cor:dim_double}
\(\dim\Hil(D(\K))=2\dim\Hil(\K)-2|\bd\K_0|\).
\end{corollary}
\proofappendixref{app:proof:cor:dim_double}

\begin{proposition}[Euler characteristic]\label{prop:euler_double}
\(\chi(D(\K))=2\chi(\K)\).
\end{proposition}
\proofappendixref{app:proof:prop:euler_double}

\begin{corollary}[Genus formula]\label{cor:genus_double}
Let $\K$ be connected, orientable, genus $g$, with $b\ge 1$ boundary components. Then
$g_{\mathrm{double}}=2g+b-1$.
\end{corollary}
\proofappendixref{app:proof:cor:genus_double}

\begin{corollary}[Quotient equality of the relative and doubled boundary models]\label{cor:pl_doubled_quotient}
Let $(\K_n)_n$ be a uniformly regular refinement family. With
$H^{\mathrm{PL},\infty}_\pm=\bigoplus_n H^{\mathrm{PL}}_\pm(\K_n)$ and
$H^{\bd,\infty}_\pm=\bigoplus_n H^{\bd}_\pm(\K_n)$ on the original-side Hilbert space,
$H^{\bd,\infty}_\pm-H^{\mathrm{PL},\infty}_\pm\in C^*(Y\subset X)$, and the relative duality defect
$\bigoplus_n(B_{\K_n}S^{\mathrm{PL}}_{\K_n}+S^{\mathrm{PL}}_{\K_n}B_{\K_n}^\dagger)$ belongs to
$C^*(Y\subset X)$. Continuous functional calculus passes both models to the same image in
$\mathcal A/\mathcal I$.
\end{corollary}
\proofappendixref{app:proof:cor:pl_doubled_quotient}

\begin{corollary}[Asymptotic bulk spectral equivalence]\label{cor:pl_bulk_spectral_equivalence}
With $d_n=\dim\Hil(\K_n)$ and $r_n=\dim\ell^2(N_{R_{\mathrm{PL}}}(\bd\K_n))$, the spectral
distribution functions of $H^{\mathrm{PL}}_\pm(\K_n)$ and $H^\partial_\pm(\K_n)$ satisfy
$\sup_t|d_n^{-1}N^{\bd}_{\pm,n}(t)-d_n^{-1}N^{\mathrm{PL}}_{\pm,n}(t)|\le r_n/d_n$.
For families with $r_n/d_n\to0$ the normalized spectral distribution functions converge to the
same limits.
\end{corollary}
\proofappendixref{app:proof:cor:pl_bulk_spectral_equivalence}

%====================================================================
\section{Restatements of main-paper body results referenced in the appendix proofs}\label{sec:body_restatements}
%====================================================================

The appendix proofs below refer back to the body lemma/theorem statements of the main paper.
For convenience and to make the supplement self-contained, the body statements whose
\emph{statements} were compressed in the main paper (kept as informal restatements rather than as
labelled environments) are repeated here with their original labels.

\begin{definition}[Propagation / locality radius]\label{def:propagation}
A bounded operator $T$ on $\Hil(\K)=\ell^2(X(\K))$ has \emph{propagation at most $R$} if
$\langle \sigma,T\tau\rangle=0$ whenever $\dist_{X(\K)}(\sigma,\tau)>R$. We also call such an
operator $R$-local.
\end{definition}

\subsection*{Section~\protect\ref{M-sec:setup} body labels (closed-case and orientation conventions)}\label{sec:Sclosed}\label{sec:orientation_signs}

The labels \(\ref{sec:Sclosed}\) and \(\ref{sec:orientation_signs}\) are used by the appendix to
refer to the closed-case cap-product definition and the orientation/sign conventions; in this
supplement they point at this restatement subsection. The conventions themselves are stated in
\Cref{M-sec:setup} of the main paper.

\subsection*{Equation-label aliases used by the appendix proofs}
\label{eq:resolvent_bound}%
\label{eq:leak_crude}%
\label{eq:lifted_cap_defect}%
\label{eq:proj_general_resolvent_bound}%
\label{eq:relative_cap_chain_identity}%
\label{eq:specA_exact}%
\label{eq:specA_crude}%

The label keys
\(\ref{eq:resolvent_bound}\), \(\ref{eq:leak_crude}\), \(\ref{eq:lifted_cap_defect}\),
\(\ref{eq:proj_general_resolvent_bound}\), \(\ref{eq:relative_cap_chain_identity}\),
\(\ref{eq:specA_exact}\), \(\ref{eq:specA_crude}\) used by the supplement's appendix proofs are
defined here as section-local aliases. The equations they correspond to are the explicit
resolvent and leakage bounds collected in \Cref{M-thm:resolvent,M-thm:eigvec_leak} and in the
relative-boundary comparison \Cref{M-thm:relative_boundary_comparison}; the aliases let the
supplement's appendix proofs cite them with the original label keys without xr round-trips.

\subsection*{Numerical-table aliases}

The certificate tables (\Cref{tab:periodic_triangular_certificate,tab:periodic_triangular_extension_certificate})
are produced by the companion script and appear in this supplement as
\Cref{sec:numerics}; the body refers to them through these label aliases.

\begin{table}[H]
\centering\small
\caption{Periodic triangular Bloch certificate (FHS check; reproduced from the companion script).}
\label{tab:periodic_triangular_certificate}
\IfFileExists{periodic_triangular_certificate_table.tex}{\begin{tabular}{llrrrr}\toprule
Hamiltonian & Grid & FHS Chern & min $|E|$ & min link & max plaquette \\
\midrule
$H_+$ & $11\times11$ & -1.000000 & 0.550897 & 0.905137 & 0.265005 \\
$H_+$ & $21\times21$ & -1.000000 & 0.540655 & 0.970926 & 0.080163 \\
$H_+$ & $31\times31$ & -1.000000 & 0.540954 & 0.986311 & 0.037594 \\
$H_+$ & $51\times51$ & -1.000000 & 0.540291 & 0.994870 & 0.014056 \\
$H_-$ & $11\times11$ & 1.000000 & 0.550897 & 0.905137 & 0.265005 \\
$H_-$ & $21\times21$ & 1.000000 & 0.540655 & 0.970926 & 0.080163 \\
$H_-$ & $31\times31$ & 1.000000 & 0.540954 & 0.986311 & 0.037594 \\
$H_-$ & $51\times51$ & 1.000000 & 0.540291 & 0.994870 & 0.014056 \\
\bottomrule
\end{tabular}
}{%
\emph{(File \texttt{periodic\_triangular\_certificate\_table.tex} not found. Run
\texttt{python periodic\_triangular\_bulk\_certificate.py --write-table}.)}}
\end{table}

\begin{table}[H]
\centering\small
\caption{Periodic triangular extension checks (FHS).}
\label{tab:periodic_triangular_extension_certificate}
\IfFileExists{periodic_triangular_extension_certificate_table.tex}{%
\resizebox{\linewidth}{!}{\begin{tabular}{llrrrrr}\toprule
Diagnostic & Model & Grid & expected & FHS & min $|E|$ & min link \\
\midrule
positive mass & $H_{+,m}$, $m=0.5$ & $31\times31$ & -1 & -1.000000 & 0.287664 & 0.959784 \\
positive mass & $H_{-,m}$, $m=2$ & $31\times31$ & 1 & 1.000000 & 0.806343 & 0.993537 \\
finite cover & $H_{+,A}$, $A=\mathrm{diag}(2,1)$ & $31\times31$ & -2 & -2.000000 & 0.540954 & 0.945837 \\
finite cover & $H_{-,A}$, $A=\mathrm{diag}(2,1)$ & $31\times31$ & 2 & 2.000000 & 0.540954 & 0.945837 \\
orientation-reversing cover & $H_{+,A}$, $A=\mathrm{diag}(1,-1)$ & $31\times31$ & 1 & 1.000000 & 0.540954 & 0.986311 \\
cell enlargement & $H_+$, $2\times1$ folded cell & $21\times21$ & -1 & -1.000000 & 0.540655 & 0.965480 \\
cell enlargement & $H_-$, $2\times1$ folded cell & $21\times21$ & 1 & 1.000000 & 0.540655 & 0.965480 \\
\bottomrule
\end{tabular}
}}{%
\emph{(File \texttt{periodic\_triangular\_extension\_certificate\_table.tex} not found. Run
\texttt{python periodic\_triangular\_bulk\_certificate.py --write-table}.)}}
\end{table}

%====================================================================
\section{Auxiliary Roe-algebra lemmas, stability, and cross-checks}\label{sec:roe_auxiliaries}
%====================================================================

The following auxiliary Roe-algebra lemmas, stability statements, and cross-check lemmas are
referenced by the body proof of \Cref{M-thm:pathb_roe_winding} in the main paper. Their
statements live here in the supplement; the corresponding proof bodies appear in
\Cref{app:pathb_proofs} below.

\begin{lemma}[Local analytic properties of \(H_\pm(K)\)]\label{lem:pathb_hp_operator_bounds}
For a bounded-geometry triangulation \(K\), the operators \(H_\pm(K)\) are self-adjoint and have
finite hopping range in the barycentric graph metric. Along a bounded-geometry family the
hopping range and operator norms are bounded uniformly.
\end{lemma}
\proofappendixref{app:proof:lem:pathb_hp_operator_bounds}

\begin{lemma}[Combes--Thomas decay for the closed-bulk Fermi projection]\label{lem:pathb_combes_thomas}
Assume the uniform hopping and gap constants in
\Cref{lem:pathb_hp_operator_bounds,M-lem:pathb_hp_gap}. Then there are constants $C,c>0$,
independent of $n$, such that
\(\bigl|\langle\delta_v,P_{\pm,n}\delta_w\rangle\bigr|\le C e^{-c\,d(v,w)}\)
for all barycentric vertices $v,w\in X(D(\K_n))$.
\end{lemma}
\proofappendixref{app:proof:lem:pathb_combes_thomas}

\begin{lemma}[Boundary map for the interface projection]\label{lem:pathb_boundary_map}
The map $\partial_Y$ in \Cref{M-lem:pathb_exact_sequence} is the standard $K$-theory connecting
map for the Roe extension. For the lift $h(J_{\pm,M}^\infty)$ of $p_J$,
\(\partial_Y([p_J])=[u_{J,h}]\in K_1(C^*(Y\subset X))\), and the class is independent of the
admissible cutoff $h$.
\end{lemma}
\proofappendixref{app:proof:lem:pathb_boundary_map}

\begin{definition}[Ewert--Meyer pairing for the original/copy partition]\label{def:pathb_em_pairing}
The Ewert--Meyer bulk--interface pairing is the coarse Mayer--Vietoris pairing of a bulk
$K_0(\mathcal A/\mathcal I)$-class with the interface cocycle obtained from the partition
$(X^+,X^-;Y)$ \cite{EwertMeyer2019}. Here this pairing is evaluated on the class $[p_J]$ through
the connecting map $\partial_Y$ above. The compatibility assertion is not an extra hypothesis: it
is the statement that, for the doubled-jammed original/copy partition, Ewert--Meyer's coarse
boundary map is represented by the same Roe extension and the same cutoff unitary $u_{J,h}$.
\end{definition}

% Stability and cross-check lemmas are now stated in the main paper (Section M-sec:pathb_internalization).
% Their proof bodies remain below in \Cref{app:pathb_proofs}.

%====================================================================
\section{A coarse interface class construction}\label{sec:bulkboundary}
%====================================================================

We place the doubled-and-jammed refinement family in a Roe-algebra construction that produces the
interface class. In the headline application, the required closed-bulk gap is supplied by
Higson--Prodan's closed theorem. We use the
quotient extension
\[
0 \longrightarrow \mathcal I \longrightarrow \mathcal A \xrightarrow{\;\pi\;} \mathcal A/\mathcal I
\longrightarrow 0,
\qquad
\mathcal A:=C^*(X),
\qquad
\mathcal I:=C^*(Y\subset X),
\]
and define the interface class from a canonical quotient projection associated to a sidewise
realization.

Here \(C^*(X)\) denotes the scalar Roe algebra: the norm closure in
\(\mathcal B(\ell^2(X))\) of locally compact finite-propagation operators. We do not use the
uniform Roe algebra in this paper. In Theorem~\ref{M-thm:pathb_roe_winding} of
Section~\ref{M-sec:pathb_internalization}, the Ewert--Meyer Mayer--Vietoris boundary map for the
original/copy partition is identified internally with the Roe connecting map used below. The
HP--Ewert--Meyer terminology below fixes the convention for that pairing; it is not an additional
input for the original/copy headline theorem.

Uniform locality, boundedness, and a closed-bulk gap give the quotient projection and the
interface \(K_1\)-class. In the HP-covered case, Higson--Prodan supplies the gap and Chern value;
in the periodic triangular case, Section~\ref{M-sec:periodicbulk} supplies them directly. The
original/copy \(K_1\)-class identity is Theorem~\ref{M-thm:pathb_roe_winding}; the integer
winding evaluation is internal in the Toeplitz periodic case
(\Cref{M-cor:toeplitz_winding}) and conditional in general
(\Cref{M-cor:conditional_winding}).

In the present scalar-coefficient setting, the Roe algebra $C^*(X)$ is \emph{unital}: the identity
operator on $\ell^2(X)$ has propagation $0$ and is locally compact (since every finite-support
compression is finite rank), hence belongs to $C^*(X)$. Similarly, every diagonal
$\{0,1\}$-valued multiplication operator on $\ell^2(X)$ belongs to $C^*(X)$. The quotient
$\mathcal A/\mathcal I$ is therefore also unital, and the continuous functional calculus in both
algebras is the standard unital one.

This convention is the scalar, ample-representation version of the Roe algebra. Ewert--Meyer's
Mayer--Vietoris boundary map is used below only through the induced connecting map for the
coarsely excisive pair \((X^+,X^-)\) and its compatible pairing with the closed bulk Chern class.
Passing between this unital scalar presentation and stabilized or nonunital presentations changes
the bookkeeping, not the connecting class; Theorem~\ref{M-thm:pathb_roe_winding} identifies the
connecting class used here with the Ewert--Meyer convention for the original/copy partition.

\begin{remark}[Notation and conventions]
Throughout this section, whenever $T$ is a bounded self-adjoint operator on $\ell^2(X)$, we write
\[
\spec_{\mathrm{op}}(T):=\spec_{\mathcal B(\ell^2(X))}(T)
\]
for its operator spectrum. Because $C^*(X)\subset\mathcal B(\ell^2(X))$ is a unital inclusion,
the $C^*(X)$-spectrum and the operator spectrum of any $T\in C^*(X)$ coincide. All spectral-gap
statements for concrete operators refer to $\spec_{\mathrm{op}}$.
\end{remark}

We define the canonical quotient projection via \emph{continuous} cutoff functions that agree with
$\chi_{(-\infty,0)}$ on the relevant bounded spectral windows. This avoids applying the
discontinuous function $\chi_{(-\infty,0)}$ across the quotient map, since the spectrum of a
quotient-algebra element can differ from the operator spectrum of its lift. In the present unital
scalar Roe algebra this is ordinary unital functional calculus; in nonunital variants one should use
compactly supported cutoffs on the same bounded spectral windows.

\subsection{Uniformly regular refinement families}

Let $(\K_n)$ be a refinement sequence of a fixed compact oriented surface with boundary. For each
$n$, write
\[
X_n:=X\bigl(D(\K_n)\bigr)
\]
for the set of barycenters of all simplices in the doubled triangulation, equipped with the
barycentric graph metric from Definition~\ref{def:propagation}.

\begin{definition}[Uniformly regular refinement family]\label{def:uniform_regular}
The refinement sequence $(\K_n)$ is \emph{uniformly regular} if the metric spaces $X_n$ have
uniformly bounded geometry: for every $r>0$ there exists $N(r)<\infty$ such that
\[
|B_{X_n}(x,r)|\le N(r)
\qquad\text{for all }n\text{ and all }x\in X_n.
\]
A sufficient concrete condition is a uniform bound on simplex incidences under refinement.
\end{definition}

For each $n$, let $X_{n,\mathrm{orig}}$ be the subset of simplex barycenters lying on the original
side together with the glued boundary simplices, and let $X_{n,\mathrm{copy}}$ be the subset of
simplex barycenters strictly on the reflected side. Let $\Gamma_n\subset X_n$ denote the set of
simplex barycenters whose simplices meet the glued boundary subcomplex. If $R>0$, define the
$R$-collar
\[
Y_n^{(R)}:=N_R(\Gamma_n)\subset X_n.
\]

\begin{definition}[Coarse disjoint unions and the interface ideal]\label{def:coarse_union}
Let
\[
X:=\bigsqcup_{n=1}^\infty X_n
\]
be a coarse disjoint union equipped with a metric that restricts to the given metric on each
component and separates distinct components by distances tending to infinity; concretely, for every
$r>0$ there is $N_r$ such that $\dist(X_n,X_m)>r$ whenever $n\neq m$ and
\(\max\{n,m\}>N_r\). For fixed $R\ge0$, let
\[
Y^{(R)}:=\bigsqcup_{n=1}^\infty Y_n^{(R)}\subset X.
\]
The Roe algebra $C^*(X)$ is the norm closure of bounded finite-propagation operators on
$\ell^2(X)$ that are locally compact. The interface ideal $C^*(Y^{(R)}\subset X)$ is the norm
closure of bounded finite-propagation operators supported within a bounded neighborhood of
$Y^{(R)}$.
\end{definition}

\begin{lemma}[Roe ideal property]\label{lem:roe_ideal}
For every $R\ge0$, the algebra $C^*(Y^{(R)}\subset X)$ is a closed two-sided ideal in $C^*(X)$.
\end{lemma}
\proofappendixref{app:proof:lem:roe_ideal}




\begin{proposition}[Independence of fixed collar radius]\label{prop:collar_radius_independence}
Let
\[
\Gamma:=\bigsqcup_{n=1}^\infty \Gamma_n\subset X,
\qquad
Y^{(R)}:=\bigsqcup_{n=1}^\infty N_R(\Gamma_n).
\]
For any fixed $R,R'\ge0$,
\[
C^*(Y^{(R)}\subset X)=C^*(Y^{(R')}\subset X).
\]
Consequently the interface ideal used below is independent of the particular fixed collar radius.
\end{proposition}
\proofappendixref{app:proof:prop:collar_radius_independence}




\begin{proposition}[Coarse invariance of the interface ideal]\label{prop:interface_ideal_coarse_invariance}
Let \(Y_0,Y_1\subset X\) be subsets at finite Hausdorff distance: there is \(R<\infty\) such that
\[
Y_0\subset N_R(Y_1),
\qquad
Y_1\subset N_R(Y_0).
\]
Then
\[
C^*(Y_0\subset X)=C^*(Y_1\subset X).
\]
\end{proposition}
\proofappendixref{app:proof:prop:interface_ideal_coarse_invariance}




\begin{remark}
For background on coarse spaces and Roe algebras, see Roe \cite{RoeLectures}. Here the only
Roe-algebra input is the concrete quotient extension
\[
0\to C^*(Y\subset X)\to C^*(X)\to C^*(X)/C^*(Y\subset X)\to 0,
\]
together with the standard $K_1$-class of a unitary in $\Id+\mathcal I$.
\end{remark}

\subsection{Uniform locality and the bulk Fermi projection}

For each $n$, write
\[
H_{\pm,n}:=H_\pm(D(\K_n)).
\]

\begin{assumption}[Uniform cap-product locality input]\label{ass:S_uniform}
There exist constants $R_S,N_S,c_S>0$, independent of $n$, such that for every $n$ and every pair
of simplex barycenters $\sigma,\tau\in X_n$:
\begin{enumerate}[label=(\roman*),leftmargin=2em]
\item if $d_{X_n}(\sigma,\tau)>R_S$, then $S_n(\sigma,\tau)=0$;
\item each row and each column of the matrix of $S_n$ has at most $N_S$ nonzero entries;
\item $|S_n(\sigma,\tau)|\le c_S$.
\end{enumerate}
\end{assumption}

\begin{remark}
Assumption~\ref{ass:S_uniform} is the uniform bounded-geometry locality input used for the
Higson--Prodan cap-product operator on the closed doubled family. For the exact
Higson--Roe/Higson--Prodan duality operator on a uniformly regular closed family, this is an
imported consequence of geometric control: the cap-product operator is a geometrically controlled
map in the sense of Higson--Roe \cite{HigsonRoeII}, hence has uniformly bounded propagation and
uniformly bounded matrix coefficients. Equivalently, it may be kept as an explicit hypothesis for
abstract variants of the construction. Concretely, the required input is a uniform finite row and
column support bound and a uniform entry bound for the cap-product matrix. In the local
Higson--Prodan formula this comes from the fact that each term pairs only simplices in a uniformly
bounded barycentric neighborhood, with coefficients drawn from a fixed finite set of incidence
signs and phase factors; the lemma below records the resulting Schur-test estimate used here.
\end{remark}

\begin{lemma}[Uniform locality and norm bound for \(S_n\)]\label{lem:S_uniform_norm}
Assume that the Higson--Prodan cap-product operators $S_n$ satisfy
Assumption~\ref{ass:S_uniform}. Then
\[
\mathrm{prop}(S_n)\le R_S,
\qquad
\opnorm{S_n}\le N_Sc_S
\]
for all $n$.
\end{lemma}
\proofappendixref{app:proof:lem:S_uniform_norm}




\begin{proposition}[Uniform locality and norm bound]\label{prop:uniform_locality_norm}
Let $(\K_n)$ be a uniformly regular refinement family, and assume
\ref{ass:S_uniform}. Then there exist constants
$R_{\mathrm{loc}},B_{\mathrm{bulk}}>0$ such that for all $n$:
\begin{enumerate}[label=(\alph*),leftmargin=2em]
\item $H_{\pm,n}$ has propagation at most $R_{\mathrm{loc}}$ on $\ell^2(X_n)$;
\item $\opnorm{H_{\pm,n}}\le B_{\mathrm{bulk}}$.
\end{enumerate}
\end{proposition}
\proofappendixref{app:proof:prop:uniform_locality_norm}




\begin{remark}[Structure of the gap condition]\label{rem:gap_structure}
The operator $S_n$ is degree-reversing, so its $\Hil_0\leftrightarrow\Hil_2$ block is a
rectangular matrix of dimensions $|D(\K_n)_2|\times|D(\K_n)_0|$. On a closed surface with Euler
characteristic $\chi$, the relation $3|K_2|=2|K_1|$ gives $|K_0|-|K_2|=\chi-|K_2|/2$, so
generically $|K_0|\neq|K_2|$ and $S_n$ has $\bigl||K_0|-|K_2|\bigr|$ zero eigenvalues from the
dimension mismatch. Hence the sufficient condition $\dist(0,\spec(S_n))>0$ from
Proposition~\ref{M-prop:gap_from_S} is \emph{not} available for typical triangulations (it holds only
when $|K_0|=|K_2|$, e.g., the tetrahedron).

Nevertheless, $H_{\pm,n}$ can still be gapped because the term $(B_n+B_n^\dagger)$ may be nonzero
on $\ker(S_n)$. More precisely:
\[
\langle\psi, H_{\pm,n}^2\psi\rangle
=
\|(B_n+B_n^\dagger)\psi\|^2+\|S_n\psi\|^2
\pm\langle\psi,(B_n^\dagger S_n+S_nB_n)\psi\rangle
\]
for each $\psi$. On $\ker(S_n)$, $\|S_n\psi\|^2=0$ and the cross term vanishes as a quadratic form:
$\langle\psi,(B_n^\dagger S_n+S_nB_n)\psi\rangle=\langle B_n\psi,S_n\psi\rangle+\langle
S_n\psi,B_n\psi\rangle=0$ since $S_n\psi=0$. So the gap on $\ker(S_n)$ is controlled by
$(B_n+B_n^\dagger)^2$. On $\ker(S_n)^\perp$, the positive term $S_n^2$ also contributes.

The uniform gap therefore depends on a joint spectral property of $B_n+B_n^\dagger$ and $S_n$
together, not on $S_n$ alone. The Higson--Roe controlled-duality machinery supplies the relevant
locality and Poincar\'e-duality input for the closed model, and the HP-covered headline application
uses Higson--Prodan's clean-gap theorem for the closed doubles. The numerics
(Figure~\ref{fig:refine}) report the measured gap on the tested refinement levels, but they are not
used as a proof of a refinement-uniform gap.
\end{remark}

The uniform bulk hypothesis used by the coarse interface theorem is the following.

\begin{assumption}[Uniform bulk hypotheses]\label{ass:uniform_bulk}
The refinement family $(\K_n)$ is uniformly regular and there exist constants
$R_{\mathrm{loc}},B_{\mathrm{bulk}},g_{\mathrm{bulk}}>0$ such that for all $n$:
\begin{enumerate}[label=(\alph*),leftmargin=2em]
\item $H_{\pm,n}$ has propagation at most $R_{\mathrm{loc}}$ on $\ell^2(X_n)$;
\item $\opnorm{H_{\pm,n}}\le B_{\mathrm{bulk}}$;
\item $\dist\bigl(0,\spec_{\mathrm{op}}(H_{\pm,n})\bigr)\ge g_{\mathrm{bulk}}$.
\end{enumerate}
\end{assumption}

\begin{remark}\label{rem:assumption_status}
Parts~(a) and~(b) are reduced in Proposition~\ref{prop:uniform_locality_norm} to uniform regularity
together with the standard cap-product locality input recorded in
Assumption~\ref{ass:S_uniform}. Part~(c) is the spectral input not proved by the elementary locality
argument. As explained in Remark~\ref{rem:gap_structure}, this is a joint spectral condition on
$B_n+B_n^\dagger$ and $S_n$ that does not reduce to invertibility of $S_n$ alone. For the
HP-covered closed doubles used in the headline theorem, Higson--Prodan supplies~(c).
\end{remark}

\begin{assumption}[Higson--Prodan closed-bulk input]\label{ass:hp_admissibility}
The original refinement family $(\K_n)_n$ is uniformly regular, and the closed doubled refinement
family $L_n=D(\K_n)$ is covered by the Higson--Prodan closed universal-model theorem. Thus the
following data are available:
\begin{enumerate}[label=(\alph*),leftmargin=2em]
\item the Higson--Roe/Higson--Prodan controlled Poincar\'e-duality operator $S_n$ with the
uniform locality, row-support, and coefficient bounds of Assumption~\ref{ass:S_uniform};
\item constants $B_{\mathrm{HP}},g_{\mathrm{HP}}>0$, independent of $n$, such that
\[
\opnorm{H_{\pm,n}}\le B_{\mathrm{HP}},
\qquad
\dist\bigl(0,\spec_{\mathrm{op}}(H_{\pm,n})\bigr)\ge g_{\mathrm{HP}};
\]
\item the Higson--Prodan/Bellissard--Connes bulk Chern pairing $\Ch$ for which the associated
closed-bulk Fermi classes have Chern number $\pm1$.
\end{enumerate}
\end{assumption}

\begin{remark}\label{rem:hp_admissibility_status}
Assumption~\ref{ass:hp_admissibility} records the cited closed-bulk theorem used for the jammed
boundary model. The geometric input is the controlled Poincar\'e-duality construction of
Higson--Roe: bounded-geometry oriented combinatorial manifolds give analytically controlled
Hilbert--Poincar\'e complexes
\cite[Theorem~3.14 and Proposition~4.1]{HigsonRoeII}. Higson and Prodan specialize this
controlled construction to the tight-binding Hamiltonians $H_\pm=(B+B^\dagger)\pm S$ and identify
the closed-bulk Chern classes in their universal Chern model
\cite{HigsonProdanUniversalPRL,HigsonProdanUniversalArxiv}. The present paper cites that closed
theory exactly as it cites the standard background theorems; after the HP gap and Chern
conclusions are in place, the doubled-and-jammed interface conclusions below are proved internally.
\end{remark}

\begin{theorem}[Closed-bulk consequences of the Higson--Prodan input]\label{thm:hp_closed_input}\label{cor:hp_uniform_bulk}
If Assumption~\ref{ass:hp_admissibility} holds, then
Assumptions~\ref{ass:S_uniform} and~\ref{ass:uniform_bulk} hold. Moreover the bulk Fermi
projection
\[
P_\pm=\chi_{(-\infty,0)}(H_\pm^\infty),
\]
constructed below in Proposition~\ref{prop:bulk_class}, is the Higson--Prodan closed-bulk Fermi
projection for the refinement family of doubles, and its imported Chern pairing satisfies
\[
\Ch([P_\pm])=\pm1.
\]
\end{theorem}
\proofappendixref{app:proof:thm:hp_closed_input}




Let
\[
H_\pm^\infty:=\bigoplus_{n=1}^\infty H_{\pm,n}
\qquad\text{on}\qquad
\ell^2(X)=\bigoplus_{n=1}^\infty \ell^2(X_n).
\]

\begin{proposition}[Bulk Roe-algebra class]\label{prop:bulk_class}
Assume \ref{ass:uniform_bulk}. Then $H_\pm^\infty\in C^*(X)$ is bounded, self-adjoint, and
finite-propagation. Let
\[
P_\pm:=\chi_{(-\infty,0)}\!\bigl(H_\pm^\infty\bigr)
\]
denote the negative operator spectral projection of $H_\pm^\infty$ on $\ell^2(X)$. Then
$P_\pm\in C^*(X)$ and therefore defines a class
\[
[P_\pm]\in K_0\bigl(C^*(X)\bigr).
\]
\end{proposition}
\proofappendixref{app:proof:prop:bulk_class}




\begin{proposition}[Gap-homotopy invariance of the bulk Fermi class]\label{prop:bulk_gap_homotopy}
Let $s\mapsto H_s$, $s\in[0,1]$, be a norm-continuous path of self-adjoint elements of $C^*(X)$.
Assume that there are constants $B,g>0$ such that
\[
\spec_{\mathrm{op}}(H_s)\subset[-B,-g]\cup[g,B]
\qquad\text{for every }s\in[0,1].
\]
Set
\[
P_s:=\chi_{(-\infty,0)}(H_s).
\]
Then $s\mapsto P_s$ is a norm-continuous path of projections in $C^*(X)$. In particular,
\[
[P_s]\in K_0(C^*(X))
\]
is independent of $s$.
\end{proposition}
\proofappendixref{app:proof:prop:bulk_gap_homotopy}




\begin{proposition}[Quantitative stability of a spectral projection]\label{prop:quantitative_projection_stability}
Let $H=H^*$ be a bounded self-adjoint operator, let $V=V^*$ be bounded, and let $\Gamma$ be a
positively oriented contour with
\[
\delta_\Gamma:=\dist(\Gamma,\spec(H))>0.
\]
If
\[
\opnorm{V}<\delta_\Gamma,
\]
then $\Gamma\subset\rho(H+V)$ and
\[
\opnorm{P_\Gamma(H+V)-P_\Gamma(H)}
\le
\frac{|\Gamma|}{2\pi}\,
\frac{\opnorm{V}}{\delta_\Gamma(\delta_\Gamma-\opnorm{V})}.
\]
\end{proposition}
\proofappendixref{app:proof:prop:quantitative_projection_stability}




Small-norm Fermi and positive-side interface stability is a variant statement collected in Appendix~\ref{app:roe_interface_variants}.

\subsection{A collar-locality theorem}

For each $n$, let $P_{\mathrm{orig},n}$ and $P_{\mathrm{copy},n}$ be the projections associated to
the decomposition
\[
\ell^2(X_n)=\Hil_{\mathrm{orig},n}\oplus\Hil_{\mathrm{copy},n}.
\]
Define the split operator
\[
H_{\pm,n}^{\mathrm{split}}
:=
P_{\mathrm{orig},n}H_{\pm,n}P_{\mathrm{orig},n}
+
P_{\mathrm{copy},n}H_{\pm,n}P_{\mathrm{copy},n}.
\]

\begin{proposition}[Finite collar-locality]\label{prop:collar_locality_finite}
Assume that $H_{\pm,n}$ has propagation at most $R_{\mathrm{loc}}$ for every $n$, and set
\[
Y_n:=Y_n^{(R_{\mathrm{loc}})}=N_{R_{\mathrm{loc}}}(\Gamma_n).
\]
Then
\begin{equation}\label{eq:finite_collar}
\chi_{X_n\setminus Y_n}\,\bigl(H_{\pm,n}-H_{\pm,n}^{\mathrm{split}}\bigr)\,\chi_{X_n\setminus Y_n}=0.
\end{equation}
Moreover, for every $M>0$,
\begin{equation}\label{eq:finite_collar_jam}
\chi_{X_n\setminus Y_n}\,\Bigl(H_{\pm,M,n}-\bigl(H_{\pm,n}^{\mathrm{split}}+MP_{\mathrm{copy},n}\bigr)\Bigr)\,\chi_{X_n\setminus Y_n}=0.
\end{equation}
\end{proposition}
\proofappendixref{app:proof:prop:collar_locality_finite}




\begin{corollary}[Support near the coarse interface]\label{cor:collar_locality_coarse}
Assume \ref{ass:uniform_bulk}. Let
\[
H_\pm^{\mathrm{split},\infty}:=\bigoplus_{n=1}^\infty H_{\pm,n}^{\mathrm{split}},
\qquad
P_{\mathrm{copy}}^\infty:=\bigoplus_{n=1}^\infty P_{\mathrm{copy},n},
\qquad
J_{\pm,M}^\infty:=\bigoplus_{n=1}^\infty H_{\pm,M,n}.
\]
Then
\[
H_\pm^\infty-H_\pm^{\mathrm{split},\infty}\in C^*(Y\subset X),
\qquad
J_{\pm,M}^\infty-\bigl(H_\pm^{\mathrm{split},\infty}+MP_{\mathrm{copy}}^\infty\bigr)\in C^*(Y\subset X),
\]
where $Y=Y^{(R_{\mathrm{loc}})}$.
\end{corollary}
\proofappendixref{app:proof:cor:collar_locality_coarse}




For later use define
\[
P_{\mathrm{orig}}^\infty:=\bigoplus_{n=1}^\infty P_{\mathrm{orig},n},
\qquad
H_{\pm,\mathrm{copy}}^\infty:=\bigoplus_{n=1}^\infty P_{\mathrm{copy},n}H_{\pm,n}P_{\mathrm{copy},n},
\]
and for each $n$,
\[
H_{\pm,n}^{\mathrm{copy}}
:=
P_{\mathrm{copy},n}H_{\pm,n}P_{\mathrm{copy},n}\big|_{\Hil_{\mathrm{copy},n}}.
\]
For each fixed sign $\pm$, define the uniform negative copy-side threshold
\[
B_{\mathrm{copy},-}
:=
\sup_n \max\bigl\{0,\,-\lambda_{\min}(H_{\pm,n}^{\mathrm{copy}})\bigr\}.
\]
By Assumption~\ref{ass:uniform_bulk}(b), this quantity is finite and satisfies
\[
B_{\mathrm{copy},-}\le B_{\mathrm{bulk}}.
\]
For each $M>0$, define the auxiliary negative-side comparison operator
\[
\widetilde T^-_{\pm,M}:=
P_{\mathrm{orig}}^\infty+H_{\pm,\mathrm{copy}}^\infty+MP_{\mathrm{copy}}^\infty.
\]

\subsection{The jammed side is topologically trivial}

\begin{proposition}[Uniform positivity of the auxiliary negative-side model]\label{prop:trivial_jam}
Assume \ref{ass:uniform_bulk}. If
\[
M>B_{\mathrm{copy},-},
\]
then the auxiliary comparison operator
\[
\widetilde T^-_{\pm,M}
=
P_{\mathrm{orig}}^\infty+H_{\pm,\mathrm{copy}}^\infty+MP_{\mathrm{copy}}^\infty
\]
is bounded, self-adjoint, local, and strictly positive as an operator on $\ell^2(X)$. Moreover,
\[
\spec_{\mathrm{op}}(\widetilde T^-_{\pm,M})\subset \{1\}\cup [M-B_{\mathrm{copy},-},\infty),
\]
so
\[
\dist\bigl(0,\spec_{\mathrm{op}}(\widetilde T^-_{\pm,M})\bigr)\ge
\min\{1,M-B_{\mathrm{copy},-}\}.
\]
In particular,
its negative operator spectral projection vanishes:
\[
\chi_{(-\infty,0)}\!\bigl(\widetilde T^-_{\pm,M}\bigr)=0.
\]
The stronger condition $M>B_{\mathrm{bulk}}$ is sufficient but not necessary.
\end{proposition}
\proofappendixref{app:proof:prop:trivial_jam}




\begin{remark}\label{rem:why_aux_negative}
The bare operator $H_{\pm,\mathrm{copy}}^\infty+MP_{\mathrm{copy}}^\infty$ is positive on the copy
sector but has a zero block on the original-side sector when viewed as an operator on all of
$\ell^2(X)$. It is therefore \emph{not} gapped at $0$ on the full Hilbert space. The added term
$P_{\mathrm{orig}}^\infty$ is inserted solely to remove that spurious zero block while preserving the
negative-side agreement away from the interface. Any fixed positive scalar multiple
$\mu P_{\mathrm{orig}}^\infty$ with $\mu>0$ could be used in place of
$P_{\mathrm{orig}}^\infty$.
\end{remark}

\subsection{Quotient decomposition away from the interface}

For each $n$, define
\[
X_n^+:=X_{n,\mathrm{orig}}\cup Y_n,
\qquad
X_n^-:=X_{n,\mathrm{copy}}\cup Y_n,
\qquad
Y_n:=Y_n^{(R_{\mathrm{loc}})}.
\]
Set
\[
X^+:=\bigsqcup_{n=1}^\infty X_n^+,
\qquad
X^-:=\bigsqcup_{n=1}^\infty X_n^-.
\]
Then $X=X^+\cup X^-$ and $X^+\cap X^-=Y$.

\begin{proposition}[Coarse excision of the doubled decomposition]\label{prop:coarse_excision}
Under Assumption~\ref{ass:uniform_bulk}, the decomposition $X=X^+\cup X^-$ is coarsely
excisive: for every $r>0$ there exists $s(r)>0$ such that
\[
N_r(X^+)\cap N_r(X^-)\subset N_{s(r)}(Y).
\]
\end{proposition}
\proofappendixref{app:proof:prop:coarse_excision}




Let
\[
\mathcal A:=C^*(X),
\qquad
\mathcal I:=C^*(Y\subset X),
\qquad
\pi:\mathcal A\to \mathcal A/\mathcal I
\]
be the quotient map. Also set
\[
Z^+:=X^+\setminus Y,
\qquad
Z^-:=X^-\setminus Y.
\]
Let
\[
Q_+:=\chi_{Z^+},
\qquad
Q_-:=\chi_{Z^-},
\qquad
Q_Y:=\chi_Y
\]
be the corresponding diagonal multiplication operators on $\ell^2(X)$.

\begin{remark}[Strict-side projections in $C^*(X)$]\label{rem:strict_side_multipliers}
Because $C^*(X)$ is unital in the scalar-coefficient setting, the diagonal $\{0,1\}$-valued
operators $Q_+$, $Q_-$, and $Q_Y$ all belong to $\mathcal A=C^*(X)$: each has propagation $0$
and is locally compact. They are pairwise orthogonal projections summing to $\Id$ and they
preserve $\mathcal I$ under left and right multiplication.
\end{remark}

\begin{proposition}[Quotient splitting by strict-side compressions]\label{prop:quotient_split}
Assume \ref{ass:uniform_bulk}. With the notation above:
\begin{enumerate}[label=(\roman*),leftmargin=2em]
\item for every algebraic Roe operator $T$ of propagation at most $R$,
\[
Q_YT,\;TQ_Y,\;Q_+TQ_-,\;Q_-TQ_+ \in \mathcal I;
\]
\item if
\[
e_+:=\pi(Q_+),
\qquad
e_-:=\pi(Q_-),
\]
then $e_+$ and $e_-$ are orthogonal central projections in $\mathcal A/\mathcal I$ with
\[
e_++e_-=\Id_{\mathcal A/\mathcal I};
\]
\item the formulas
\[
E_+(\pi(a)):=\pi(Q_+a),
\qquad
E_-(\pi(a)):=\pi(Q_-a)
\qquad (a\in\mathcal A)
\]
equivalently $E_\pm(b)=e_\pm b$, 
define well-defined $*$-homomorphisms
\[
E_\pm:\mathcal A/\mathcal I\longrightarrow \mathcal A/\mathcal I
\]
such that
\[
E_\pm^2=E_\pm,
\qquad
E_+E_-=E_-E_+=0,
\qquad
E_+ + E_- = \Id_{\mathcal A/\mathcal I},
\]
and
\[
E_\pm(b_1)b_2=b_1E_\pm(b_2)
\qquad\text{for all }b_1,b_2\in\mathcal A/\mathcal I.
\]
Consequently,
\[
\mathcal A/\mathcal I
=
E_+(\mathcal A/\mathcal I)\oplus E_-(\mathcal A/\mathcal I)
\]
as a direct sum of closed two-sided ideals.
\end{enumerate}
\end{proposition}
\proofappendixref{app:proof:prop:quotient_split}




\begin{proposition}[Quotient-side description]\label{prop:quotient_sidewise}
Assume \ref{ass:uniform_bulk}. Let $J,T^+,T^-$ be bounded self-adjoint local operators on
\(\ell^2(X)\). Since \(X\) is a discrete bounded-geometry scalar Roe space, these operators belong
to \(\mathcal A=C^*(X)\). Define
\[
P^\pm:=\chi_{(-\infty,0)}(T^\pm),
\]
and assume:
\begin{enumerate}[label=(\alph*),leftmargin=2em]
\item $0<\gamma_\pm:=\dist\bigl(0,\spec_{\mathrm{op}}(T^\pm)\bigr)$;
\item there exists $R_{\mathrm{ag}}>0$ such that
\[
\chi_{X^+\setminus N_{R_{\mathrm{ag}}}(Y)}(J-T^+)\chi_{X^+\setminus N_{R_{\mathrm{ag}}}(Y)}=0,
\]
\[
\chi_{X^-\setminus N_{R_{\mathrm{ag}}}(Y)}(J-T^-)\chi_{X^-\setminus N_{R_{\mathrm{ag}}}(Y)}=0.
\]
\end{enumerate}
Then:
\begin{enumerate}[label=(\roman*),leftmargin=2em]
\item
\[
\pi(J)=E_+(\pi(T^+))+E_-(\pi(T^-));
\]
\item for every $0<\gamma<\min\{\gamma_+,\gamma_-\}$, every
$B\ge \max\{\opnorm{T^+},\opnorm{T^-}\}$, and every $h\in C_c(\mathbb R,[0,1])$ satisfying
\[
h(t)=1\quad\text{for }t\in[-B,-\gamma],
\qquad
h(t)=0\quad\text{for }t\in[\gamma,B],
\]
one has
\[
h(\pi(J))
=
E_+(\pi(P^+))+E_-(\pi(P^-));
\]
\item the element
\[
p_J:=E_+(\pi(P^+))+E_-(\pi(P^-))
\]
is a projection in $\mathcal A/\mathcal I$, independent of the choice of $\gamma$ and $h$, and for every such $h$
\[
p_J=h(\pi(J))=\pi(h(J)).
\]
\end{enumerate}
\end{proposition}
\proofappendixref{app:proof:prop:quotient_sidewise}




\begin{remark}[Relation to the literal quotient spectral projection]\label{rem:literal_quotient_projection}
Since $\mathcal A/\mathcal I$ is unital, the projection $p_J$ from
Proposition~\ref{prop:quotient_sidewise} coincides with $\chi_{(-\infty,0)}(\pi(J))$ whenever
$0\notin\spec_{\mathcal A/\mathcal I}(\pi(J))$. The continuous-cutoff formulation is used because
it avoids the need to verify that last spectral condition. The compactly supported version above is
also the form compatible with nonunital or stabilized Roe-algebra presentations.
\end{remark}

\begin{definition}[Sidewise class-A interface]\label{def:sidewise_interface}
Let $J$, $T^+$, and $T^-$ be bounded self-adjoint local operators on $\ell^2(X)$. We say that
$(J;T^+,T^-)$ is a \emph{sidewise class-A interface along $Y$} if:
\begin{enumerate}[label=(\roman*),leftmargin=2em]
\item $0$ lies in a gap of the operator spectra $\spec_{\mathrm{op}}(T^\pm)$;
\item there exists $R>0$ such that
\[
\chi_{X^+\setminus N_R(Y)}(J-T^+)\chi_{X^+\setminus N_R(Y)}=0,
\qquad
\chi_{X^-\setminus N_R(Y)}(J-T^-)\chi_{X^-\setminus N_R(Y)}=0;
\]
\item the operator Fermi projections
\[
P^+:=\chi_{(-\infty,0)}(T^+),
\qquad
P^-:=\chi_{(-\infty,0)}(T^-)
\]
belong to $\mathcal A=C^*(X)$.
\end{enumerate}
\end{definition}

\subsection{Boundary class of a liftable quotient projection}

\begin{proposition}[Boundary class of a liftable quotient projection]\label{prop:boundary_projection}
Let $p\in M_n(\mathcal A/\mathcal I)$ be a projection that admits a self-adjoint lift
$\widetilde p\in M_n(\mathcal A)$. Then
\[
\exp(2\pi i\,\widetilde p)-\Id \in M_n(\mathcal I),
\]
so $\exp(2\pi i\,\widetilde p)\in \Id+M_n(\mathcal I)$ defines a class in $K_1(\mathcal I)$. This
class depends only on $p$, not on the chosen self-adjoint lift. We denote it by
\[
\partial_Y(p)\in K_1(\mathcal I).
\]
This fixes our sign convention for the connecting map; some references use the inverse class.

Moreover, if $a\in M_n(\mathcal A)$ is self-adjoint and $h\in C(\mathbb R,[0,1])$ satisfies
\[
h|_{\spec(\pi(a))}=\chi_{(-\infty,0)}|_{\spec(\pi(a))},
\]
then $\pi(h(a))$ is a projection in $M_n(\mathcal A/\mathcal I)$ and
\[
\partial_Y\!\bigl(\pi(h(a))\bigr)
=
\bigl[\exp(2\pi i\,h(a))\bigr]\in K_1(\mathcal I).
\]
\end{proposition}
\proofappendixref{app:proof:prop:boundary_projection}




\begin{proposition}[Homotopy invariance of positive-side boundary classes]\label{prop:positive_boundary_homotopy}
Assume \ref{ass:uniform_bulk}, so that the quotient-side map $E_+$ is defined. Let
$s\mapsto P_s$, $s\in[0,1]$, be a norm-continuous path of projections in $\mathcal A=C^*(X)$. Then
\[
\partial_Y\!\bigl(E_+(\pi(P_s))\bigr)\in K_1(\mathcal I)
\]
is independent of $s$.
\end{proposition}
\proofappendixref{app:proof:prop:positive_boundary_homotopy}




\begin{definition}[Interface class of a sidewise realization]\label{def:interface_class}
If $(J;T^+,T^-)$ is a sidewise class-A interface along $Y$, let
\[
P^+:=\chi_{(-\infty,0)}(T^+),
\qquad
P^-:=\chi_{(-\infty,0)}(T^-),
\]
and define the canonical quotient projection
\[
p_J:=E_+(\pi(P^+))+E_-(\pi(P^-))\in \mathcal A/\mathcal I.
\]
By Proposition~\ref{prop:quotient_sidewise}, if
\[
0<\gamma<\min\Bigl\{\dist\bigl(0,\spec_{\mathrm{op}}(T^+)\bigr),\dist\bigl(0,\spec_{\mathrm{op}}(T^-)\bigr)\Bigr\}
\]
and, for some $B\ge\max\{\opnorm{T^+},\opnorm{T^-}\}$, the function
$h\in C_c(\mathbb R,[0,1])$ satisfies
\[
h(t)=1\ \text{ for }t\in[-B,-\gamma],
\qquad
h(t)=0\ \text{ for }t\in[\gamma,B],
\]
then
\[
p_J=h(\pi(J))=\pi(h(J)).
\]
Its \emph{interface class} is
\[
\Indint(J;T^+,T^-):=\partial_Y(p_J)\in K_1(\mathcal I).
\]

Equivalently, for every such cutoff $h$,
\[
\Indint(J;T^+,T^-)=\bigl[\exp(2\pi i\,h(J))\bigr]\in K_1(\mathcal I),
\]
and this class is independent of the choice of $h$.
\end{definition}

Sidewise interface-supported changes, finite-component changes, norm-vanishing component defects, sidewise homotopies, and positive copy-confinement variants are collected in Appendix~\ref{app:roe_interface_variants}.

\begin{remark}
The interface class is built from the quotient
projection $p_J$ and the genuine extension
\[
0\to C^*(Y\subset X)\to C^*(X)\to C^*(X)/C^*(Y\subset X)\to 0.
\]
\end{remark}

\subsection{Interface class for the jammed family}

\begin{theorem}[Doubled-and-jammed interface class]\label{thm:bulkboundary_main}
Assume \ref{ass:uniform_bulk}. Let
\[
J_{\pm,M}^\infty=\bigoplus_{n=1}^\infty H_{\pm,M,n}
\]
be the direct-sum jammed family on the coarse disjoint union $X$, and let
\[
Y=Y^{(R_{\mathrm{loc}})}
\]
be the coarse interface from Corollary~\ref{cor:collar_locality_coarse}. If
\[
M>B_{\mathrm{copy},-},
\]
then with
\[
T^+:=H_\pm^\infty,
\qquad
T^-:=\widetilde T^-_{\pm,M},
\]
the triple $(J_{\pm,M}^\infty;T^+,T^-)$ is a sidewise class-A interface along $Y$ in the sense of
Definition~\ref{def:sidewise_interface}. Its canonical quotient projection is
\begin{equation}\label{eq:quotient_fermi_projection}
p_{J_{\pm,M}^\infty}=E_+\!\bigl(\pi(P_\pm)\bigr),
\end{equation}
where
\[
P_\pm:=\chi_{(-\infty,0)}(H_\pm^\infty)\in C^*(X)
\]
is the bulk operator Fermi projection from Proposition~\ref{prop:bulk_class}. Equivalently, if
\[
0<\gamma<\min\Bigl\{g_{\mathrm{bulk}},\,1,\,M-B_{\mathrm{copy},-}\Bigr\}
\]
and, for some $B\ge\max\{\opnorm{T^+},\opnorm{T^-}\}$, the function
$h\in C_c(\mathbb R,[0,1])$ satisfies
\[
h(t)=1\ \text{ for }t\in[-B,-\gamma],
\qquad
h(t)=0\ \text{ for }t\in[\gamma,B],
\]
then
\[
\pi(h(J_{\pm,M}^\infty))=E_+\!\bigl(\pi(P_\pm)\bigr).
\]
Consequently,
\begin{equation}\label{eq:main_interface_class}
\Indint(J_{\pm,M}^\infty;T^+,T^-)
=
\partial_Y\!\bigl(E_+(\pi(P_\pm))\bigr)
\in K_1\!\bigl(C^*(Y\subset X)\bigr).
\end{equation}
In particular, the jammed side contributes no negative-energy quotient projection.
\end{theorem}
\proofappendixref{app:proof:thm:bulkboundary_main}


This theorem uses the closed-bulk gap in Assumption~\ref{ass:uniform_bulk}(c), together with the
uniform locality and boundedness hypotheses stated there. In the headline application,
these hypotheses are supplied by Theorem~\ref{thm:hp_closed_input}. The finite-dimensional
decoupling estimates of Section~\ref{M-sec:decoupling} do not use a refinement-limit bulk gap; they
require only their stated finite-matrix separation hypotheses.



Large-jam-scale independence is recorded with the other Roe variants in Appendix~\ref{app:roe_interface_variants}.

\subsection{Identification of the bulk class}

The Hamiltonians $H_{\pm,n}=H_\pm(D(\K_n))$ are exactly the Higson--Prodan universal Chern
Hamiltonians on the closed oriented surfaces $D(\K_n)$. For closed doubles covered by the
Higson--Prodan theorem, the bulk Fermi projection
$P_\pm=\chi_{(-\infty,0)}(H_\pm^\infty)$ is the corresponding Higson--Prodan Fermi projection.

In this subsection \(\Ch\) denotes the numerical bulk Chern pairing supplied by the
Higson--Prodan theorem for the closed refinement family. We do not construct a new Chern cocycle on
an arbitrary Roe algebra here. Theorem~\ref{M-thm:pathb_roe_winding} uses the Higson--Prodan bulk
pairing associated with the closed doubled family and identifies the original/copy coarse interface
pairing internally.

\begin{theorem}[Bulk Chern number of the doubled family, imported from Higson--Prodan]\label{thm:bulk_chern}
Assume \ref{ass:uniform_bulk}. Assume moreover that the refinement family of closed doubles
$(D(\K_n))_n$ is an admissible Higson--Prodan refinement family satisfying the hypotheses of the
Higson--Prodan bulk theorem, and that the Chern pairing $\Ch$ used here is the corresponding
Higson--Prodan bulk Chern pairing. Then the $K_0$-class
$[P_\pm]\in K_0(C^*(X))$ is the Higson--Prodan bulk class for $(D(\K_n))_n$.
Consequently, by the Higson--Prodan theorem
\cite{HigsonProdanUniversalPRL,HigsonProdanUniversalArxiv},
\[
\Ch([P_\pm])=\pm1.
\]
\end{theorem}
\proofappendixref{app:proof:thm:bulk_chern}




The next definition records the Ewert--Meyer convention used to phrase the restatement of
Theorem~\ref{M-thm:pathb_roe_winding} in bulk--interface language. The intermediate Chern-\(\pm1\)
and HP corollaries are collected in Appendix~\ref{app:roe_interface_variants}.

\begin{definition}[HP--Ewert--Meyer boundary convention]\label{def:hp_em_admissible}\label{ass:coarse_bec}
A boundary refinement family \((\K_n)_n\) has the \emph{HP--Ewert--Meyer boundary convention} when
\(D(\K_n)\) satisfies the Higson--Prodan closed-bulk input in
Assumption~\ref{ass:hp_admissibility}, and the associated original/copy coarse boundary data
\[
(X;X^+,X^-;Y)
\]
fit the Roe-algebraic bulk--interface framework of Ewert--Meyer \cite{EwertMeyer2019}.
Concretely, the required data are:
\begin{enumerate}[label=(\alph*),leftmargin=2em]
\item \(X=X^+\cup X^-\) is a coarse partition with interface \(Y\); in particular, for every
  \(r>0\) there is \(s(r)>0\) such that
  \[
  N_r(X^+)\cap N_r(X^-)\subset N_{s(r)}(Y).
  \]
\item Away from \(Y\), the positive side agrees with the Higson--Prodan bulk model and the
  negative side agrees with a trivial gapped comparison model.
\item The connecting map
  \[
  \partial_Y:K_0(C^*(X)/C^*(Y\subset X))\longrightarrow K_1(C^*(Y\subset X))
  \]
  is paired with the Ewert--Meyer interface cocycle using the sign convention of
  Theorem~\ref{M-thm:pathb_roe_winding}; a quotient class with Chern number \(c_+\) on the positive
  side and \(c_-\) on the negative side has winding number \(c_+-c_-\).
\end{enumerate}
\end{definition}

Convention checks and tail variants are collected in Appendix~\ref{app:roe_interface_variants}.

The Ewert--Meyer input used here is the coarse Mayer--Vietoris bulk--interface boundary map and
its identification with the bulk--edge map
\cite[Propositions~3.2.1--3.2.2 and Corollary~3.2.3]{EwertMeyer2019}. Kubota's controlled
framework and the Ludewig--Thiang edge-following theorem give closely related formulations
\cite{KubotaControlled,LudewigThiangEdge}; the restatement below uses the Ewert--Meyer convention
fixed in Definition~\ref{def:hp_em_admissible}.

\begin{center}
\fbox{\begin{minipage}{0.92\linewidth}
\textbf{Status of the interface \(K_1\)-class identity and integer winding.}
For the doubled-jammed original/copy partition, Section~\ref{M-sec:pathb_internalization}
proves the \(K_1\)-class identity internally. The restatement below records the same
\(K_1\)-class identity in Ewert--Meyer language, with Higson--Prodan supplying the
closed-bulk gap and Chern value. The integer winding evaluation is internal in the
straight periodic Toeplitz case (\Cref{M-cor:toeplitz_winding}) and is the conditional
corollary \Cref{M-cor:conditional_winding} for general partitions.
\end{minipage}}
\end{center}

\begin{theorem}[Ewert--Meyer restatement of the Roe interface theorem]\label{thm:full_hp_bec}
Let \((\K_n)_n\) be a uniformly regular boundary refinement family whose closed doubles satisfy the
Higson--Prodan universal-model theorem, and use the original/copy coarse partition. Then, for every
\(M>B_{\mathrm{copy},-}\),
\[
\Indint(J_{\pm,M}^\infty;T^+,T^-)
=
\partial_Y\!\bigl(E_+(\pi(P_\pm))\bigr)
\in K_1(C^*(Y\subset X)).
\]
The projection \(P_\pm\) is the Higson--Prodan closed-bulk Fermi projection and
\[
\Ch([P_\pm])=\pm1.
\]
In the straight periodic Toeplitz case (\Cref{M-cor:toeplitz_winding}), the integer interface
winding pairing is internal and evaluates to
\[
\operatorname{Wind}\!\left(
\Indint(J_{\pm,M}^\infty;T^+,T^-)
\right)
=\pm1.
\]
For general partitions, the integer interface winding pairing is recorded as a conditional
corollary (\Cref{M-cor:conditional_winding}) requiring an integer-valued Chern-jump
pairing on \(K_1(\mathcal I)\), whose construction outside the periodic class is not carried
out in this paper. The \(K_1\)-class identity itself is unconditional modulo HP.  The same
class is independent of the jam scale \(M\) for all \(M>B_{\mathrm{copy},-}\).
\end{theorem}
\proofappendixref{app:proof:thm:full_hp_bec}




Relabeling, orientation reversal, and robustness of the compatible winding are stated in Appendix~\ref{app:roe_interface_variants}.

\begin{remark}[From bulk Chern number to numerical interface index]\label{rem:chern_to_interface}
Theorem~\ref{M-thm:pathb_roe_winding} is the headline winding result: for the doubled-jammed
original/copy partition it supplies the Ewert--Meyer identification internally and uses the
Higson--Prodan closed-bulk theorem for the gap and Chern value. Theorem~\ref{thm:full_hp_bec}
records the same conclusion in Ewert--Meyer language. This identification is not a boundary-gap
statement and does not assert that the negative quotient of the closed doubled bulk projection
vanishes. In the corresponding six-term exact
sequence
\[
K_0(\mathcal I)\to K_0(\mathcal A)\to K_0(\mathcal A/\mathcal I)
\xrightarrow{\;\partial\;}
K_1(\mathcal I)\to K_1(\mathcal A)\to K_1(\mathcal A/\mathcal I)
\]
the connecting map is evaluated by the interface winding pairing. The positive side has Chern
number \(\pm1\), and the jammed negative side is trivial, so the Chern jump is \(\pm1\).
\end{remark}

\begin{remark}[What the global theorem does and does not say]\label{rem:componentwise}
Theorem~\ref{thm:bulkboundary_main} is a \emph{global} statement about the interface class
supported near the entire glued boundary. It does not by itself imply that each boundary
component carries its own local index or local boundary map. A componentwise statement requires
an additional localized pairing. In the continuum Hall setting, Drouot and Zhu prove a
curved-interface formula in which the local edge conductance is an intersection number times the
bulk Hall conductance \cite{DrouotZhuBEC,DrouotZhuEdgeSpectrum}. We view that as the correct
model for any future local theorem in the present lattice setting.
\end{remark}

\FloatBarrier



\section{Detailed appendix proofs}\label{app:proofs}


\subsection{Sign ledger}\label{app:sign_ledger}
The table records the orientation conventions used in the periodic triangular crossing and in the
Toeplitz edge map. It is a ledger for the proof of Theorem~\ref{M-thm:periodic_triangular_chern}, not
a separate numerical input.

\paragraph{Explicit sign chain.}
The sign convention used in the headline corollary \(\operatorname{wind}=\mp\det[t\ n]\) for the
periodic triangular case is fully determined by composing the following four signed steps; a
referee can verify each step by referring to the row of the table below and the cited
proposition.
(i) \emph{Bloch torus orientation.} The Bloch torus carries the orientation
\(dk_1\wedge dk_2\); the Chern integer
\(\Ch(P)=\frac{1}{2\pi i}\int_{\mathbb T^2}\tr(P\,dP\wedge dP)\) uses this orientation, giving
\(\Ch(E_-(H_+^\triangle))=-1\) and \(\Ch(E_-(H_-^\triangle))=+1\) by
Theorem~\ref{M-thm:periodic_triangular_chern}.
(ii) \emph{Boundary map sign.} The Higson--Roe \(K\)-theory connecting map \(\partial_Y\) is
fixed by Convention~\ref{M-conv:pathb_sign} of the main paper: for any self-adjoint lift
\(\tilde P\) of a quotient projection \(\pi(P)\),
\(\partial[\pi(P)]=[\exp(2\pi i\,\tilde P)]\), with no sign flip.
(iii) \emph{Toeplitz extension sign.} The half-space Toeplitz extension in the positive-normal
direction \(n\) is fixed by Theorem~\ref{M-thm:periodic_toeplitz_pairing}: the winding-number
pairing satisfies \(\operatorname{wind}(\partial_{\mathrm T,t,n}([P]))=\varepsilon(t,n)\,\Ch(P)\)
where \(\varepsilon(t,n)=\det[t\ n]\in\{\pm 1\}\) records the change-of-basis Jacobian from the
canonical \((k_1,k_2)\) frame to the tangential/positive-normal frame \((\kappa,\eta)\).
(iv) \emph{Composite winding sign.} Composing (ii) and (iii) gives
\(\operatorname{wind}=\varepsilon(t,n)\,\Ch(E_-(H_\pm^\triangle))=\mp\,\varepsilon(t,n)
=\mp\det[t\ n]\), in agreement with both
\Cref{M-cor:toeplitz_winding} and \Cref{M-cor:periodic_straight_edge_winding}.

If any one of the four sign choices above were flipped, the displayed winding would be the
opposite integer; for instance, replacing the Bloch orientation by \(dk_2\wedge dk_1\) would
swap the Chern signs of step (i) (and hence flip the right-hand side in step (iv)), while
replacing \(n\) by \(-n\) in step (iii) would flip \(\varepsilon(t,n)\). Because all four
choices are fixed once for all in Convention~\ref{M-conv:pathb_sign} and in the table below,
the cited final integer is uniquely determined.
\begin{center}
\small
\begin{tabular}{>{\raggedright\arraybackslash}p{0.25\linewidth}
                >{\raggedright\arraybackslash}p{0.30\linewidth}
                >{\raggedright\arraybackslash}p{0.32\linewidth}}
\toprule
Object & Orientation convention & Sign consequence \\
\midrule
Bloch torus & \(dk_1\wedge dk_2\) & FHS plaquettes use the same ordered coordinate square. \\
Crossing space & \(dq_1\wedge dq_2\wedge d\mu\) & The local degree is the sign of \(\det Du\). \\
Self-dual basis & \(e_1,e_2,e_3,e_4\) oriented by \(e_1\wedge e_2\wedge e_3\wedge e_4\) & The self-dual compression has degree \(-1\). \\
Mass reversal & \(k\mapsto -k\) on \(\mathbb T^2\) & The base map preserves orientation; complex conjugation reverses \(c_1\). \\
Toeplitz edge & Tangential/positive-normal basis \((t,n)\) & Winding is \(\det[t\ n]\Ch(P)\). \\
\bottomrule
\end{tabular}
\end{center}
\subsection{Routine perturbation proofs}\label{app:routine_perturbation_proofs}
The two auxiliary finite-dimensional perturbation statements below are retained for reference and
kept out of the body. Their proofs are included with the other appendix proof blocks.
%====================================================================
\subsection{Two-band hexagonal and square benchmark symbols}\label{app:benchmark_hex_square}\label{app:benchmark_hexagonal}\label{app:benchmark_square}
%====================================================================

\begin{remark}[Status of the two-band benchmark symbols]\label{rem:two_band_benchmarks}
The two-band hexagonal and square symbols collected in this subsection are auxiliary benchmark
classes; they are not derivations of the full Higson--Prodan universal symbol for arbitrary
hexagonal or square triangulations. They feed only the cross-check
Remark~\ref{M-lem:pathb_hex_square}, whose role is a sanity check against the periodic triangular
calculation in Section~\ref{M-sec:periodicbulk}. The reader interested only in the headline
Theorem~\ref{M-thm:pathb_roe_winding} may skip this subsection on first reading. The underlying
proofs are collected in
\Cref{app:proof:lem:hex_chern_positive,app:proof:lem:hex_chern_negative,app:proof:lem:hex_toeplitz_winding,app:proof:thm:hex_bulk_edge,app:proof:lem:square_chern_positive,app:proof:lem:square_toeplitz_winding,app:proof:thm:square_bulk_edge}.
\end{remark}

\paragraph{Hexagonal symbol.}
With \(k=(k_1,k_2)\in\mathbb T^2\), \(a_1(k)=\sin k_1\), \(a_2(k)=\sin k_2\),
\(a_3^{\hexagon}(k)=\cos k_1+\cos k_2-1\):

\begin{definition}[Hexagonal two-band Bloch representative]\label{def:hex_bloch_symbol}
For \(m\in\mathbb R\),
\(H_{\hexagon}(k;m)=
\begin{pmatrix} m\,a_3^{\hexagon}(k) & a_1(k)-i a_2(k)\\ a_1(k)+i a_2(k) & -m\,a_3^{\hexagon}(k)\end{pmatrix}\).
\end{definition}

\begin{remark}[Massless hexagonal closure]\label{lem:hex_massless_closure}
\(H_{\hexagon}(0,0;0)\) is gapless at zero (direct substitution; see
\Cref{app:proof:lem:hex_massless_closure} for the one-line verification).
\end{remark}

\begin{lemma}[Hexagonal determinant and bulk gap]\label{lem:hex_bulk_gap}
For \(m\ne0\) and \(k\in\mathbb T^2\),
\(|\det H_{\hexagon}(k;m)|\ge c_{\hexagon}(m):=\min\{1/4,(\sqrt{3}-1)^2m^2\}\), hence
\(\dist(0,\spec H_{\hexagon}(k;m))\ge \gamma_{\hexagon}(m):=\sqrt{c_{\hexagon}(m)}\).
\end{lemma}
\proofappendixref{app:proof:lem:hex_bulk_gap}

\begin{lemma}[Hexagonal Chern integer for \(m>0\)]\label{lem:hex_chern_positive}
For \(m>0\), with orientation \(dk_1\wedge dk_2\),
\(\Ch(E_-(H_{\hexagon}(\cdot;m)))=-1\).
\end{lemma}
\proofappendixref{app:proof:lem:hex_chern_positive}

\begin{lemma}[Hexagonal mass reversal]\label{lem:hex_chern_negative}
For \(m<0\), \(\Ch(E_-(H_{\hexagon}(\cdot;m)))=+1\).
\end{lemma}
\proofappendixref{app:proof:lem:hex_chern_negative}

\begin{lemma}[Hexagonal Toeplitz winding]\label{lem:hex_toeplitz_winding}
Let \(P_{\hexagon,m}=\chi_{(-\infty,0)}(H_{\hexagon}(\cdot;m))\); for a primitive
tangential/positive-normal lattice basis \((t,n)\),
\(\operatorname{wind}(\partial_{\mathrm T,t,n}([P_{\hexagon,m}]))=\det[t\ n]\,\Ch(P_{\hexagon,m})\),
which gives \(-1\) for \(m>0\) and \(+1\) for \(m<0\).
\end{lemma}
\proofappendixref{app:proof:lem:hex_toeplitz_winding}

\begin{theorem}[Hexagonal two-band bulk-edge correspondence]\label{thm:hex_bulk_edge}
For every \(m\ne0\), \(H_{\hexagon}(k;m)\) is uniformly gapped at zero with gap at least
\(\gamma_{\hexagon}(m)\), \(\Ch(P_{\hexagon,m})=-\sgn(m)\), and for every primitive straight
periodic edge with basis \((t,n)\),
\(\operatorname{wind}(\partial_{\mathrm T,t,n}([P_{\hexagon,m}]))=-\det[t\ n]\,\sgn(m)\).
\end{theorem}
\proofappendixref{app:proof:thm:hex_bulk_edge}

\paragraph{Square symbol.}
With \(b_1(k)=\sin k_1\), \(b_2(k)=\sin k_2\), \(b_3^{\square}(k)=1+\cos k_1+\cos k_2\):

\begin{definition}[Square two-band Bloch representative]\label{def:square_bloch_symbol}
For \(m\in\mathbb R\),
\(H_{\square}(k;m)=
\begin{pmatrix} m\,b_3^{\square}(k) & b_1(k)-i b_2(k)\\ b_1(k)+i b_2(k) & -m\,b_3^{\square}(k)\end{pmatrix}\).
\end{definition}

\begin{remark}[Massless square closure]\label{lem:square_massless_closure}
\(H_{\square}(0,0;0)\) is gapless at zero (direct substitution; see
\Cref{app:proof:lem:square_massless_closure} for the verification).
\end{remark}

\begin{lemma}[Square determinant and bulk gap]\label{lem:square_bulk_gap}
For \(m\ne0\) and \(k\in\mathbb T^2\),
\(|\det H_{\square}(k;m)|\ge c_{\square}(m):=\min\{1/4,(\sqrt{3}-1)^2m^2\}\), hence
\(\dist(0,\spec H_{\square}(k;m))\ge\gamma_{\square}(m):=\sqrt{c_{\square}(m)}\).
\end{lemma}
\proofappendixref{app:proof:lem:square_bulk_gap}

\begin{lemma}[Square Chern integer]\label{lem:square_chern_positive}
For \(m>0\), with orientation \(dk_1\wedge dk_2\),
\(\Ch(E_-(H_{\square}(\cdot;m)))=+1\).
\end{lemma}
\proofappendixref{app:proof:lem:square_chern_positive}

\begin{lemma}[Square Toeplitz winding]\label{lem:square_toeplitz_winding}
Let \(P_{\square,m}=\chi_{(-\infty,0)}(H_{\square}(\cdot;m))\); for a primitive
tangential/positive-normal lattice basis \((t,n)\),
\(\operatorname{wind}(\partial_{\mathrm T,t,n}([P_{\square,m}]))=\det[t\ n]\,\Ch(P_{\square,m})\),
which gives \(+1\) for \(m>0\) and \(-1\) for \(m<0\).
\end{lemma}
\proofappendixref{app:proof:lem:square_toeplitz_winding}

\begin{theorem}[Square two-band bulk-edge correspondence]\label{thm:square_bulk_edge}
For every \(m\ne0\), \(H_{\square}(k;m)\) is uniformly gapped at zero with gap at least
\(\gamma_{\square}(m)\), \(\Ch(P_{\square,m})=\sgn(m)\), and for every primitive straight periodic
edge with basis \((t,n)\),
\(\operatorname{wind}(\partial_{\mathrm T,t,n}([P_{\square,m}]))=\det[t\ n]\,\sgn(m)\).
\end{theorem}
\proofappendixref{app:proof:thm:square_bulk_edge}




\subsection{Section 5 variants}\label{app:section5_variants}
The periodic triangular section in the body keeps the closed-bulk calculation, the Chern-sign
calculation, and the primitive straight Toeplitz edge map. The following subsections collect the
variant statements that use the same inputs with extra quotient, tail, finite-defect, finite-corner,
or annular hypotheses.


\subsection{Robustness variants: finite defects}\label{app:robustness_finite_defects}
\begin{corollary}[Robust straight-periodic Toeplitz winding]\label{cor:straight_toeplitz_robustness}\label{cor:finite_defect_tail_toeplitz}
In the setting of Corollary~\ref{M-cor:periodic_straight_edge_winding}, the Toeplitz winding values
are unchanged by the following operations, provided the stated sidewise comparison gaps remain
open:
\begin{enumerate}[label=(\alph*),leftmargin=2em]
\item replacing the closed periodic triangular bulk by \(H_{\pm,m}^\triangle\) with fixed \(m>0\);
\item adding a self-adjoint finite-range periodic bulk perturbation \(V\) to \(H_{\pm,m}^\triangle\)
with \(\sup_k\opnorm{V(k)}<g_{\pm,m}\), and in particular adding \(V\) to \(H_\pm^\triangle\) with
\(\sup_k\opnorm{V(k)}<2^{-13}\);
\item adding any bounded finite-propagation self-adjoint perturbation supported in a fixed
straight-edge collar, including periodic edge or collar terms that are not norm-small;
\item replacing scalar jamming by a bounded componentwise scalar jam above the positivity
threshold, or by a sufficiently positive bounded finite-propagation copy-side confining potential;
\item changing finitely many initial components, adding norm-vanishing component defects, or passing
  through a norm-continuous sidewise homotopy with uniformly open comparison gaps.
\item replacing the straight interface subset \(Y\) by \(Y'\) at finite Hausdorff distance from \(Y\),
  with the same sidewise realization outside a fixed neighborhood of \(Y\cup Y'\);
\item adding a bounded self-adjoint finite-propagation perturbation to the jammed operator, or to the
  sidewise data with the comparison gaps still open, supported in a fixed edge collar and in a
  finite tangential region;
\item after deleting finitely many coarse-disjoint-union components, making arbitrary bounded
  finite-propagation changes on the deleted components when the tail has a split, finite-corner, or
  \(K_1\)-trivial-corner Toeplitz description.
\end{enumerate}
The same conclusion holds componentwise for every split finite family covered by
Theorem~\ref{thm:split_component_toeplitz} and for every finite-corner split family covered by
Theorem~\ref{thm:finite_corner_toeplitz}. The finite-defect and tail-split cases have the same
winding values as the associated straight periodic tail.
\end{corollary}
\proofappendixref{app:proof:cor:straight_toeplitz_robustness}





\begin{theorem}[Bounded-distance rough and eventually periodic collars]\label{thm:rough_eventual_toeplitz_equivalence}
Let \(Y_0\) be a primitive straight periodic interface for a finite-range periodic bulk projection
\(P\), and let \(Y_1\subset X\) be another interface subset at finite Hausdorff distance from
\(Y_0\). Then
\[
C^*(Y_0\subset X)=C^*(Y_1\subset X)=:\mathcal I.
\]
Suppose \((J_i;T_i^+,T_i^-)\), \(i=0,1\), are sidewise class-A interfaces along \(Y_i\), and that
\[
J_1-J_0,\qquad T_1^+-T_0^+,\qquad T_1^--T_0^-\in\mathcal I
\]
after viewing both realizations on the same Roe space. Then the quotient Fermi data agree:
\[
\pi(\chi_{(-\infty,0)}(T_1^\pm))
=
\pi(\chi_{(-\infty,0)}(T_0^\pm)),
\qquad
p_{J_1}=p_{J_0}.
\]
Consequently
\[
\Indint(J_1;T_1^+,T_1^-)=\Indint(J_0;T_0^+,T_0^-)\in K_1(\mathcal I),
\]
and the rough-edge winding is the straight Toeplitz winding
\[
\det[t\ n]\,\Ch(P)
\]
for the primitive tangential/positive-normal basis \((t,n)\) of the straight tail.

If the difference between the rough collar and the straight periodic collar is supported in a fixed
normal collar and in a finite tangential interval, then the displayed differences are finite rank.
Thus eventually periodic straight or bounded-distance rough collars with a finite transition region
have the same quotient class and the same winding as their periodic tail, provided the sidewise
comparison gaps remain open.
\end{theorem}
\proofappendixref{app:proof:thm:rough_eventual_toeplitz_equivalence}




The fixed-collar robustness checks in
\Cref{fig:disorder,fig:copy_confinement,tab:copy_confinement} are finite-matrix checks of the
same ideal-stability mechanism; the winding equality here is the quotient-ideal argument, not a
separate finite-volume index.





\subsection{Robustness variants: finite corners}\label{app:robustness_finite_corners}
\begin{theorem}[Finite-corner split Toeplitz additivity]\label{thm:finite_corner_toeplitz}
Let \(H\) and \(P\) be as in Theorem~\ref{M-thm:periodic_toeplitz_pairing}. Suppose an interface
ideal \(\mathcal I_{\mathrm{poly}}\) contains a finite-dimensional corner ideal
\(\mathcal F\) and fits into an exact sequence
\[
0\longrightarrow \mathcal F\longrightarrow \mathcal I_{\mathrm{poly}}
\xrightarrow{\;q\;}
\bigoplus_{j=1}^r \mathcal I_j
\longrightarrow0,
\]
where each \(\mathcal I_j\) is identified with the straight periodic Toeplitz edge ideal
\(C(S^1,\mathcal K)\) for a primitive tangential/normal lattice basis \((t_j,n_j)\). Assume that,
after quotienting by \(\mathcal F\), the corresponding quotient extension is the direct sum of
these \(r\) Toeplitz extensions. If
\[
\partial_{\mathrm{poly}}([P])\in K_1(\mathcal I_{\mathrm{poly}})
\]
is the boundary class for the polygonal or cornered interface extension, then
\[
q_*\partial_{\mathrm{poly}}([P])
=
\bigoplus_{j=1}^r\partial_{\mathrm T,t_j,n_j}([P]).
\]
Moreover \(q_*:K_1(\mathcal I_{\mathrm{poly}})\to
K_1(\bigoplus_j\mathcal I_j)\) is injective. Thus the finite corner ideal contributes no
\(K_1\)-coordinate, and the \(j\)-th component winding is
\[
\operatorname{wind}_j\bigl(\partial_{\mathrm{poly}}([P])\bigr)
=
\det[t_j\ n_j]\,\Ch(P).
\]
\end{theorem}
\proofappendixref{app:proof:thm:finite_corner_toeplitz}




\begin{corollary}[K1-trivial corner ideals]\label{cor:k1_trivial_corner_toeplitz}
In Theorem~\ref{thm:finite_corner_toeplitz}, the finite-dimensional hypothesis on
\(\mathcal F\) may be replaced by the weaker condition
\[
K_1(\mathcal F)=0.
\]
Under the same quotient-extension hypothesis, \(q_*\) is injective on \(K_1\), the finite-corner
class is determined by its straight-edge quotient coordinates, and
\[
\operatorname{wind}_j\bigl(\partial_{\mathrm{poly}}([P])\bigr)
=
\det[t_j\ n_j]\,\Ch(P).
\]
This applies, in particular, when \(\mathcal F\) is finite-dimensional, a \(c_0\)-direct sum of
finite-dimensional matrix algebras, or a bounded product of finite-dimensional matrix algebras.
\end{corollary}
\proofappendixref{app:proof:cor:k1_trivial_corner_toeplitz}




\begin{corollary}[Collar-separation criterion for finite-corner quotients]\label{cor:finite_corner_geometric_criterion}
Let \(Y=C\cup Y^{(1)}\cup\cdots\cup Y^{(r)}\) be a cornered interface collar in a bounded-geometry
Roe space \(X\). Suppose:
\begin{enumerate}[label=(\roman*),leftmargin=2em]
\item the corner-supported ideal \(\mathcal F_C\), obtained as the closure of operators supported
  in bounded neighborhoods of \(C\), has \(K_1(\mathcal F_C)=0\);
\item the edge pieces separate off the corner set: for every propagation bound \(R\) there is
  \(S(R)\) such that the sets
  \(N_R(Y^{(i)})\setminus N_{S(R)}(C)\) and
  \(N_R(Y^{(j)})\setminus N_{S(R)}(C)\) are disjoint whenever \(i\ne j\);
\item outside a bounded neighborhood of \(C\), each \(Y^{(j)}\) is identified with a primitive
  straight periodic collar, and restriction to that collar identifies its edge quotient ideal with
  the Toeplitz ideal \(\mathcal I_j\cong C(S^1,\mathcal K)\).
\end{enumerate}
Then the interface ideal \(\mathcal I_Y=C^*(Y\subset X)\) fits into an exact sequence
\[
0\longrightarrow\mathcal F_C\longrightarrow\mathcal I_Y
\xrightarrow{\;q\;}
\bigoplus_{j=1}^r\mathcal I_j\longrightarrow0.
\]
If the quotient extension is the corresponding direct sum of primitive straight Toeplitz
extensions, the component winding formula of Corollary~\ref{cor:k1_trivial_corner_toeplitz}
applies.
\end{corollary}
\proofappendixref{app:proof:cor:finite_corner_geometric_criterion}




\begin{remark}
The split hypotheses in Theorems~\ref{thm:split_component_toeplitz}
and~\ref{thm:finite_corner_toeplitz} are genuine extra localization inputs. They hold for finite
direct sums of independent straight periodic half-space collars and for cornered models where
removing a \(K_1\)-trivial corner ideal leaves such a direct sum, for instance under the
collar-separation criterion in Corollary~\ref{cor:finite_corner_geometric_criterion}. They are not
asserted for arbitrary curved or nonperiodic boundary components.
\end{remark}



\subsection{Robustness variants: tail-split and annular}\label{app:robustness_tail_annular}
\begin{corollary}[Finite covers and finite-index cell enlargements]\label{cor:periodic_covers_sublattices}
Let \(A\in M_2(\mathbb Z)\) have \(\det A\ne0\), and define the finite-cover pullback
\[
H_{\pm,A}^\triangle(\kappa):=H_\pm^\triangle(A\kappa),
\qquad \kappa\in\mathbb T^2 .
\]
Then \(H_{\pm,A}^\triangle\) is gapped at zero and
\[
\Ch\bigl(E_-(H_{+,A}^\triangle)\bigr)=-\det A,
\qquad
\Ch\bigl(E_-(H_{-,A}^\triangle)\bigr)=\det A .
\]
The same conclusion holds for the positive-mass symbols \(H_{\pm,m}^\triangle(A\kappa)\), \(m>0\).
Small self-adjoint finite-range perturbations with norm below the corresponding pulled-back gap
preserve these Chern numbers.

If the same periodic triangular operator is only rewritten with respect to a finite-index
sublattice of periods, equivalently with a larger real-space fundamental cell and the folded Bloch
Hamiltonian, the zero gap and the Chern signs \(-1\) for \(H_+^\triangle\) and \(+1\) for
\(H_-^\triangle\) are unchanged.
\end{corollary}
\proofappendixref{app:proof:cor:periodic_covers_sublattices}



\begin{corollary}[Finite periodic triangular tori]\label{cor:finite_periodic_triangular_tori}
Every finite \(N_1\times N_2\) periodic quotient obtained from the symbol above has a zero gap
bounded below by \(2^{-13}\), uniformly in \(N_1,N_2\). For quotients large enough to be embedded
as simplicial tori, this is the corresponding finite periodic triangular torus. Its occupied
finite-volume Bloch states are samples of the same rank-three Bloch bundle.
\end{corollary}
\proofappendixref{app:proof:cor:finite_periodic_triangular_tori}


The next interface statement uses the internal periodic Bloch gap and Chern sign from
Theorems~\ref{M-thm:periodic_triangular_bulk} and~\ref{M-thm:periodic_triangular_chern}. Straight
periodic winding values are evaluated by the Toeplitz extension inside the paper; nonstraight local
pairings are not asserted in this periodic block.



\begin{corollary}[Periodic annular closed-bulk input and interface class]\label{cor:self_contained_annular}
Let \((\K_N)_N\) be a sequence of triangular-lattice cylinders whose doubles are the periodic
triangular tori of Corollary~\ref{cor:finite_periodic_triangular_tori}. Then the closed-bulk gap
and Chern value needed for the doubled-and-jammed boundary construction are supplied by
Theorems~\ref{M-thm:periodic_triangular_bulk}
and~\ref{M-thm:periodic_triangular_chern}, not by the Higson--Prodan closed-bulk theorem. Hence, for
every sufficiently large jam scale \(M\), the doubled-and-jammed interface class is
\[
\partial_Y(E_+(\pi(P_\pm)))\in K_1(C^*(Y\subset X)).
\]
For primitive straight periodic components, the Toeplitz pairing gives winding signs \(-1\) for
\(H_+^\triangle\) and \(+1\) for \(H_-^\triangle\).
\end{corollary}
\proofappendixref{app:proof:cor:self_contained_annular}




The FHS computation in \Cref{tab:periodic_triangular_certificate} independently confirms the Bloch
Chern signs in Theorem~\ref{M-thm:periodic_triangular_chern} on the displayed grids, and
\Cref{tab:periodic_triangular_extension_certificate} checks the finite-cover, folded-cell, and
fixed-positive-mass variants used below.

\begin{corollary}[Robust periodic annular input]\label{cor:robust_periodic_annular_input}
The conclusion of Corollary~\ref{cor:self_contained_annular} remains valid if the closed periodic
triangular bulk Hamiltonian on the doubled cylinders is replaced by one of the following
finite-range periodic bulk models:
\begin{enumerate}[label=(\roman*),leftmargin=2em]
\item \(H_{\pm,m}^\triangle(k)=B(k)+B(k)^*\pm mS(k)\), for any fixed \(m>0\);
\item \(H_{\pm,m}^\triangle(k)+V(k)\), where \(m>0\) is fixed, \(V\) is self-adjoint finite-range
periodic, and \(\sup_k\opnorm{V(k)}<g_{\pm,m}\) in the notation of
Corollary~\ref{M-cor:periodic_chern_stability};
\item \(H_\pm^\triangle(k)+V(k)\), where \(V\) is self-adjoint, finite-range periodic, and
\[
\sup_{k\in\mathbb T^2}\opnorm{V(k)}<2^{-13}.
\]
\item the finite-cover pullbacks and finite-index cell enlargements covered by
Corollary~\ref{cor:periodic_covers_sublattices}.
\end{enumerate}
The coarse interface collar \(Y\) is enlarged, if necessary, by the finite propagation radius of
the chosen periodic bulk model.
For each such model, the closed-bulk gap and the closed-bulk Chern sign are supplied internally by
Proposition~\ref{M-prop:periodic_positive_mass_homotopy} and
Corollary~\ref{M-cor:periodic_chern_stability}. Therefore, for every jam scale above the corresponding
copy-side positivity threshold, the doubled-and-jammed family has the interface class
\[
\partial_Y(E_+(\pi(P_\pm^{\mathrm{per}})))\in K_1(C^*(Y\subset X)),
\]
where \(P_\pm^{\mathrm{per}}\) is the negative Fermi projection of the chosen periodic closed-bulk
model.
This class is independent of the large scalar jam scale and is unchanged if the scalar jam is
replaced by any sufficiently positive bounded finite-propagation copy-side confining potential.
For primitive straight periodic components, the Toeplitz winding value is the corresponding
periodic Bloch Chern number supplied by Theorem~\ref{M-thm:periodic_triangular_chern} and
Corollary~\ref{cor:periodic_covers_sublattices}.
\end{corollary}
\proofappendixref{app:proof:cor:robust_periodic_annular_input}




\begin{corollary}[Periodic boundary families with internal bulk input]\label{cor:verified_periodic_families}
Let \((\K_n)_n\) be any sequence of triangular-lattice cylinders whose doubles are finite periodic
triangular tori, with arbitrary periods tending to infinity along a coarse disjoint union. The
following choices give boundary families for which the closed-bulk locality, gap, and Chern-sign
inputs are verified internally in this paper:
\begin{enumerate}[label=(\alph*),leftmargin=2em]
\item the unperturbed periodic triangular Hamiltonians \(H_\pm^\triangle\);
\item the positive-duality-strength Hamiltonians \(H_{\pm,m}^\triangle\), for any fixed \(m>0\);
\item the perturbed Hamiltonians \(H_{\pm,m}^\triangle+V\), for any fixed \(m>0\) and any
  self-adjoint finite-range periodic \(V\) satisfying
  \(\sup_k\opnorm{V(k)}<g_{\pm,m}\);
\item the explicit-radius perturbations \(H_\pm^\triangle+V\), for any self-adjoint finite-range
  periodic \(V\) satisfying \(\sup_k\opnorm{V(k)}<2^{-13}\);
\item finite-cover pullbacks of the triangular symbols, whose Chern numbers are multiplied by the
  signed covering degree, and finite-index cell enlargements, whose Chern signs are unchanged.
\end{enumerate}
For every family in (a)--(e), every sufficiently large scalar jam scale, every bounded
componentwise scalar jam scale bounded below by the corresponding positivity threshold, and every
sufficiently positive bounded finite-propagation copy-side confining potential produces the same
interface class
\[
\partial_Y(E_+(\pi(P_\pm^{\mathrm{per}})))\in K_1(C^*(Y\subset X)),
\]
after enlarging the coarse interface collar by the propagation radius of the selected periodic bulk
model and copy-side confinement. The class is also unchanged by altering finitely many initial
components of the cylinder sequence, or by passing through any norm-continuous sidewise homotopy
whose comparison gaps stay uniformly open. For primitive straight periodic components, its Toeplitz
winding value is the corresponding periodic Bloch Chern number from
Theorem~\ref{M-thm:periodic_triangular_chern} and
Corollary~\ref{cor:periodic_covers_sublattices}.
\end{corollary}
\proofappendixref{app:proof:cor:verified_periodic_families}




\begin{corollary}[Finite-cover primitive Toeplitz formula]\label{cor:covered_primitive_toeplitz}
Let \(H\) and \(P\) be as in Theorem~\ref{M-thm:periodic_toeplitz_pairing}, and let
\(A\in M_2(\mathbb Z)\) have \(\det A\ne0\). Set
\[
H_A(\kappa):=H(A\kappa),\qquad P_A(\kappa):=P(A\kappa).
\]
For every primitive tangential/positive-normal basis \((t,n)\) in the cover coordinates,
\[
\operatorname{wind}\!\left(\partial_{\mathrm T,t,n}([P_A])\right)
=
\det[t\ n]\,\det A\,\Ch(P).
\]
For the triangular symbols this gives
\[
\operatorname{wind}(H^\triangle_{+,A})=-\det[t\ n]\,\det A,
\qquad
\operatorname{wind}(H^\triangle_{-,A})=\det[t\ n]\,\det A,
\]
with the same formula for positive-mass symbols and for the small periodic perturbations covered by
Corollary~\ref{M-cor:periodic_chern_stability}.
\end{corollary}
\proofappendixref{app:proof:cor:covered_primitive_toeplitz}





\begin{corollary}[Gapped periodic homotopies and Toeplitz winding]\label{cor:gapped_periodic_toeplitz_homotopy}
Let \(s\mapsto H_s(k)\), \(0\le s\le1\), be a norm-continuous family of finite-range self-adjoint
periodic Bloch Hamiltonians with a common finite hopping range. Assume that
\[
\inf_{s,k}\dist\bigl(0,\spec H_s(k)\bigr)>0,
\qquad
P_s(k):=\chi_{(-\infty,0)}(H_s(k)).
\]
Then, for every primitive tangential/positive-normal basis \((t,n)\),
\[
\partial_{\mathrm T,t,n}([P_s])\in K_1(C(S^1,\mathcal K))
\]
is independent of \(s\), and
\[
\operatorname{wind}\!\left(\partial_{\mathrm T,t,n}([P_s])\right)
=
\det[t\ n]\,\Ch(P_0)
=
\det[t\ n]\,\Ch(P_1).
\]
Thus any finite-range periodic model connected to one of the internally verified triangular
families by such a zero-gapped homotopy has the same straight-edge Toeplitz winding.
\end{corollary}
\proofappendixref{app:proof:cor:gapped_periodic_toeplitz_homotopy}





\begin{theorem}[Split componentwise Toeplitz winding]\label{thm:split_component_toeplitz}
Let \(H\) and \(P\) be as in Theorem~\ref{M-thm:periodic_toeplitz_pairing}. Suppose a periodic
straight-edge model has finitely many interface components \(Y^{(1)},\ldots,Y^{(r)}\), and suppose
the edge ideal splits as a \(C^*\)-direct sum of the component ideals,
\[
C^*(Y\subset X)
\cong
\bigoplus_{j=1}^r C^*(Y^{(j)}\subset X),
\]
with the \(j\)-th summand identified with the Toeplitz edge ideal
\(C(S^1,\mathcal K)\) for a primitive tangential/normal lattice basis \((t_j,n_j)\). Assume also
that the quotient extension is the direct sum of these component Toeplitz extensions. Then the
global boundary class decomposes as
\[
\partial_Y([P])
=
\bigoplus_{j=1}^r \partial_{\mathrm T,t_j,n_j}([P])
\quad\text{in}\quad
K_1(C^*(Y\subset X)).
\]
Consequently the \(j\)-th component winding is
\[
\operatorname{wind}_j\bigl(\partial_Y([P])\bigr)
=
\det[t_j\ n_j]\,\Ch(P).
\]
If each component is oriented by the induced boundary orientation, the tangential coordinate is
reversed exactly when the normal sign is reversed, so each induced-orientation component winding
records the same bulk Chern number \(\Ch(P)\).
\end{theorem}
\proofappendixref{app:proof:thm:split_component_toeplitz}





\begin{theorem}[Tail split component winding maps]\label{thm:tail_split_component_windings}
Let \(X=\bigsqcup_n X_n\) be the coarse disjoint union of a straight-periodic boundary family from
Corollary~\ref{cor:verified_periodic_families}. Suppose that, after deleting finitely many initial
components, the tail interface ideal has either
\begin{enumerate}[label=(\roman*),leftmargin=2em]
\item a split direct-sum Toeplitz decomposition as in
  Theorem~\ref{thm:split_component_toeplitz}, or
\item a finite-corner or \(K_1\)-trivial-corner quotient decomposition as in
  Theorem~\ref{thm:finite_corner_toeplitz} and
  Corollary~\ref{cor:k1_trivial_corner_toeplitz}.
\end{enumerate}
Let \((t_j,n_j)\), \(1\le j\le r\), be the primitive tangential/positive-normal bases of the tail
components, and let \(\Omega_j\) be the component map obtained by deleting the finite initial
components, passing to the \(j\)-th Toeplitz quotient coordinate, and applying the usual winding
homomorphism. Then the maps \(\Omega_j\) are well defined on the original interface class and
\[
\Omega_j\!\left(\Indint(J;T^+,T^-)\right)
=
\det[t_j\ n_j]\,\Ch(P).
\]
For a finite-cover pullback by \(A\), the right-hand side is
\[
\det[t_j\ n_j]\,\det A\,\Ch(P).
\]
The same formulas are unchanged by arbitrary bounded finite-propagation changes on finitely many
initial components.
\end{theorem}
\proofappendixref{app:proof:thm:tail_split_component_windings}




\begin{corollary}[Internal straight-periodic classes]\label{cor:fully_verified_straight_periodic}
For the following boundary/interface families, with the split, quotient, tail, and sidewise-gap
hypotheses stated in the cited results, no Higson--Prodan closed-bulk input and no Ewert--Meyer
pairing assumption are used. Their zero gap, Chern value, interface class, and Toeplitz winding are
supplied by the internal periodic calculation and the Toeplitz evaluation above:
\begin{enumerate}[label=(\alph*),leftmargin=2em]
\item primitive straight periodic edges of the unperturbed triangular models
  \(H^\triangle_\pm\);
\item primitive straight periodic edges of \(H^\triangle_{\pm,m}\), for every fixed \(m>0\);
\item primitive straight periodic edges of \(H^\triangle_{\pm,m}+V\), for every fixed \(m>0\) and
  every self-adjoint finite-range periodic perturbation satisfying
  \(\sup_k\opnorm{V(k)}<g_{\pm,m}\); in particular this includes
  \(H^\triangle_\pm+V\) with \(\sup_k\opnorm{V(k)}<2^{-13}\);
\item finite-cover pullbacks of the families in (a)--(c), with Chern numbers multiplied by the
  signed covering degree, and finite-index cell enlargements of the same operators, with the
  original Chern signs unchanged;
\item any finite split direct sum of the families in (a)--(d), with component ideals and quotient
  extension splitting as in Theorem~\ref{thm:split_component_toeplitz};
\item any finite-corner or \(K_1\)-trivial-corner split sum of the families in (a)--(d), with
  quotient extension as in Theorem~\ref{thm:finite_corner_toeplitz},
  Corollary~\ref{cor:k1_trivial_corner_toeplitz}, or
  Corollary~\ref{cor:finite_corner_geometric_criterion};
\item finite-defect, bounded-distance straight-edge, and tail-split variants covered by
  Theorem~\ref{thm:rough_eventual_toeplitz_equivalence},
  Corollary~\ref{cor:finite_defect_tail_toeplitz}, and
  Theorem~\ref{thm:tail_split_component_windings};
\item the same families after replacing scalar jamming by any bounded componentwise scalar jam
  uniformly above the positivity threshold, or by any sufficiently positive bounded
  finite-propagation copy-side confinement;
\item the same families after adding norm-vanishing component defects, with the sidewise comparison
  gaps still open.
\end{enumerate}
For the triangular families before applying a finite cover, the winding values are \(-1\) for
\(H_+\) and \(+1\) for \(H_-\), up to the orientation factor \(\det[t_j\ n_j]\) on each primitive
component. For a
finite-cover pullback, these values are further multiplied by the signed covering degree. Fixed-collar
finite-propagation edge perturbations, norm-vanishing component defects, and norm-continuous
sidewise homotopies preserve these values whenever the sidewise comparison gaps remain open; this
last gap condition is a stability condition for the perturbation path, not an imported
bulk--interface theorem.
\end{corollary}
\proofappendixref{app:proof:cor:fully_verified_straight_periodic}





\subsection{Roe interface variants}\label{app:roe_interface_variants}
This subsection records stability, convention, relabeling, tail, and confinement variants for the
Roe interface \(K_1\)-class. Where an integer winding value appears, it refers to the
internal Toeplitz integer pairing in the straight periodic case
(\Cref{M-cor:toeplitz_winding}); for general partitions the integer pairing is the
conditional statement \Cref{M-cor:conditional_winding}.
\begin{corollary}[Small-norm stability of Fermi and interface classes]\label{cor:bulk_small_perturbation}\label{cor:positive_interface_small_perturbation}
Let $H=H^*\in C^*(X)$ satisfy
\[
\spec_{\mathrm{op}}(H)\subset[-B,-g]\cup[g,B]
\]
for some $B,g>0$. Let $V=V^*\in C^*(X)$ with
\[
\opnorm{V}<g.
\]
Then
\[
H+V
\]
is gapped at zero, with
\[
\spec_{\mathrm{op}}(H+V)\cap\bigl(-(g-\opnorm V),\,g-\opnorm V\bigr)=\emptyset,
\]
and
\[
\bigl[\chi_{(-\infty,0)}(H+V)\bigr]
=
\bigl[\chi_{(-\infty,0)}(H)\bigr]
\qquad\text{in }K_0(C^*(X)).
\]
If, in addition, \(H=H_\pm^\infty\) under Assumption~\ref{ass:uniform_bulk} and \(V=V^*\in C^*(X)\)
satisfies
\[
\opnorm{V}<g_{\mathrm{bulk}},
\]
then, with
\[
P_V:=\chi_{(-\infty,0)}(H_\pm^\infty+V).
\]
one has
\[
\partial_Y\!\bigl(E_+(\pi(P_V))\bigr)
=
\partial_Y\!\bigl(E_+(\pi(P_\pm))\bigr)
\qquad\text{in }K_1(C^*(Y\subset X)).
\]
\end{corollary}
\proofappendixref{app:proof:cor:bulk_small_perturbation}





\begin{proposition}[Interface-supported sidewise changes]\label{prop:interface_supported_invariance}\label{prop:sidewise_data_interface_invariance}
Let $(J;T^+,T^-)$ be a sidewise class-A interface along $Y$. Let
\[
J'=J+V,\qquad T^{\prime +}=T^+ + V^+,\qquad T^{\prime -}=T^-+V^-,
\]
where $V,V^+,V^-\in\mathcal I$ are self-adjoint. Assume that
$(J';T^{\prime +},T^{\prime -})$ is again a sidewise class-A interface along $Y$. Then
\[
p_{J'}=p_J,
\qquad
\Indint(J';T^{\prime +},T^{\prime -})=\Indint(J;T^+,T^-).
\]
The special case \(V^+=V^-=0\) gives invariance under perturbing only the interface Hamiltonian
\(J\) by a self-adjoint element of \(\mathcal I\).
\end{proposition}
\proofappendixref{app:proof:prop:sidewise_data_interface_invariance}




\begin{proposition}[Stability under finite-component changes]\label{prop:finite_component_stability}
Let $X=\bigsqcup_n X_n$ be the coarse disjoint union, and assume each $Y_n$ is nonempty. Let
$V=V^*$ be a bounded finite-propagation operator supported on only finitely many components: there
is $N<\infty$ such that
\[
V\chi_{X_n}=\chi_{X_n}V=0
\qquad\text{for all }n>N.
\]
Then $V\in C^*(Y\subset X)$. Consequently, changing any of the sidewise data
$J,T^+,T^-$ on finitely many initial components does not change the interface class, provided the
sidewise gap hypotheses remain valid after the change.
\end{proposition}
\proofappendixref{app:proof:prop:finite_component_stability}




\begin{proposition}[Norm-vanishing component defects]\label{prop:norm_vanishing_component_defects}
Let \(X=\bigsqcup_nX_n\) be the coarse disjoint union and assume each \(Y_n\) is nonempty. Let
\[
V=\bigoplus_{n=1}^\infty V_n
\]
be a bounded self-adjoint block-diagonal operator with \(V_n\in\mathcal B(\ell^2(X_n))\) and
\(\|V_n\|\to0\). Then
\[
V\in C^*(Y\subset X).
\]
Consequently, adding such norm-vanishing component defects to \(J,T^+\), or \(T^-\) does not change
the sidewise interface class, provided the resulting comparison operators remain gapped and the
sidewise class-A hypotheses remain valid.
\end{proposition}
\proofappendixref{app:proof:prop:norm_vanishing_component_defects}




\begin{proposition}[Homotopy invariance of the sidewise interface class]\label{prop:sidewise_homotopy}
Let
\[
s\longmapsto (J_s;T_s^+,T_s^-),
\qquad s\in[0,1],
\]
be a family of sidewise class-A interfaces along the same interface $Y$. Assume that
$s\mapsto T_s^\pm$ are norm-continuous and that there is a uniform comparison gap
\[
\dist(0,\spec_{\mathrm{op}}(T_s^\pm))\ge \gamma>0
\]
for all $s\in[0,1]$. Then
\[
\Indint(J_s;T_s^+,T_s^-)
\]
is independent of $s$.
\end{proposition}
\proofappendixref{app:proof:prop:sidewise_homotopy}




\begin{corollary}[Topological independence of the confining copy potential]\label{cor:copy_potential_independence}
Assume \ref{ass:uniform_bulk}. Let $s\mapsto V_s$, $s\in[0,1]$, be a norm-continuous family of bounded self-adjoint
finite-propagation operators on $\ell^2(X)$ such that
\[
V_s=P_{\mathrm{copy}}^\infty V_sP_{\mathrm{copy}}^\infty
\]
and, on the copy sector,
\[
V_s\ge M_0\Id
\qquad\text{for a constant }M_0>B_{\mathrm{copy},-}.
\]
Define
\[
J_s:=H_\pm^\infty+V_s,
\qquad
T_s^+:=H_\pm^\infty,
\qquad
T_s^-:=P_{\mathrm{orig}}^\infty+H_{\pm,\mathrm{copy}}^\infty+V_s.
\]
Then the sidewise interface class of $(J_s;T_s^+,T_s^-)$ is independent of $s$.
\end{corollary}
\proofappendixref{app:proof:cor:copy_potential_independence}




\begin{theorem}[Interface theorem for general confining copy potentials]\label{thm:general_copy_interface}
Assume \ref{ass:uniform_bulk}. Let $V=V^*$ be a bounded finite-propagation operator on $\ell^2(X)$
such that
\[
V=P_{\mathrm{copy}}^\infty V P_{\mathrm{copy}}^\infty
\]
and, on the copy sector,
\[
V\ge M_0\Id
\qquad\text{for some }M_0>B_{\mathrm{copy},-}.
\]
Define
\[
J_V:=H_\pm^\infty+V,
\qquad
T^+:=H_\pm^\infty,
\qquad
T_V^-:=P_{\mathrm{orig}}^\infty+H_{\pm,\mathrm{copy}}^\infty+V.
\]
Then $(J_V;T^+,T_V^-)$ is a sidewise class-A interface along $Y$. Its canonical quotient
projection is
\[
p_{J_V}=E_+\!\bigl(\pi(P_\pm)\bigr),
\]
the negative-side comparison has the quantitative positivity margin
\[
\dist\bigl(0,\spec_{\mathrm{op}}(T_V^-)\bigr)\ge \min\{1,M_0-B_{\mathrm{copy},-}\},
\]
and hence
\[
\Indint(J_V;T^+,T_V^-)
=
\partial_Y\!\bigl(E_+(\pi(P_\pm))\bigr)
\in K_1(C^*(Y\subset X)).
\]
In particular, the interface class is unchanged by replacing scalar jamming by any sufficiently
positive bounded finite-propagation copy-side confining potential.
\end{theorem}
\proofappendixref{app:proof:thm:general_copy_interface}





\begin{corollary}[Independence of the interface class for large jam scales]\label{cor:jam_independence}
Assume \ref{ass:uniform_bulk}. If $M_1,M_2>B_{\mathrm{copy},-}$, then
\[
\Indint(J_{\pm,M_1}^\infty;T_1^+,T_1^-)
=
\Indint(J_{\pm,M_2}^\infty;T_2^+,T_2^-)
=
\partial_Y\!\bigl(E_+(\pi(P_\pm))\bigr),
\]
where
\[
T_j^+=H_\pm^\infty,
\qquad
T_j^-=\widetilde T^-_{\pm,M_j}.
\]
\end{corollary}
\proofappendixref{app:proof:cor:jam_independence}





Combining Theorems~\ref{thm:bulkboundary_main} and~\ref{thm:bulk_chern} with
Corollary~\ref{cor:jam_independence}, we obtain the following explicit statement.

\begin{corollary}[Interface class from a Chern-$\pm1$ bulk projection]\label{cor:interface_from_chern}
Assume \ref{ass:uniform_bulk}, the additional Higson--Prodan hypotheses from
Theorem~\ref{thm:bulk_chern}, and $M>B_{\mathrm{copy},-}$. Then
\[
\Indint(J_{\pm,M}^\infty;T^+,T^-)
=
\partial_Y\!\bigl(E_+(\pi(P_\pm))\bigr)
\in K_1\!\bigl(C^*(Y\subset X)\bigr),
\]
where $P_\pm$ is the Higson--Prodan bulk Fermi projection with $\Ch([P_\pm])=\pm1$. The interface
class depends only on the bulk projection and the quotient decomposition, not on the jam scale $M$.
\end{corollary}
\proofappendixref{app:proof:cor:interface_from_chern}


\begin{corollary}[Higson--Prodan input implies the doubled-and-jammed interface class]\label{cor:hp_admissible_interface}
Assume that the closed doubled refinement family satisfies the Higson--Prodan closed-bulk input in
Assumption~\ref{ass:hp_admissibility}. Then
Assumptions~\ref{ass:S_uniform} and~\ref{ass:uniform_bulk} hold.
Consequently, for every $M>B_{\mathrm{copy},-}$,
\[
\Indint(J_{\pm,M}^\infty;T^+,T^-)
=
\partial_Y\!\bigl(E_+(\pi(P_\pm))\bigr)
\in K_1\!\bigl(C^*(Y\subset X)\bigr),
\]
where $P_\pm$ is the Higson--Prodan bulk Fermi projection and
\[
\Ch([P_\pm])=\pm1.
\]
The construction of this $K_1$-class is internal to the present paper; in the original/copy
partition the $K_1$-class identity is Theorem~\ref{M-thm:pathb_roe_winding}. The integer
winding evaluation is internal in the straight periodic Toeplitz case
(\Cref{M-cor:toeplitz_winding}) and is the conditional corollary
\Cref{M-cor:conditional_winding} for general partitions.
\end{corollary}
\proofappendixref{app:proof:cor:hp_admissible_interface}





\begin{proposition}[Ewert--Meyer convention criterion]\label{prop:hp_em_criterion}
Let \((\K_n)_n\) be a uniformly regular full-boundary refinement family. Suppose that the closed
doubled family \(D(\K_n)\) satisfies the Higson--Prodan closed-bulk input in
Assumption~\ref{ass:hp_admissibility}. Then the original/copy coarse data satisfy the
HP--Ewert--Meyer boundary convention of Definition~\ref{def:hp_em_admissible}, with the
Roe/Ewert--Meyer \(K_1\)-class identity supplied by Theorem~\ref{M-thm:pathb_roe_winding}.
No boundary spectral gap is required.
\end{proposition}
\proofappendixref{app:proof:prop:hp_em_criterion}




\begin{corollary}[HP-covered boundary refinement families]\label{cor:concrete_hp_em_classes}
Let \((\K_n)_n\) be a refinement sequence of a fixed compact oriented PL surface with full
boundary subcomplex. Suppose the closed doubles \(D(\K_n)\), after the standard rescaling used in
the closed Higson--Prodan refinement limit, are covered by the Higson--Prodan closed universal-model
theorem:
\begin{enumerate}[label=(\roman*),leftmargin=2em]
\item the closed doubles have the HP closed-bulk gap and Chern value;
\item the original/copy decomposition gives the coarse partition of
  Definition~\ref{def:hp_em_admissible}.
\end{enumerate}
Then the doubled-and-jammed original/copy family has the interface \(K_1\)-class of
Theorem~\ref{M-thm:pathb_roe_winding}. The following common refinement schemes are natural
HP-covered families when they satisfy the bounded-geometry hypotheses of the closed
theorem:
\begin{enumerate}[label=(\alph*),leftmargin=2em]
\item iterated barycentric refinements;
\item bounded-valence stellar or edge refinements;
\item quasi-uniform shape-regular triangulations used in finite-element refinements.
\end{enumerate}
For any such family, and every \(M>B_{\mathrm{copy},-}\), the doubled-and-jammed interface
\(K_1\)-class is given by Theorem~\ref{M-thm:pathb_roe_winding}; in the straight periodic
Toeplitz case the integer winding is \(\pm1\) internally (\Cref{M-cor:toeplitz_winding}),
and for general partitions the integer winding is the conditional statement
\Cref{M-cor:conditional_winding}.
\end{corollary}
\proofappendixref{app:proof:cor:concrete_hp_em_classes}




\begin{remark}[Only the tail of a refinement family matters]\label{cor:tail_admissibility}
Suppose that, after deleting finitely many initial refinements, the tail \((\K_n)_{n\ge N_0}\) is
HP-covered in the sense of Corollary~\ref{cor:concrete_hp_em_classes}. Extend the tail by any
finite set of initial components with bounded
local sidewise data for which the comparison operators are gapped at zero. Then the
doubled-and-jammed interface class and its Ewert--Meyer winding value are the same as those of the
tail. In particular, the HP-covered condition only has to hold eventually along a refinement
sequence; see \Cref{app:proof:cor:tail_admissibility} for the verification.
\end{remark}





\begin{remark}[Relabeling and orientation reversal]\label{cor:relabel_orientation}
The HP--Ewert--Meyer boundary convention, the interface class, and the Ewert--Meyer winding value
are invariant under orientation-preserving boundary-compatible simplicial isomorphisms of
refinement families. If the surface orientation is reversed, then the roles of \(H_+\) and \(H_-\)
are interchanged and the winding sign is reversed; see
\Cref{app:proof:cor:relabel_orientation}.
\end{remark}




\begin{corollary}[Robustness of the interface class and winding]\label{cor:interface_winding_robustness}
Under the hypotheses of Theorem~\ref{thm:full_hp_bec}, the interface \(K_1\)-class is
unchanged by the following operations, whenever the stated sidewise comparison gaps remain
open. In the straight periodic Toeplitz case (\Cref{M-cor:toeplitz_winding}), the resulting
integer winding is also unchanged.
\begin{enumerate}[label=(\alph*),leftmargin=2em]
\item replacing the scalar jam scale \(M\) by any other scalar \(M'>B_{\mathrm{copy},-}\);
\item replacing scalar jamming by a bounded finite-propagation copy-side confinement \(V\)
  satisfying
  \[
  V=P_{\mathrm{copy}}^\infty V P_{\mathrm{copy}}^\infty,\qquad
  V\ge M_0\Id,\qquad M_0>B_{\mathrm{copy},-};
  \]
  this includes bounded componentwise scalar jam scales \(V=\bigoplus_n M_nP_{\mathrm{copy},n}\)
  with \(B_{\mathrm{copy},-}<\inf_nM_n\le\sup_nM_n<\infty\);
\item changing the sidewise data by self-adjoint interface-supported operators;
\item changing the boundary/copy coupling by uniformly bounded finite-propagation perturbations
  supported in any fixed boundary collar;
\item replacing the chosen fixed collar model for \(Y\) by another fixed-radius collar;
\item changing the operators on finitely many initial refinement components;
\item adding norm-vanishing component defects;
\item passing through a norm-continuous sidewise homotopy with a uniform comparison gap.
\end{enumerate}
\end{corollary}
\proofappendixref{app:proof:cor:interface_winding_robustness}






\subsection{Symbolic verification appendix}\label{app:symbolic_verification}
The certificate is implemented in the companion script
\texttt{periodic\_triangular\_bulk\_certificate.py} in the project repository
(\href{https://github.com/yspennstate/Doubling-and-Jamming-A-Universal-Chern-Hamiltonian-on-Triangulated-Surfaces-with-Boundary}{GitHub project repository}).
The script numerically evaluates the closed-bulk Chern number on the triangular fundamental
cell at four grid sizes and writes the certificate table
\texttt{periodic\_triangular\_certificate\_table.tex}. The minimal SymPy snippet below is
illustrative: it builds the Bloch boundary operator, the cap-product blocks, and the
Hermitian duality block \(S(k)\) using the same convention as the companion script
(\(S[4:6,0:1]=\tfrac{i}{2}(P_2^*+P_0)\), \(S[1:4,1:4]=\tfrac{i}{2}(P_1^*-P_1)\),
\(S[0:1,4:6]=-\tfrac{i}{2}(P_0^*+P_2)\)), and prints the Hermitian Hamiltonians at a sample
point. The Chern-sign verification itself uses the numerical Fukui--Hatsugai--Suzuki
procedure in the companion script, which produces the integer table reproduced as
\Cref{tab:periodic_triangular_certificate} below.
\begin{verbatim}
# Minimal symbolic check. Builds B(k), P_p(k), S(k) and the two Hermitian
# Bloch Hamiltonians H_pm(k) = B(k) + B(k)^* +- S(k). Same conventions as
# the companion script periodic_triangular_bulk_certificate.py.
import sympy as sp
q1, q2 = sp.symbols('q1 q2', real=True)
i = sp.I

# Ordered basis: v, a, b, c, f, g.
z1 = sp.exp(i*q1)
z2 = sp.exp(i*q2)
d1 = sp.Matrix([[z1-1, z2-1, z1*z2-1]])
d2 = sp.Matrix([[1, -z2], [z1, -1], [-1, 1]])

# Bloch boundary operator B(k) (degree-lowering).
B = sp.zeros(6)
B[0:1, 1:4] = d1
B[1:4, 4:6] = d2

# Cap-product blocks (same convention as the companion script).
P0 = sp.Matrix([[sp.conjugate(z1*z2)], [sp.conjugate(z2)]])
P1 = sp.Matrix([[0, sp.conjugate(z1), 0], [0, 0, 0],
                [-sp.conjugate(z2), 0, 0]])
P2 = sp.Matrix([[1, 1]])

# Hermitian Bloch duality block S(k), single consistent assignment per block:
#   S[4:6, 0:1] =  (i/2) * (P2^* + P0)
#   S[1:4, 1:4] =  (i/2) * (P1^* - P1)
#   S[0:1, 4:6] = -(i/2) * (P0^* + P2)
S = sp.zeros(6)
S[4:6, 0:1] = (i/2) * (P2.conjugate().T + P0)
S[1:4, 1:4] = (i/2) * (P1.conjugate().T - P1)
S[0:1, 4:6] = (-i/2) * (P0.conjugate().T + P2)

H_plus  = B + B.conjugate().T + S
H_minus = B + B.conjugate().T - S

# Hermiticity sanity check at a sample point.
sample = {q1: sp.Rational(1, 7), q2: sp.Rational(2, 9)}
herm_plus = sp.simplify(H_plus.subs(sample) - H_plus.subs(sample).H)
print("H_+(k) Hermiticity defect at sample k:")
print(herm_plus)
\end{verbatim}

\noindent\emph{Note on this listing.} The script
\texttt{periodic\_triangular\_bulk\_certificate.py} in the project repository
(\href{https://github.com/yspennstate/Doubling-and-Jamming-A-Universal-Chern-Hamiltonian-on-Triangulated-Surfaces-with-Boundary}{GitHub project repository}) is the load-bearing implementation: it runs the
Fukui--Hatsugai--Suzuki Chern computation on a finite \(k\)-grid and writes the certificate
table (\Cref{tab:periodic_triangular_certificate}). The SymPy snippet above is an
illustrative excerpt showing the convention; it builds the Hermitian Bloch matrices
\(H_\pm(k)\) and is short enough to run unmodified.
The proof of the Chern signs is the analytic harmonic-compression and
crossing-orientation argument in \Cref{app:proof:thm:periodic_triangular_chern}, not the
symbolic computation.

The proof of each result in the main text is collected here in the order in which the results appear.

\subsection{Proof of Result~\ref{M-prop:gap_from_S}}\label{app:proof:prop:gap_from_S}
\begin{proof}[Proof of \Cref{M-prop:gap_from_S}]
Expand $H_\pm^2$ and use the assumed anticommutation to eliminate the cross term. Since $(B+B^\dagger)^2\succeq0$ and $S^2\succeq \mu_S^2\Id$, we obtain $H_\pm^2\succeq \mu_S^2\Id$.
\end{proof}

\subsection{Proof of Result~\ref{M-prop:gap_defect}}\label{app:proof:prop:gap_defect}
\begin{proof}[Proof of \Cref{M-prop:gap_defect}]
Since
\[
H_\pm^2=Q^2+S^2\pm(QS+SQ),
\]
we have $Q^2\succeq0$ and $S^2\succeq \mu_S^2\Id$. Also,
\[
\pm(QS+SQ)\succeq -\opnorm{QS+SQ}\,\Id=-\e_{\mathrm{ac}}\Id.
\]
Hence $H_\pm^2\succeq (\mu_S^2-\e_{\mathrm{ac}})\Id$. The spectral-gap conclusion follows immediately.
\end{proof}

\subsection{Proof of Result~\ref{M-prop:double_closed}}\label{app:proof:prop:double_closed}
\begin{proof}[Proof of \Cref{M-prop:double_closed}]
We verify that $D(\K)$ is a well-defined abstract simplicial complex. A simplex is determined by its
vertex set, so we must check that no two distinct simplices (of any dimension) in the doubled
complex share the same vertex set.

\emph{Faces.} Let $F=\{v_0,v_1,v_2\}$ be a face of $\K$. By the fullness condition in
Definition~\ref{M-def:surfwithbdry}, $F$ has at least one interior vertex $v_0$. Every reflected face
contains at least one copy-side vertex $\bar w\in\overline\K_0^{\mathrm{int}}$, and no original face
contains any copy-side vertex. Hence no original face coincides with any reflected face. Since faces
within each copy are distinct, all $2$-simplices of $D(\K)$ have distinct vertex sets.

\emph{Edges.} Interior edges of $\K$ have at least one interior endpoint (by fullness: both
endpoints on $\bd\K$ would force the edge to be a boundary edge). Hence every reflected interior
edge acquires at least one copy-side vertex and has a different vertex set from the original. Boundary
edges are identified (not duplicated), so they appear once. Thus all $1$-simplices of $D(\K)$ have
distinct vertex sets.

\emph{Edge valences.} Every interior edge of $\K$ is incident to exactly two faces in $\K$. Its
reflected copy (with a different vertex set) is incident to exactly two reflected faces. So each
such pair contributes two edges, each incident to two faces. Every boundary edge is incident to one
original face and one reflected face, hence to exactly two faces. Thus every edge of $D(\K)$ is
incident to exactly two faces.

At an interior vertex of either copy, the link is already a cycle. At a boundary vertex $v$, the link in $\K$ is a path and the link in $\overline\K$ is a second path with the opposite boundary orientation. After the boundary identification these two paths glue end-to-end to form a cycle. Thus every vertex link in $D(\K)$ is a cycle, so $D(\K)$ is a closed triangulated surface.

Finally, orient the faces in the original copy by the given surface orientation and orient the faces in the reflected copy by the opposite orientation. Along every glued boundary edge, the two induced edge orientations are opposite, so the face orientations are compatible across the gluing. Hence $D(\K)$ is oriented.
\end{proof}

\subsection{Proof of Result~\ref{M-thm:counts}}\label{app:proof:thm:counts}
\begin{proof}[Proof of \Cref{M-thm:counts}]
Faces are duplicated and never identified (each reflected face has a distinct vertex set by the
fullness condition), so $|D(\K)_2|=2|\K_2|$. Every non-boundary edge has at least one interior endpoint
(by fullness), so its reflected copy has a distinct vertex set; non-boundary edges are therefore
duplicated. Boundary edges are identified, not duplicated. Hence
$|D(\K)_1|=2|\K_1^{\mathrm{int}}|+|\bd\K_1|=2(|\K_1|-|\bd\K_1|)+|\bd\K_1|=2|\K_1|-|\bd\K_1|$.
The same reasoning applies to vertices: interior vertices are duplicated, boundary vertices are
identified, giving $|D(\K)_0|=2|\K_0|-|\bd\K_0|$. Since $\bd\K$ is a disjoint union of cycles, each
connected component has as many vertices as edges, so $|\bd\K_0|=|\bd\K_1|$.
\end{proof}

\subsection{Proof of Result~\ref{cor:dim_double}}\label{app:proof:cor:dim_double}
\begin{proof}[Proof of \Cref{cor:dim_double}]
Use
\[
\dim\Hil(\K)=|\K_0|+|\K_1|+|\K_2|,
\qquad
\dim\Hil(D(\K))=|D(\K)_0|+|D(\K)_1|+|D(\K)_2|,
\]
insert Theorem~\ref{M-thm:counts}, and simplify with $|\bd\K_0|=|\bd\K_1|$.
\end{proof}

\subsection{Proof of Result~\ref{prop:euler_double}}\label{app:proof:prop:euler_double}
\begin{proof}[Proof of \Cref{prop:euler_double}]
Insert the identities from Theorem~\ref{M-thm:counts} into $\chi=|\K_0|-|\K_1|+|\K_2|$ and use $|\bd\K_0|=|\bd\K_1|$.
\end{proof}

\subsection{Proof of Result~\ref{cor:genus_double}}\label{app:proof:cor:genus_double}
\begin{proof}[Proof of \Cref{cor:genus_double}]
For an orientable surface with boundary, $\chi(\K)=2-2g-b$ \cite{Hatcher}. By Proposition~\ref{prop:euler_double},
\[
\chi(D(\K))=4-4g-2b.
\]
Since a closed orientable surface of genus $G$ has Euler characteristic $2-2G$, one gets $G=2g+b-1$.
\end{proof}

\subsection{Proof of Result~\ref{M-lem:relative_cap_identity}}\label{app:proof:lem:relative_cap_identity}
\begin{proof}[Proof of \Cref{M-lem:relative_cap_identity}]
For a \(p\)-cochain \(\alpha\), identified with a \(p\)-chain by the simplex basis, the standard
cap-product boundary formula \cite[Section~3.3]{Hatcher} is
\[
\partial(c\cap\alpha)=(-1)^p\bigl((\partial c)\cap\alpha-c\cap\delta\alpha\bigr).
\]
Taking \(c=[\K,\bd\K]\), whose boundary is zero in the relative chain complex, gives
\eqref{eq:relative_cap_chain_identity}. Applying \(j_{1-p}\) to that identity and subtracting
from \(B_{\K,2-p}j_{2-p}P^{\mathrm{rel}}_p\) gives \eqref{eq:lifted_cap_defect}.

The operator \(B_{\K,r}j_r-j_{r-1}\partial^{\mathrm{rel}}_r\) is zero on every relative simplex
whose codimension-one faces are all represented by non-boundary simplices. When it is nonzero, its
output is a boundary simplex, and its input simplex meets \(\bd\K\). Since the cap-product maps
\(P^{\mathrm{rel}}_p\) have a fixed finite propagation radius, a nonzero coefficient in
\eqref{eq:lifted_cap_defect} can occur only with source and target inside a uniformly bounded
collar of \(\bd\K\). The operator \(S^{\mathrm{PL}}_\K\) is obtained from \(\widehat P\) and
\(\widehat P^\dagger\) by fixed degree phases and Hermitian symmetrization; taking adjoints
preserves the same support condition. Hence the displayed duality defect is collar-supported.
\end{proof}

\subsection{Proof of Result~\ref{M-lem:cap_coefficient_agreement}}\label{app:proof:lem:cap_coefficient_agreement}
\begin{proof}[Proof of \Cref{M-lem:cap_coefficient_agreement}]
With the coefficient convention in Section~\ref{sec:Sclosed}, a matrix coefficient of either
cap-product operator is a signed count of oriented faces for which the source simplex appears as
the prescribed suffix and the target simplex appears as the complementary prefix. If both simplices
are outside the \(R_\frown\)-collar, every face contributing to the coefficient is an original-side
face disjoint from the glued boundary. On such faces the closed fundamental cycle of \(D(\K)\)
restricts exactly to the relative fundamental chain \([\K,\bd\K]\), and the quotient map \(q\)
does not remove the complementary prefix. Hence the two signed coefficient sums are identical.
\end{proof}

\subsection{Proof of Result~\ref{M-thm:relative_boundary_comparison}}\label{app:proof:thm:relative_boundary_comparison}
\begin{proof}[Proof of \Cref{M-thm:relative_boundary_comparison}]
Self-adjointness is built into the Hermitian symmetrization. Locality follows from the chain-level
cap product: a matrix coefficient of \(P^{\mathrm{rel}}\) can be nonzero only when the two
simplices are incident inside a common face or star. The representative map \(j\) removes boundary
representatives but does not enlarge support, so \(S^{\mathrm{PL}}_\K\), and hence
\(H^{\mathrm{PL}}_{\pm}(\K)\), has finite propagation.

The relative duality defect is collar-supported by Lemma~\ref{M-lem:relative_cap_identity}.

The \(B+B^\dagger\) part of \(H_{\pm}^{\bd}(\K)\) agrees exactly with \(B_\K+B_\K^\dagger\).
Compressing the doubled boundary operator to the original side keeps the original faces, edges,
and vertices and removes only reflected simplices not contained in \(\bd\K\). Thus the difference
between \(H_{\pm}^{\bd}(\K)\) and \(H^{\mathrm{PL}}_{\pm}(\K)\) is entirely in the cap-product
term.

Lemma~\ref{M-lem:cap_coefficient_agreement} gives coefficient-by-coefficient equality between the
doubled cap-product operator, compressed to the original side, and the lifted relative cap-product
operator away from a fixed boundary collar. Thus any discrepancy can occur only when the local
cap-product calculation involves a simplex meeting \(\bd\K\), or when the relative quotient has
discarded a boundary representative. Since the cap-product rule has a fixed finite propagation
radius, all such discrepancies have both source and target inside a fixed boundary collar. Defining
\[
R_{\partial,\pm}:=
H_{\pm}^{\bd}(\K)-H^{\mathrm{PL}}_{\pm}(\K)
\]
gives \eqref{M-eq:relative_boundary_comparison}. The difference of two self-adjoint operators is
self-adjoint, and the collar-support statement follows from the preceding support analysis. The
uniform statement follows because the cap-product propagation radius is fixed over a uniformly
regular refinement family.
\end{proof}

\subsection{Proof of Result~\ref{cor:pl_doubled_quotient}}\label{app:proof:cor:pl_doubled_quotient}
\begin{proof}[Proof of \Cref{cor:pl_doubled_quotient}]
The difference on the \(n\)-th component is \(R_{\partial,\pm,n}\), supported in the fixed boundary
collar from Theorem~\ref{M-thm:relative_boundary_comparison}. The direct sum has uniformly bounded
propagation and uniformly bounded matrix coefficients by uniform regularity and the fixed local
cap-product convention. Hence it belongs to the algebraic interface ideal, and therefore to
\(C^*(Y\subset X)\). The quotient equality follows immediately. The same theorem gives the fixed
collar support of the relative duality defect, so its direct sum is in the same ideal. Finally,
quotient maps commute with continuous functional calculus for bounded self-adjoint elements.
\end{proof}

\subsection{Proof of Result~\ref{M-lem:collar_rank_counts}}\label{app:proof:lem:collar_rank_counts}
\begin{proof}[Proof of \Cref{M-lem:collar_rank_counts}]
Since \(R=P_WRP_W\), \(\rank R\le r\). Weyl's rank inequality for spectral distribution
functions \cite[Chapter~III]{BhatiaMatrixAnalysis} gives
\[
N_A(t)-\rank R\le N_{A+R}(t)\le N_A(t)+\rank R,
\]
which proves the first estimate. The interval estimate follows by subtracting the distribution
functions at \(a\) and \(b\). The final statement is Theorem~\ref{M-thm:relative_boundary_comparison}.
\end{proof}

\subsection{Proof of Result~\ref{M-prop:fixed_collar_negligible}}\label{app:proof:prop:fixed_collar_negligible}
\begin{proof}[Proof of \Cref{M-prop:fixed_collar_negligible}]
Let \(h_n\to0\) be the quasi-uniform mesh size. Uniform regularity gives, for each fixed
\(R\), a constant \(C_R\) such that every radius-\(R\) ball in the barycentric graph contains at
most \(C_R\) simplices. Hence the \(R\)-collar of the boundary is covered by radius-\(R\) balls
centered at boundary vertices and boundary edges, so
\[
r_n(R)
\le C_R\bigl(|(\bd\K_n)_0|+|(\bd\K_n)_1|\bigr).
\]
Since \(\bd\K_n\) is a disjoint union of cycles, \(|(\bd\K_n)_0|=|(\bd\K_n)_1|\). Boundary edge
lengths are comparable to \(h_n\), while the total boundary length of the fixed PL surface is
finite and independent of \(n\). Therefore
\[
|(\bd\K_n)_1|=O(h_n^{-1}),
\qquad
r_n(R)=O(h_n^{-1}).
\]
On the other hand, shape regularity gives face areas comparable to \(h_n^2\), while the total
surface area is fixed and positive on every nondegenerate component. Thus
\[
|(\K_n)_2|=\Theta(h_n^{-2}).
\]
Uniform incidence bounds give \(|(\K_n)_0|+|(\K_n)_1|+|(\K_n)_2|=\Theta(|(\K_n)_2|)\), and hence
\[
d_n=\dim\Hil(\K_n)=|(\K_n)_0|+|(\K_n)_1|+|(\K_n)_2|
=\Theta(h_n^{-2}).
\]
It follows that
\[
\frac{r_n(R)}{d_n}=O(h_n)\longrightarrow0 .
\]
Taking \(R=R_{\mathrm{PL}}\) is legitimate because
\Cref{M-thm:relative_boundary_comparison} supplies a single fixed collar radius over uniformly
regular refinement families.
\end{proof}

\subsection{Proof of Result~\ref{cor:pl_bulk_spectral_equivalence}}\label{app:proof:cor:pl_bulk_spectral_equivalence}
\begin{proof}[Proof of \Cref{cor:pl_bulk_spectral_equivalence}]
Apply Lemma~\ref{M-lem:collar_rank_counts} with
\(W=\ell^2(N_{R_{\mathrm{PL}}}(\bd\K_n))\), divide by \(d_n\), and take the supremum in \(t\).
For quasi-uniform fixed-surface two-dimensional refinements,
\Cref{M-prop:fixed_collar_negligible} gives \(r_n/d_n\to0\); outside such geometric refinement
regimes it should be checked as an explicit hypothesis. The statement about continuous test
functions follows because the operators have uniformly bounded norm under uniform regularity, so
the empirical spectral measures are supported in a common compact interval and the
distribution-function estimate forces the same weak subsequential limits.
\end{proof}

\subsection{Proof of Result~\ref{M-lem:schur}}\label{app:proof:lem:schur}
\begin{proof}[Proof of \Cref{M-lem:schur}]
Since $D$ is self-adjoint and $\Im z\neq0$, the operator $D-z$ is invertible. Standard block Gaussian elimination gives the Schur-complement formula.
\end{proof}

\subsection{Proof of Result~\ref{M-thm:resolvent}}\label{app:proof:thm:resolvent}
\begin{proof}[Proof of \Cref{M-thm:resolvent}]
By Lemma~\ref{M-lem:schur},
\[
P_{\mathrm{orig}}(H_{\pm,M}-z)^{-1}P_{\mathrm{orig}}
=\bigl(A-z-K_M(z)\bigr)^{-1},
\qquad
K_M(z)=C(D_M-z)^{-1}C^\dagger.
\]
Hence
\[
\opnorm{K_M(z)}\le \frac{\opnorm{C}^2}{d_M(z)}.
\]
Now set $R_0=(A-z)^{-1}$ and $R_M=(A-z-K_M(z))^{-1}$. The resolvent identity gives
\[
R_M-R_0=R_MK_M(z)R_0.
\]
Moreover $R_M$ is the $(1,1)$ block of $(H_{\pm,M}-z)^{-1}$, so
\[
\opnorm{R_M}\le \opnorm{(H_{\pm,M}-z)^{-1}}\le \frac1{|\Im z|},
\qquad
\opnorm{R_0}\le \frac1{|\Im z|}.
\]
Combining these bounds yields \eqref{M-eq:resolvent_bound_exact}. Since
\[
\spec(D_M)\subset [M-\|D\|,M+\|D\|],
\]
we also have $d_M(z)\ge M-\|D\|-|\Re z|$ whenever $M>\|D\|+|\Re z|$, which gives \eqref{eq:resolvent_bound}.
\end{proof}

\subsection{Proof of Result~\ref{M-cor:boundary_hamiltonian_limit}}\label{app:proof:cor:boundary_hamiltonian_limit}
\begin{proof}[Proof of \Cref{M-cor:boundary_hamiltonian_limit}]
This is the first case in \Cref{app:proof:schur_routine_corollaries}.
\end{proof}

\subsection{Proof of Result~\ref{M-cor:full_resolvent}}\label{app:proof:cor:full_resolvent}
\begin{proof}[Proof of \Cref{M-cor:full_resolvent}]
This is the second case in \Cref{app:proof:schur_routine_corollaries}.
\end{proof}

\subsection{Routine Schur-complement corollaries}\label{app:proof:schur_routine_corollaries}
\begin{proof}
For \Cref{M-cor:boundary_hamiltonian_limit}, apply Theorem~\ref{M-thm:resolvent} with
\(A=H_{\pm}^{\bd}(\K)\).

For \Cref{M-cor:full_resolvent}, set \(R_D(z)=(D_M-z)^{-1}\) and
\(S_M(z)=A-z-CR_D(z)C^\dagger\). The block inverse formula gives
\[
(H_{\pm,M}-z)^{-1}=
\begin{pmatrix}
S_M(z)^{-1}&-S_M(z)^{-1}CR_D(z)\\[0.4em]
-R_D(z)C^\dagger S_M(z)^{-1}&R_D(z)+R_D(z)C^\dagger S_M(z)^{-1}CR_D(z)
\end{pmatrix}.
\]
The upper-left block is controlled by Theorem~\ref{M-thm:resolvent}. The off-diagonal blocks have norm
at most \(\opnorm{C}/(|\Im z|\,d_M(z))\), while the lower-right block is bounded by
\[
\frac1{d_M(z)}+\frac{\opnorm{C}^2}{d_M(z)^2|\Im z|}.
\]
Adding these four bounds gives \Cref{M-cor:full_resolvent}.
\end{proof}

\subsection{Proof of Result~\ref{M-thm:first_order_eff}}\label{app:proof:thm:first_order_eff}
\begin{proof}[Proof of \Cref{M-thm:first_order_eff}]
Write
\[
D_M-z=M\Bigl(\Id+\frac{D-z}{M}\Bigr).
\]
If $M>\|D-z\|$, the Neumann series gives
\[
\Bigl(\Id+\frac{D-z}{M}\Bigr)^{-1}
=
\Id-\frac{D-z}{M}+\Bigl(\frac{D-z}{M}\Bigr)^2-\cdots,
\]
which yields the stated expansion and remainder estimate. Multiplying on both sides by $C$ and $C^\dagger$ gives the formula for $K_M(z)$.

Set
\[
R_M=P_{\mathrm{orig}}(H_{\pm,M}-z)^{-1}P_{\mathrm{orig}}=(A-z-K_M(z))^{-1},
\qquad
R_M^{(1)}=(A_M^{(1)}-z)^{-1}.
\]
Then
\[
R_M-R_M^{(1)}=R_M\,\widetilde R_M(z)\,R_M^{(1)}.
\]
Since both $H_{\pm,M}$ and $A_M^{(1)}$ are self-adjoint,
\[
\opnorm{R_M}\le \frac1{|\Im z|},
\qquad
\opnorm{R_M^{(1)}}\le \frac1{|\Im z|}.
\]
Combining these bounds with the estimate on $\widetilde R_M(z)$ gives \eqref{M-eq:first_order_bound}.
\end{proof}

\subsection{Proof of Result~\ref{M-thm:proj_general}}\label{app:proof:thm:proj_general}
\begin{proof}[Proof of \Cref{M-thm:proj_general}]
For $t\in[0,1]$ define the homotopy
\[
H_{t,M}:=
\begin{pmatrix}
A&tC\\
tC^\dagger&D_M
\end{pmatrix}.
\]
Fix $z\in\Gamma$. Since $\delta_D(\Gamma)>0$, the operator $D_M-z$ is invertible. The Schur complement of $H_{t,M}-z$ with respect to the lower-right block is
\[
S_{t,M}(z)=A-z-t^2C(D_M-z)^{-1}C^\dagger.
\]
Because $z\in\Gamma$, we have
\[
\opnorm{(A-z)^{-1}}\le \frac{1}{\delta_A(\Gamma)},
\qquad
\opnorm{(D_M-z)^{-1}}\le \frac{1}{\delta_D(\Gamma)}.
\]
Hence
\[
\opnorm{(A-z)^{-1}t^2C(D_M-z)^{-1}C^\dagger}
\le
\frac{\opnorm{C}^2}{\delta_A(\Gamma)\,\delta_D(\Gamma)}
<1
\]
by the hypothesis $\delta_H(\Gamma)>0$. Therefore
\[
S_{t,M}(z)
=
(A-z)\Bigl[\Id-(A-z)^{-1}t^2C(D_M-z)^{-1}C^\dagger\Bigr]
\]
is invertible, because the bracketed term is invertible by the Neumann series. Since $D_M-z$ is invertible as well, the block matrix $H_{t,M}-z$ is invertible. This proves $\Gamma\subset\rho(H_{t,M})$ for all $t\in[0,1]$, and hence $\Gamma\subset\rho(H_{\pm,M})$.

The map $t\mapsto P_\Gamma(H_{t,M})$ is norm-continuous, so its rank is constant on the connected interval $[0,1]$. At $t=0$,
\[
H_{0,M}=A\oplus D_M,
\]
and therefore
\[
P_\Gamma(H_{0,M})=P_\Gamma(A)\oplus P_\Gamma(D_M).
\]
This proves the rank identity.

For the norm estimate, abbreviate
\[
K_M(z):=C(D_M-z)^{-1}C^\dagger,
\qquad
S_M(z):=A-z-K_M(z).
\]
The Schur-complement inverse formula gives
Set \(R_D(z)=(D_M-z)^{-1}\). The block inverse formula gives
\[
(H_{\pm,M}-z)^{-1}=
\begin{pmatrix}
S_M(z)^{-1}&-S_M(z)^{-1}CR_D(z)\\[0.4em]
-R_D(z)C^\dagger S_M(z)^{-1}&R_D(z)+R_D(z)C^\dagger S_M(z)^{-1}CR_D(z)
\end{pmatrix}.
\]
Because $\delta_H(\Gamma)>0$,
\[
\opnorm{S_M(z)^{-1}}\le \frac{1}{\delta_H(\Gamma)},
\qquad
\opnorm{(A-z)^{-1}}\le \frac{1}{\delta_A(\Gamma)},
\qquad
\opnorm{(D_M-z)^{-1}}\le \frac{1}{\delta_D(\Gamma)}.
\]
Also,
\[
S_M(z)^{-1}-(A-z)^{-1}=S_M(z)^{-1}K_M(z)(A-z)^{-1},
\]
so
\[
\opnorm{S_M(z)^{-1}-(A-z)^{-1}}
\le
\frac{\opnorm{C}^2}{\delta_A(\Gamma)\,\delta_H(\Gamma)\,\delta_D(\Gamma)}.
\]
The off-diagonal blocks satisfy
\[
\opnorm{S_M(z)^{-1}C(D_M-z)^{-1}}
\le
\frac{\opnorm{C}}{\delta_H(\Gamma)\,\delta_D(\Gamma)},
\]
and the extra lower-right block is bounded by
\[
\opnorm{(D_M-z)^{-1}C^\dagger S_M(z)^{-1}C(D_M-z)^{-1}}
\le
\frac{\opnorm{C}^2}{\delta_H(\Gamma)\,\delta_D(\Gamma)^2}.
\]
Subtracting the block-diagonal resolvent
\[
\begin{pmatrix}
(A-z)^{-1}&0\\[0.2em]
0&(D_M-z)^{-1}
\end{pmatrix}
\]
gives the pointwise bound \eqref{eq:proj_general_resolvent_bound}. Integrating around $\Gamma$
then gives \eqref{M-eq:proj_general_bound}.
\end{proof}

\subsection{Proof of Result~\ref{M-thm:spectral_localization}}\label{app:proof:thm:spectral_localization}
\begin{proof}[Proof of \Cref{M-thm:spectral_localization}]
Let $\psi=(\psi_{\mathrm{orig}},\psi_{\mathrm{copy}})$ be a normalized eigenvector with eigenvalue $E$. If $\psi_{\mathrm{orig}}=0$, then $(D_M)\psi_{\mathrm{copy}}=E\psi_{\mathrm{copy}}$, which is impossible because $|E|\le E_{\max}$ and $g_M(E_{\max})>0$. Thus $\psi_{\mathrm{orig}}\neq0$.

Solving the second block equation gives
\[
\psi_{\mathrm{copy}}=-(D_M-E)^{-1}C^\dagger\psi_{\mathrm{orig}}.
\]
Substituting into the original-block equation yields
\[
\bigl(A-\Sigma_M(E)\bigr)\psi_{\mathrm{orig}}=E\psi_{\mathrm{orig}},
\qquad
\Sigma_M(E)=C(D_M-E)^{-1}C^\dagger.
\]
Hence $E\in\spec(A-\Sigma_M(E))$. Since
\[
\opnorm{\Sigma_M(E)}
\le
\opnorm{C}^2\,\opnorm{(D_M-E)^{-1}}
\le \frac{\opnorm{C}^2}{g_M(E_{\max})},
\]
Weyl's inequality implies \eqref{eq:specA_exact}. If $M>\opnorm{D}+E_{\max}$, then
\[
g_M(E_{\max})\ge M-\opnorm{D}-E_{\max},
\]
which gives \eqref{eq:specA_crude}.
\end{proof}

\subsection{Proof of Result~\ref{M-cor:relative_spectral_localization}}\label{app:proof:cor:relative_spectral_localization}
\begin{proof}[Proof of \Cref{M-cor:relative_spectral_localization}]
By Theorem~\ref{M-thm:relative_boundary_comparison}, the original-side block is
\(A=H^{\mathrm{PL}}_{\pm}(\K)+R_{\partial,\pm}\). The result is therefore
Theorem~\ref{M-thm:spectral_localization} with this identification of \(A\).
\end{proof}

\subsection{Proof of Result~\ref{M-thm:first_order_spectral_localization}}\label{app:proof:thm:first_order_spectral_localization}
\begin{proof}[Proof of \Cref{M-thm:first_order_spectral_localization}]
Let $\psi=(\psi_{\mathrm{orig}},\psi_{\mathrm{copy}})$ be a normalized eigenvector of $H_{\pm,M}$ with eigenvalue $E$, where $|E|\le E_{\max}$. As in the proof of Theorem~\ref{M-thm:spectral_localization},
\[
\bigl(A-\Sigma_M(E)\bigr)\psi_{\mathrm{orig}}=E\psi_{\mathrm{orig}},
\qquad
\Sigma_M(E)=C(D_M-E)^{-1}C^\dagger.
\]
Now
\[
(D_M-E)^{-1}-\frac1M\Id
=
-\frac1M(D-E)(D_M-E)^{-1},
\]
so
\[
\opnorm{(D_M-E)^{-1}-\frac1M\Id}
\le
\frac{\opnorm{D-E}}{M}\,\opnorm{(D_M-E)^{-1}}
\le
\frac{\opnorm{D}+E_{\max}}{M\,g_M(E_{\max})}.
\]
Hence
\[
\Sigma_M(E)=\frac1MCC^\dagger+\widetilde\Sigma_M(E),
\qquad
\opnorm{\widetilde\Sigma_M(E)}
\le
\frac{\opnorm{C}^2(\opnorm{D}+E_{\max})}{M\,g_M(E_{\max})}.
\]
Therefore
\[
A-\Sigma_M(E)=A_M^{(1)}-\widetilde\Sigma_M(E).
\]
Since $E\in\spec(A_M^{(1)}-\widetilde\Sigma_M(E))$, Weyl's inequality yields \eqref{M-eq:specA1_exact}. The crude bound \eqref{M-eq:specA1_crude} follows from $g_M(E_{\max})\ge M-\opnorm{D}-E_{\max}$ when $M>\opnorm{D}+E_{\max}$.
\end{proof}

\subsection{Proof of Result~\ref{M-thm:eigvec_leak}}\label{app:proof:thm:eigvec_leak}
\begin{proof}[Proof of \Cref{M-thm:eigvec_leak}]
From the second block equation,
\[
(D_M-E)\psi_{\mathrm{copy}}=-C^\dagger\psi_{\mathrm{orig}}.
\]
Because $g_M(E_{\max})>0$, the operator $D_M-E$ is invertible for every $|E|\le E_{\max}$, so
\[
\|\psi_{\mathrm{copy}}\|
\le
\opnorm{(D_M-E)^{-1}}\,\opnorm{C}\,\|\psi_{\mathrm{orig}}\|
\le
\frac{\opnorm{C}}{g_M(E_{\max})}\,\|\psi_{\mathrm{orig}}\|.
\]
If $M>\opnorm{D}+E_{\max}$, then $g_M(E_{\max})\ge M-\opnorm{D}-E_{\max}$, which gives \eqref{eq:leak_crude}.
\end{proof}

\subsection{Proof of Result~\ref{M-thm:inertia_reduction}}\label{app:proof:thm:inertia_reduction}
\begin{proof}[Proof of \Cref{M-thm:inertia_reduction}]
Since \(D_M>0\), block Gaussian elimination gives the congruence
\[
\begin{pmatrix}
\Id&-CD_M^{-1}\\
0&\Id
\end{pmatrix}
\begin{pmatrix}
A&C\\ C^\dagger&D_M
\end{pmatrix}
\begin{pmatrix}
\Id&0\\
-D_M^{-1}C^\dagger&\Id
\end{pmatrix}
=
\begin{pmatrix}
S_M(0)&0\\
0&D_M
\end{pmatrix}.
\]
Sylvester's law of inertia gives the first two identities because the lower-right block is
positive. The norm estimate follows from
\(\opnorm{D_M^{-1}}\le(M-\opnorm D)^{-1}\). If this norm is strictly smaller than
\(\delta_A\), then the straight-line path \(A-tCD_M^{-1}C^\dagger\), \(0\le t\le1\), remains
invertible by Weyl's inequality, so its negative spectral count is constant. Combining this with
the first identity proves the final claim.
\end{proof}

\subsection{Proof of Result~\ref{M-cor:scalar_jam_monotonicity}}\label{app:proof:cor:scalar_jam_monotonicity}
\begin{proof}[Proof of \Cref{M-cor:scalar_jam_monotonicity}]
One has
\[
H_{\pm,M_2}-H_{\pm,M_1}=(M_2-M_1)P_{\mathrm{copy}}\succeq0.
\]
The min-max principle therefore gives the monotonicity of each ordered eigenvalue. The spectral
count statement is the same assertion written for distribution functions.
\end{proof}

\subsection{Proof of Result~\ref{M-cor:no_copy_spectral_subspace}}\label{app:proof:cor:no_copy_spectral_subspace}
\begin{proof}[Proof of \Cref{M-cor:no_copy_spectral_subspace}]
Write \(\psi=\psi_{\mathrm{orig}}\oplus\psi_{\mathrm{copy}}\). Since
\(\psi\in\Ran Q_M\), the spectral theorem gives \(\|H_{\pm,M}\psi\|\le E_{\max}\|\psi\|\).
Taking the copy block of \(H_{\pm,M}\psi\) yields
\[
D_M\psi_{\mathrm{copy}}+C^\dagger\psi_{\mathrm{orig}}
=P_{\mathrm{copy}}H_{\pm,M}\psi .
\]
Because \(D_M\ge (M-\opnorm D)\Id\),
\[
(M-\opnorm D)\|\psi_{\mathrm{copy}}\|
\le E_{\max}\|\psi\|+\opnorm C\,\|\psi_{\mathrm{orig}}\|,
\]
which proves the quantitative estimate. If \(\psi_{\mathrm{orig}}=0\), the same block equation
gives
\[
(M-\opnorm D)\|\psi_{\mathrm{copy}}\|\le E_{\max}\|\psi_{\mathrm{copy}}\|.
\]
The strict inequality \(M-\opnorm D>E_{\max}\) forces \(\psi_{\mathrm{copy}}=0\), so
\(\psi=0\). Thus \(P_{\mathrm{orig}}\) is injective on \(\Ran Q_M\), and the rank bound follows.
\end{proof}

\subsection{Proof of Result~\ref{M-prop:inherited_gap_A}}\label{app:proof:prop:inherited_gap_A}
\begin{proof}[Proof of \Cref{M-prop:inherited_gap_A}]
Suppose, toward a contradiction, that $E\in\spec(H_{\pm,M})$ and
$|E|<g_A-\Delta_M$. Then $|E|\le g_A$, so Theorem~\ref{M-thm:spectral_localization} with
$E_{\max}=g_A$ gives
\[
\dist(E,\spec(A))\le \Delta_M.
\]
Since $\spec(A)\cap(-g_A,g_A)=\emptyset$ and $|E|<g_A$, we also have
\[
\dist(E,\spec(A))\ge g_A-|E|>\Delta_M,
\]
a contradiction.
\end{proof}

\subsection{Proof of Result~\ref{M-cor:uniform_direct_sum_jam}}\label{app:proof:cor:uniform_direct_sum_jam}
\begin{proof}[Proof of \Cref{M-cor:uniform_direct_sum_jam}]
Apply Theorems~\ref{M-thm:resolvent}, \ref{M-thm:first_order_eff},
\ref{M-thm:spectral_localization}, and~\ref{M-thm:eigvec_leak} componentwise and take suprema over
\(n\), using the bounds \(\|C_n\|\le\bar C\) and \(\|D_n\|\le\bar D\). The full resolvent
convergence follows similarly from Corollary~\ref{M-cor:full_resolvent}. Since the resolvent
functions \((t-z)^{-1}\) generate \(C_0(\mathbb R)\) as a \(C^*\)-algebra, the \(C_0\)-functional
calculus convergence follows from the norm convergence of the resolvents.
\end{proof}

\subsection{Proof of Result~\ref{M-thm:internal_bulk_spectral_equivalence}}\label{app:proof:thm:internal_bulk_spectral_equivalence}
\begin{proof}[Proof of \Cref{M-thm:internal_bulk_spectral_equivalence}]
Set
\[
\bar C=\sup_n\|C_{\pm,n}\|,\qquad \bar D=\sup_n\|D_{\pm,n}\|.
\]
Since \(M_n\to\infty\), the full block-resolvent estimate in
Corollary~\ref{M-cor:full_resolvent}, applied with the uniform bounds \(\bar C,\bar D\), gives
\[
\|(H_{\pm,n,M_n}-z)^{-1}-((A_{\pm,n}-z)^{-1}\oplus0)\|\longrightarrow0
\]
for every \(z\in\mathbb C\setminus\mathbb R\). The algebra generated by the resolvent functions is
dense in \(C_0(\mathbb R)\), so
\[
\|f(H_{\pm,n,M_n})-(f(A_{\pm,n})\oplus0)\|\longrightarrow0
\qquad(f\in C_0(\mathbb R)).
\]
The doubled Hilbert space has dimension at most \(2d_n\). Therefore
\[
\left|
\frac1{d_n}\operatorname{Tr} f(H_{\pm,n,M_n})
-
\frac1{d_n}\operatorname{Tr} f(A_{\pm,n})
\right|
\le
2\,\|f(H_{\pm,n,M_n})-(f(A_{\pm,n})\oplus0)\|
\longrightarrow0 .
\]

It remains to compare \(A_{\pm,n}\) with \(H^{\mathrm{PL}}_{\pm}(\K_n)\). By
Theorem~\ref{M-thm:relative_boundary_comparison}, their difference is
\(R_{\partial,\pm,n}=P_{W_n}R_{\partial,\pm,n}P_{W_n}\) for a collar subspace \(W_n\) of
dimension \(r_n\). Lemma~\ref{M-lem:collar_rank_counts} gives
\[
\sup_{t\in\mathbb R}
\left|
\frac1{d_n}N_{A_{\pm,n}}(t)
-
\frac1{d_n}N_{H^{\mathrm{PL}}_{\pm}(\K_n)}(t)
\right|
\le
\frac{r_n}{d_n}\longrightarrow0 .
\]
The two sequences of operators have uniformly bounded norm by hypothesis, so their normalized
spectral measures are supported in a common compact interval. Uniform convergence of the
distribution functions up to an \(o(1)\) error implies that the integrals of every continuous
function on this compact interval have the same limit; applying this to the restriction of
\(f\in C_0(\mathbb R)\) proves
\[
\frac1{d_n}\operatorname{Tr} f(A_{\pm,n})
-
\frac1{d_n}\operatorname{Tr} f(H^{\mathrm{PL}}_{\pm}(\K_n))
\longrightarrow0 .
\]
Combining the two limits gives the claimed normalized-trace equivalence.
\end{proof}

\subsection{Proof of Result~\ref{M-prop:periodic_triangular_charpoly}}\label{app:proof:prop:periodic_triangular_charpoly}
\begin{proof}[Proof of \Cref{M-prop:periodic_triangular_charpoly}]
Substituting the displayed matrices for \(B(k)\) and \(S(k)\) and expanding
\(\det(\lambda\Id-H_\pm^\triangle(k))\) gives the stated factorization. Equivalently, with
\(z_j=e^{ik_j}\), clearing denominators first factors the characteristic polynomial into
\[
4\lambda^2-F(k)
\]
times
\[
4\lambda^4-26\lambda^2+F(k)\pm4\lambda G(k),
\]
and division by the leading coefficient gives the monic form above.
\end{proof}

\subsection{Proof of Result~\ref{M-thm:periodic_triangular_bulk}}\label{app:proof:thm:periodic_triangular_bulk}
\begin{proof}[Proof of \Cref{M-thm:periodic_triangular_bulk}]
Proposition~\ref{M-prop:periodic_triangular_charpoly} gives
\[
\det H_\pm^\triangle(k)
=
-\frac1{16}
\left(
-28+8\cos k_1+6\cos k_2+6\cos(k_1+k_2)
\right)^2.
\]
Since each cosine is at most \(1\),
\[
-28+8\cos k_1+6\cos k_2+6\cos(k_1+k_2)\le -8,
\]
so the determinant is at most \(-4\). The absolute row sums of \(H_\pm^\triangle(k)\) are bounded
by \(8\), so \(\|H_\pm^\triangle(k)\|\le8\). If \(\lambda_1,\ldots,\lambda_6\) are the eigenvalues,
then \(|\lambda_j|\le8\) and
\[
4\le|\det H_\pm^\triangle(k)|
\le
\dist(0,\spec H_\pm^\triangle(k))\,8^5,
\]
which gives the stated gap. At \(k=0\), the characteristic polynomial is
\[
(\lambda^2-2)(2\lambda^4-13\lambda^2+4)/2,
\]
so exactly three eigenvalues are negative. The gap prevents the rank from changing over
\(\mathbb T^2\).

The same matrices satisfy \(B(k)S(k)+S(k)B(k)^*=0\). The stronger anticommutator
\((B(k)+B(k)^*)S(k)+S(k)(B(k)+B(k)^*)\) is not used here.
\end{proof}

\subsection{Proof of Result~\ref{M-cor:periodic_small_perturbation_gap}}\label{app:proof:cor:periodic_small_perturbation_gap}
\begin{proof}[Proof of \Cref{M-cor:periodic_small_perturbation_gap}]
Theorem~\ref{M-thm:periodic_triangular_bulk} gives
\[
\dist(0,\spec H_\pm^\triangle(k))\ge 2^{-13}
\qquad(k\in\mathbb T^2).
\]
Weyl's inequality gives
\[
\dist\bigl(0,\spec(H_\pm^\triangle(k)+V(k))\bigr)
\ge
2^{-13}-\opnorm{V(k)}>0.
\]
The negative spectral projection therefore varies continuously in \(k\), and its rank is constant.
At \(V=0\) this rank is \(3\), so it remains \(3\) along the straight-line homotopy
\(H_\pm^\triangle(k)+sV(k)\), \(0\le s\le1\).
\end{proof}

\subsection{Proof of Result~\ref{M-prop:periodic_positive_mass_homotopy}}\label{app:proof:prop:periodic_positive_mass_homotopy}
\begin{proof}[Proof of \Cref{M-prop:periodic_positive_mass_homotopy}]
The determinant identity is an exact computation from the displayed Bloch matrices. Since
\(c_1+c_2+c_{12}\le3\) and \(2+c_2+c_{12}\ge0\), the bracket is nonnegative. If the first term
vanishes, then \(c_1=c_2=c_{12}=1\), and the second term equals \(8m^2>0\). Thus the bracket is
positive for every \(m>0\), so the determinant never vanishes. The path
\((m,k)\mapsto H_{\pm,m}^\triangle(k)\) is norm-continuous on each compact interval
\(m\in[m_0,m_1]\subset(0,\infty)\).
\end{proof}

\subsection{Proof of Result~\ref{M-thm:periodic_triangular_chern}}\label{app:proof:thm:periodic_triangular_chern}
\begin{proof}[Proof of \Cref{M-thm:periodic_triangular_chern}]
The proof uses the orientation conventions fixed in Section~\ref{sec:orientation_signs}. The
Brillouin torus is oriented by \(dk_1\wedge dk_2\), the crossing parameter space is oriented by
\(dq_1\wedge dq_2\wedge d\mu\), and Chern numbers are computed as
\[
\Ch(P)=\frac{1}{2\pi i}\int_{\mathbb T^2}\tr(P\,dP\wedge dP).
\]
The parameter \(\mu\) is the coefficient of \(S\) in
\[
H_\mu^\triangle(k)=B(k)+B(k)^*+\mu S(k).
\]
Proposition~\ref{M-prop:periodic_positive_mass_homotopy} shows that \(H_\mu^\triangle(k)\) is
gapped at zero for \(\mu\ne0\), so the Chern number can change only at the crossing
\((k,\mu)=(0,0)\).

\paragraph{Step 1: harmonic compression.}
At \(k=0\), the boundary matrices are
\[
\partial_1(0)=\begin{bmatrix}0&0&0\end{bmatrix},\qquad
\partial_2(0)=\begin{bmatrix}1&-1\\1&-1\\-1&1\end{bmatrix}.
\]
Thus \(\operatorname{rank}\partial_1(0)=0\) and \(\operatorname{rank}\partial_2(0)=1\). The
kernel of the Hodge--Dirac operator \(D_0=B(0)+B(0)^*\) is the direct sum of the harmonic parts in
degrees \(0,1,2\): one vertex vector, the two-dimensional orthogonal complement of
\(\operatorname{im}\partial_2(0)\subset\mathbb C\langle a,b,c\rangle\), and the kernel of
\(\partial_2(0)\) in degree \(2\). With the ordered basis
\(v,a,b,c,f,g\), choose the orthonormal harmonic basis
\[
h_0=v,\qquad
h_1=\frac{a-b}{\sqrt2},\qquad
h_2=\frac{a+b+2c}{\sqrt6},\qquad
h_3=\frac{f+g}{\sqrt2}.
\]
Let \(U:\mathbb C^4\to\mathcal H(0)\) be the isometry \(Ue_j=h_{j-1}\), and let
\[
\Pi_{\mathrm h}=UU^*
\]
be the harmonic projector. The complement \((1-\Pi_{\mathrm h})\mathcal H(0)\) is separated from
zero for \(D_0\), because \(D_0\) is self-adjoint and has no kernel on that finite-dimensional
complement.

Put \(q_1=k_1\), \(q_2=k_2\). Since \(z_j=e^{iq_j}=1+iq_j+O(|q|^2)\), the first-order harmonic
compression is
\[
L(q_1,q_2,\mu)
=\Pi_{\mathrm h}\bigl(q_1\partial_{q_1}H_0^\triangle(0)+q_2\partial_{q_2}H_0^\triangle(0)
+\mu\partial_\mu H_\mu^\triangle(0)|_{\mu=0}\bigr)\Pi_{\mathrm h}\big|_{\mathcal H(0)}.
\]
Equivalently,
\[
\langle h_r,L(q_1,q_2,\mu)h_s\rangle
=i\bigl(q_1X_1+q_2X_2+\mu X_3\bigr)_{r+1,s+1},
\]
where the real skew-symmetric matrices are
\[
X_1=
\begin{bmatrix}
0&1/\sqrt2&\sqrt{3/2}&0\\
-1/\sqrt2&0&0&-1/2\\
-\sqrt{3/2}&0&0&1/(2\sqrt3)\\
0&1/2&-1/(2\sqrt3)&0
\end{bmatrix},
\]
\[
X_2=
\begin{bmatrix}
0&-1/\sqrt2&\sqrt{3/2}&0\\
1/\sqrt2&0&0&-1/2\\
-\sqrt{3/2}&0&0&-1/(2\sqrt3)\\
0&1/2&1/(2\sqrt3)&0
\end{bmatrix},
\qquad
X_3=
\begin{bmatrix}
0&0&0&-\sqrt2\\
0&0&-1/\sqrt3&0\\
0&1/\sqrt3&0&0\\
\sqrt2&0&0&0
\end{bmatrix}.
\]
These entries are obtained by differentiating the displayed Bloch matrices in
Section~\ref{M-sec:periodicbulk} and applying the matrix elements
\(\langle h_r,(\partial H)h_s\rangle\). The factor \(i\) is forced by Hermiticity: the derivatives
of the unitary phases at \(1\) are imaginary, while \(X_j\) are real skew matrices.

It remains to justify that this compression has the same local Chern jump as the full crossing.
Write the full local operator, relative to
\(\Pi_{\mathrm h}\mathcal H\oplus(1-\Pi_{\mathrm h})\mathcal H\), as
\[
\begin{pmatrix}
A(x)&K(x)\\ K(x)^*&G(x)
\end{pmatrix},\qquad x=(q_1,q_2,\mu).
\]
Here \(A(x)=L(x)+O(|x|^2)\), \(K(x)=O(|x|)\), and \(G(0)\) is invertible. Shrinking to
\(|x|\le r\), \(G(x)\) remains invertible and
\[
A(x)-K(x)G(x)^{-1}K(x)^*=L(x)+O(|x|^2).
\]
On the sphere \(|x|=r\), the linear model below is invertible with a lower bound \(c r\) for some
\(c>0\). Choose \(r\) so that the quadratic error has norm at most \(cr/2\). Then the straight-line
homotopy from the Schur complement to \(L(x)\) is invertible on the sphere, and the complementary
block remains invertible throughout. Therefore the full crossing and the displayed harmonic model
have the same linking Chern number.

\paragraph{Step 2: self-dual and anti-self-dual splitting.}
Identify a real skew \(4\times4\) matrix with a two-form by the ordered basis
\[
e_{12},e_{13},e_{14},e_{23},e_{24},e_{34}
\]
for \(\Lambda^2\mathbb R^4\), where \(e_{rs}=e_r\wedge e_s\) and
\(e_1\wedge e_2\wedge e_3\wedge e_4\) is the positive orientation. The Hodge involution is
\[
*e_{12}=e_{34},\quad *e_{13}=-e_{24},\quad *e_{14}=e_{23},
\]
together with the same equations read backwards. Hence the self-dual basis is
\[
\omega_1^+=\frac{e_{12}+e_{34}}{\sqrt2},\quad
\omega_2^+=\frac{e_{13}-e_{24}}{\sqrt2},\quad
\omega_3^+=\frac{e_{14}+e_{23}}{\sqrt2},
\]
and the anti-self-dual basis is
\[
\omega_1^-=\frac{e_{12}-e_{34}}{\sqrt2},\quad
\omega_2^-=\frac{e_{13}+e_{24}}{\sqrt2},\quad
\omega_3^-=\frac{e_{14}-e_{23}}{\sqrt2}.
\]
For
\[
X(q_1,q_2,\mu)=q_1X_1+q_2X_2+\mu X_3,
\]
write \(X=X_+(u)+X_-(\beta)\) in these bases. Direct projection onto the \(\omega_j^+\) coordinates
gives
\[
u(q_1,q_2,\mu)=
\left(
\frac{1+1/\sqrt6}{2}(q_1-q_2),
\left(\frac{\sqrt3}{2}+\frac{1}{2\sqrt2}\right)(q_1+q_2),
-(1+1/\sqrt6)\mu
\right).
\]
The anti-self-dual vector is a scalar multiple of the same coordinate vector,
\[
\beta(q_1,q_2,\mu)=\rho\,u(q_1,q_2,\mu),\qquad
\rho=\frac{\sqrt6-1}{\sqrt6+1}<1.
\]
For a skew form with self-dual and anti-self-dual coordinates \((\alpha,\beta)\), the Hermitian
matrix \(iX(\alpha,\beta)\) has eigenvalues
\[
\pm(|\alpha|+|\beta|),\qquad \pm(|\alpha|-|\beta|).
\]
Along the homotopy \((\alpha,\beta)=(u,t\rho u)\), \(0\le t\le1\), the smallest absolute eigenvalue
is \((1-t\rho)|u|\). This is positive on every linking sphere because \(\rho<1\) and
\(u\ne0\) there. Thus the anti-self-dual part can be removed without changing the local Chern
jump.

The self-dual coordinate map is linear with Jacobian
\[
\frac{\partial u}{\partial(q_1,q_2,\mu)}=
\begin{bmatrix}
a&-a&0\\ b&b&0\\0&0&-c
\end{bmatrix},
\quad
a=\frac{1+1/\sqrt6}{2},\quad
b=\frac{\sqrt3}{2}+\frac{1}{2\sqrt2},\quad
c=1+\frac1{\sqrt6}.
\]
Therefore
\[
\det\frac{\partial u}{\partial(q_1,q_2,\mu)}=-2abc<0.
\]
With the crossing orientation \(dq_1\wedge dq_2\wedge d\mu\), the self-dual local model is
orientation-reversing.

\paragraph{Step 3: Chern jump and FHS sign convention.}
Under the standard spinor identification of self-dual two-forms, \(iX_+(u)\) is the direct sum of
two identical Pauli Hamiltonians \(u\cdot\sigma\). For the identity crossing
\(u=(q_1,q_2,\mu)\), the negative line over the linking sphere has Chern number \(+1\) with the
boundary orientation induced by \(dq_1\wedge dq_2\wedge d\mu\). Pullback by a linear map multiplies
this number by the degree of that map. Since \(\det Du<0\), one Pauli block contributes \(-1\), and
two identical blocks contribute
\[
\Ch(E_-(H^\triangle_{\mu>0}))-\Ch(E_-(H^\triangle_{\mu<0}))=-2.
\]
The FHS convention used in
\Cref{tab:periodic_triangular_certificate,tab:periodic_triangular_extension_certificate} orders each
plaquette by the positively oriented coordinate square in the \(k_1,k_2\) directions. Its lattice
integer therefore approximates the same
\((2\pi i)^{-1}\int\tr(P\,dP\wedge dP)\) used here; no sign conversion is inserted between the
proof and the tables.

Finally,
\[
H^\triangle_{-\mu}(k)=\overline{H^\triangle_{\mu}(-k)}.
\]
The base involution \(k\mapsto-k\) has Jacobian determinant \(+1\) on the two-torus, so it preserves
the orientation \(dk_1\wedge dk_2\). Complex conjugation of a complex vector bundle reverses the
first Chern class. Hence
\[
\Ch(E_-(H^\triangle_{-\mu}))=-\Ch(E_-(H^\triangle_{\mu})).
\]
Combining this symmetry with the jump \(-2\) gives
\[
\Ch(E_-(H^\triangle_{\mu}))=-1\quad(\mu>0),\qquad
\Ch(E_-(H^\triangle_{\mu}))=+1\quad(\mu<0).
\]
Taking \(\mu=1\) gives the stated signs for \(H^\triangle_+\) and \(H^\triangle_-\).
\end{proof}

\subsection{Proof of Result~\ref{M-cor:periodic_chern_stability}}\label{app:proof:cor:periodic_chern_stability}
\begin{proof}[Proof of \Cref{M-cor:periodic_chern_stability}]
The positive-mass statement follows from Proposition~\ref{M-prop:periodic_positive_mass_homotopy}
and Theorem~\ref{M-thm:periodic_triangular_chern}: changing \(m>0\) gives a norm-continuous gapped
homotopy for \(H_{+,m}^\triangle\), and \(H_{-,m}^\triangle=H_{-m}^\triangle\) is handled by the
same argument on the negative-mass side. Thus the Chern number is unchanged. The number
\(g_{\pm,m}\) is positive by compactness of \(\mathbb T^2\) and the positive-mass gap. If
\(\|V\|_\infty<g_{\pm,m}\), Weyl's inequality keeps the straight-line homotopy
\(H_{\pm,m}^\triangle+sV\) gapped at zero, so the Chern number is unchanged along it.

For the uniform perturbation statement at \(m=1\), the straight-line
homotopy
\[
H_{\pm,s}^\triangle(k)=H_\pm^\triangle(k)+sV(k),\qquad 0\le s\le1,
\]
is gapped by Theorem~\ref{M-thm:periodic_triangular_bulk} and Weyl's inequality:
\[
\dist\bigl(0,\spec H_{\pm,s}^\triangle(k)\bigr)
\ge
2^{-13}-s\sup_{q\in\mathbb T^2}\opnorm{V(q)}>0 .
\]
The Chern number of the negative spectral projection is constant along this gapped homotopy.
\end{proof}

\subsection{Proof of Result~\ref{cor:periodic_covers_sublattices}}\label{app:proof:cor:periodic_covers_sublattices}
\begin{proof}[Proof of \Cref{cor:periodic_covers_sublattices}]
The map \(p_A:\mathbb T^2\to\mathbb T^2\), \(p_A(\kappa)=A\kappa\), is a finite covering of degree
\(\det A\) with sign \cite{Hatcher}. The pulled-back Hamiltonian has the same fiber spectra as the original one,
so the gap is unchanged. Its negative Bloch bundle is \(p_A^*E_-(H_\pm^\triangle)\), and naturality
of the Chern form gives
\[
\Ch(p_A^*E)=\deg(p_A)\Ch(E)=\det A\,\Ch(E).
\]
The displayed signs follow from Theorem~\ref{M-thm:periodic_triangular_chern}; the positive-mass and
small-perturbation statements follow in the same way from
Corollary~\ref{M-cor:periodic_chern_stability}.

For a finite-index sublattice \(\Lambda'\subset\mathbb Z^2\), the folded Bloch fiber over the
smaller Brillouin torus is the direct sum of the original fibers over the finitely many momenta
lying above it. The Chern form of this direct sum is the sum of the pulled-back Chern forms over
the sheets, and integration over the smaller Brillouin torus gives exactly the integral over the
original Brillouin torus. The gap is the minimum of the same original fiber gaps. Thus changing the
period cell does not change the Chern number of the underlying periodic operator.
\end{proof}

\subsection{Proof of Result~\ref{cor:finite_periodic_triangular_tori}}\label{app:proof:cor:finite_periodic_triangular_tori}
\begin{proof}[Proof of \Cref{cor:finite_periodic_triangular_tori}]
The finite periodic torus diagonalizes by the discrete Bloch transform. Its fiber momenta are
points of \(\mathbb T^2\), so the gap bound follows from Theorem~\ref{M-thm:periodic_triangular_bulk}.
The Chern data are thermodynamic Bloch-bundle data proved in
Theorem~\ref{M-thm:periodic_triangular_chern} and numerically checked in
\Cref{tab:periodic_triangular_certificate}; they are not finite-volume invariants of the discrete
sample alone.
\end{proof}

\subsection{Proof of Result~\ref{cor:self_contained_annular}}\label{app:proof:cor:self_contained_annular}
\begin{proof}[Proof of \Cref{cor:self_contained_annular}]
The exact gap, finite-range symbol, and analytic Chern calculation provide the locality, norm,
closed-bulk gap, and closed-bulk Chern values needed in
Assumptions~\ref{ass:S_uniform} and~\ref{ass:uniform_bulk}; here these inputs come from the
periodic calculation rather than from Assumption~\ref{ass:hp_admissibility}. With the gap and
sidewise inputs fixed, the sidewise and quotient-interface conclusions are exactly
Theorem~\ref{thm:bulkboundary_main}. For primitive straight periodic components, the final winding
evaluation is the Toeplitz calculation of Theorem~\ref{M-thm:periodic_toeplitz_pairing}, with the
periodic Chern numbers supplied by Theorem~\ref{M-thm:periodic_triangular_chern}.
\end{proof}

\subsection{Proof of Result~\ref{cor:robust_periodic_annular_input}}\label{app:proof:cor:robust_periodic_annular_input}
\begin{proof}[Proof of \Cref{cor:robust_periodic_annular_input}]
For \(H_{\pm,m}^\triangle\), the gap follows from
Proposition~\ref{M-prop:periodic_positive_mass_homotopy} and the Chern signs from
Corollary~\ref{M-cor:periodic_chern_stability}. The fixed-\(m\) perturbation case and the
\(m=1\) explicit-radius perturbation case are also exactly
Corollary~\ref{M-cor:periodic_chern_stability}. Finite-cover and finite-index cell choices are
Corollary~\ref{cor:periodic_covers_sublattices}. Finite-range periodicity gives the
same uniform locality and bounded-geometry hypotheses as in the unperturbed periodic annular
family, after replacing the fixed collar by the enlarged collar in the statement. The proofs of
Theorem~\ref{thm:bulkboundary_main}, Corollary~\ref{cor:jam_independence}, and
Theorem~\ref{thm:general_copy_interface} then apply verbatim with \(H_\pm^\infty\) replaced by the
chosen finite-range periodic bulk family: the only inputs used there are uniform locality,
boundedness, the closed-bulk gap, sidewise agreement away from the collar, and positivity of the
copy-side comparison. For primitive straight periodic components, the final winding statement is
the Toeplitz evaluation used in Corollary~\ref{cor:self_contained_annular}.
\end{proof}

\subsection{Proof of Result~\ref{cor:verified_periodic_families}}\label{app:proof:cor:verified_periodic_families}
\begin{proof}[Proof of \Cref{cor:verified_periodic_families}]
The unperturbed case is Corollary~\ref{cor:self_contained_annular}. The positive-mass,
small finite-range perturbation, finite-cover, and cell-enlargement cases are
Corollary~\ref{cor:robust_periodic_annular_input}. Independence
of the scalar jam scale is Corollary~\ref{cor:jam_independence}; bounded componentwise scalar jams
and nonconstant positive copy-side confinements are covered by
Theorem~\ref{thm:general_copy_interface} and
Corollary~\ref{cor:copy_potential_independence}. Changing finitely many initial components changes
the data by a finite-component perturbation, so Proposition~\ref{prop:finite_component_stability}
applies. Uniformly gapped sidewise homotopies are covered by
Proposition~\ref{prop:sidewise_homotopy}. The last sentence is the straight-periodic Toeplitz
evaluation already isolated in Corollaries~\ref{cor:self_contained_annular}
and~\ref{cor:robust_periodic_annular_input}.
\end{proof}

\subsection{Proof of Result~\ref{M-thm:periodic_toeplitz_pairing}}\label{app:proof:thm:periodic_toeplitz_pairing}
\begin{proof}[Proof of \Cref{M-thm:periodic_toeplitz_pairing}]
The proof is the standard Toeplitz calculation, included here so the periodic straight-edge family
is evaluated internally. The extension above is the spatial tensor product of \(C(S^1)\), in the
tangential variable \(k_1\), with the one-variable Toeplitz extension in the normal variable
\(k_2\). Hence its boundary map is a group homomorphism
\[
K^0(\mathbb T^2)\longrightarrow K^1(S^1).
\]
The group \(K^0(\mathbb T^2)\) is generated by the trivial line bundle and one Bott class
\cite{AtiyahKTheory,BlackadarKTheory}. A
constant projection lifts to a constant projection in \(\mathcal T_{\mathrm{per}}\), so the
boundary map vanishes on the rank summand.

It remains to compute the Bott generator. Use the clutching description of a Chern-one line
bundle: split the \(k_2\)-circle into upper and lower arcs, trivialize the bundle over the two
cylinders \(S^1_{k_1}\times I_\pm\), and take transition function \(z(k_1)=e^{ik_1}\) on one
overlap component and \(1\) on the other. The first Chern number is the winding difference of
these two transition functions, hence \(1\) with orientation \(dk_1\wedge dk_2\). In the Toeplitz
extension, lifting the projection to the two Toeplitz charts and comparing the lifts on the
overlap gives exactly this boundary unitary \(z\). Therefore
\[
\partial_{\mathrm T}([\mathrm{Bott}])=[z]\in K_1(C(S^1)),
\qquad
\operatorname{wind}(z)=1.
\]
By additivity and the standard \(K^0(\mathbb T^2)\) decomposition, any projection class decomposes as
\[
[P]=\operatorname{rank}(P)[1]+\Ch(P)[\mathrm{Bott}]
\quad\text{in }K^0(\mathbb T^2),
\]
and applying \(\partial_{\mathrm T}\) gives
\[
\partial_{\mathrm T}([P])=\Ch(P)[z].
\]
Pairing with the standard winding homomorphism on \(K_1(C(S^1))\) gives the claimed equality.
\end{proof}

\subsection{Proof of Result~\ref{M-cor:primitive_periodic_toeplitz}}\label{app:proof:cor:primitive_periodic_toeplitz}
\begin{proof}[Proof of \Cref{M-cor:primitive_periodic_toeplitz}]
The map \((\kappa,\eta)\mapsto \kappa t+\eta n\) is a torus automorphism. Pulling \(H\) back by
this automorphism preserves finite range, self-adjointness, and the zero gap, so
Theorem~\ref{M-thm:periodic_toeplitz_pairing} applies in the \((\kappa,\eta)\)-coordinates. The
first Chern number of the pulled-back projection is multiplied by the degree of the torus
automorphism, namely \(\det[t\ n]\). Therefore the coordinate Toeplitz winding equals
\(\varepsilon(t,n)\Ch(P)\).
\end{proof}

\subsection{Proof of Result~\ref{cor:covered_primitive_toeplitz}}\label{app:proof:cor:covered_primitive_toeplitz}
\begin{proof}[Proof of \Cref{cor:covered_primitive_toeplitz}]
The map \(p_A:\mathbb T^2\to\mathbb T^2\), \(p_A(\kappa)=A\kappa\), has signed degree
\(\det A\) \cite{Hatcher}, and \(P_A=p_A^*P\). Hence
\[
\Ch(P_A)=\det A\,\Ch(P).
\]
Applying Corollary~\ref{M-cor:primitive_periodic_toeplitz} to \(P_A\) in the primitive cover
coordinates gives the displayed formula. The triangular values are Theorem~\ref{M-thm:periodic_triangular_chern}
and Corollaries~\ref{M-cor:periodic_chern_stability} and~\ref{cor:periodic_covers_sublattices}.
\end{proof}

\subsection{Proof of Result~\ref{cor:gapped_periodic_toeplitz_homotopy}}\label{app:proof:cor:gapped_periodic_toeplitz_homotopy}
\begin{proof}[Proof of \Cref{cor:gapped_periodic_toeplitz_homotopy}]
The uniform gap gives a norm-continuous Riesz projection path \(s\mapsto P_s\) in
\(C(\mathbb T^2,M_N)\). Hence \([P_s]\in K_0(C(\mathbb T^2))\) is constant. Naturality of the
Toeplitz connecting map gives a constant class
\(\partial_{\mathrm T,t,n}([P_s])\), and Corollary~\ref{M-cor:primitive_periodic_toeplitz} evaluates
its winding. The common finite hopping range is the locality condition needed to realize the whole
homotopy by the same straight periodic edge extension.
\end{proof}

\subsection{Proof of Result~\ref{thm:split_component_toeplitz}}\label{app:proof:thm:split_component_toeplitz}
\begin{proof}[Proof of \Cref{thm:split_component_toeplitz}]
For a direct sum of extensions, the six-term exact sequence is natural and the connecting map is
the direct sum of the connecting maps of the summands. Hence the boundary map for the split edge
ideal sends the quotient projection class to
\[
\bigoplus_{j=1}^r \partial_{\mathrm T,t_j,n_j}([P]).
\]
The \(j\)-th coordinate projection
\[
K_1\!\left(\bigoplus_{m=1}^r C^*(Y^{(m)}\subset X)\right)
\longrightarrow
K_1(C^*(Y^{(j)}\subset X))
\cong K_1(C(S^1))
\]
therefore recovers exactly the \(j\)-th Toeplitz class. Applying
Corollary~\ref{M-cor:primitive_periodic_toeplitz} gives the displayed component winding. Reversing
the normal direction changes the sign of the Toeplitz extension for a fixed tangential coordinate;
the induced boundary orientation reverses the tangential coordinate at the same time, which
multiplies the winding homomorphism by the opposite sign and gives the final statement.
\end{proof}

\subsection{Proof of Result~\ref{thm:finite_corner_toeplitz}}\label{app:proof:thm:finite_corner_toeplitz}
\begin{proof}[Proof of \Cref{thm:finite_corner_toeplitz}]
Naturality of the six-term exact sequence for the quotient map of extensions
\cite{BlackadarKTheory} gives
\[
q_*\partial_{\mathrm{poly}}([P])
=
\partial_{\oplus\mathrm T}([P]),
\]
where \(\partial_{\oplus\mathrm T}\) is the connecting map for the quotient extension obtained
after deleting the finite corner ideal. By hypothesis this quotient extension is the direct sum of
the \(r\) primitive straight Toeplitz extensions, so
\[
\partial_{\oplus\mathrm T}([P])
=
\bigoplus_{j=1}^r\partial_{\mathrm T,t_j,n_j}([P]).
\]

Since \(\mathcal F\) is finite-dimensional, \(K_1(\mathcal F)=0\). Exactness of the six-term
sequence for
\[
0\to\mathcal F\to\mathcal I_{\mathrm{poly}}\xrightarrow{q}\bigoplus_j\mathcal I_j\to0
\]
therefore makes \(q_*\) injective on \(K_1(\mathcal I_{\mathrm{poly}})\). Applying the \(j\)-th
coordinate projection and then the standard winding homomorphism gives
Corollary~\ref{M-cor:primitive_periodic_toeplitz}, hence the displayed formula.
\end{proof}

\subsection{Proof of Result~\ref{cor:k1_trivial_corner_toeplitz}}\label{app:proof:cor:k1_trivial_corner_toeplitz}
\begin{proof}[Proof of \Cref{cor:k1_trivial_corner_toeplitz}]
The preceding proof used finite-dimensionality only to obtain \(K_1(\mathcal F)=0\). With that
condition assumed directly, exactness of
\[
0\to\mathcal F\to\mathcal I_{\mathrm{poly}}
\xrightarrow{q}\bigoplus_j\mathcal I_j\to0
\]
again makes
\[
q_*:K_1(\mathcal I_{\mathrm{poly}})\to K_1\Bigl(\bigoplus_j\mathcal I_j\Bigr)
\]
injective, and naturality identifies the quotient boundary class with the direct sum of the
straight Toeplitz classes. Finite-dimensional algebras and \(c_0\)-sums of them have vanishing
\(K_1\). For a bounded
product of finite-dimensional algebras, every unitary has a bounded self-adjoint logarithm with
norm at most \(\pi\) in each finite-dimensional factor, so the unitary group is path connected
after stabilization and \(K_1\) again vanishes.
\end{proof}

\subsection{Proof of Result~\ref{cor:finite_corner_geometric_criterion}}\label{app:proof:cor:finite_corner_geometric_criterion}
\begin{proof}[Proof of \Cref{cor:finite_corner_geometric_criterion}]
For an algebraic finite-propagation operator \(T\) supported near \(Y\), condition (ii) implies that
all matrix coefficients coupling two different edge pieces are supported inside a bounded
neighborhood of \(C\). Modulo the corner ideal, \(T\) is therefore the direct sum of its edge
restrictions. Condition (iii) identifies these restrictions with elements of the Toeplitz edge
ideals and gives the quotient map \(q\). Extension by zero from the separated edge collars gives
surjectivity of \(q\), again modulo corner-supported terms. If all edge restrictions vanish, then
the support of \(T\) is contained in a bounded neighborhood of \(C\), so the kernel is exactly
\(\mathcal F_C\) after taking norm closures. The final statement is
Corollary~\ref{cor:k1_trivial_corner_toeplitz}.
\end{proof}

\subsection{Proof of Result~\ref{M-cor:periodic_straight_edge_winding}}\label{app:proof:cor:periodic_straight_edge_winding}
\begin{proof}[Proof of \Cref{M-cor:periodic_straight_edge_winding}]
Corollary~\ref{cor:verified_periodic_families} identifies the doubled-and-jammed straight-interface
class with the boundary image of the periodic bulk Fermi projection. For a straight periodic
component, this boundary image is exactly the Toeplitz boundary map in
Theorem~\ref{M-thm:periodic_toeplitz_pairing}; for a primitive direction it is the lattice-change
version in Corollary~\ref{M-cor:primitive_periodic_toeplitz}. The Chern signs of the unperturbed
triangular model, the positive-duality-strength models, and the allowed small periodic
perturbations are supplied by Theorem~\ref{M-thm:periodic_triangular_chern} and
Corollary~\ref{M-cor:periodic_chern_stability}; finite-cover and finite-index cell variants are
Corollary~\ref{cor:periodic_covers_sublattices}.
Substituting those signs into Theorem~\ref{M-thm:periodic_toeplitz_pairing} gives the displayed
winding values.

The invariance under scalar jams, componentwise scalar jams, positive copy-side confinement,
finite initial-component changes, and uniformly gapped sidewise homotopies is exactly the
robustness statement in Corollary~\ref{cor:verified_periodic_families}; norm-vanishing component
defects are Proposition~\ref{prop:norm_vanishing_component_defects}. The split componentwise
statement is Theorem~\ref{thm:split_component_toeplitz}; the finite-corner version is
Theorem~\ref{thm:finite_corner_toeplitz}.
\end{proof}

\subsection{Proof of Result~\ref{cor:straight_toeplitz_robustness}}\label{app:proof:cor:straight_toeplitz_robustness}\label{app:proof:cor:finite_defect_tail_toeplitz}
\begin{proof}[Proof of \Cref{cor:straight_toeplitz_robustness}]
Parts (a) and (b) preserve the closed-bulk Chern signs by
Corollary~\ref{M-cor:periodic_chern_stability}. A fixed-collar perturbation in part (c) belongs to
the straight-edge interface ideal, so Proposition~\ref{prop:interface_supported_invariance} and
Proposition~\ref{prop:sidewise_data_interface_invariance} leave the interface class unchanged
whenever the sidewise comparison gaps remain open. Part (d) is
Theorem~\ref{thm:general_copy_interface} together with
Corollary~\ref{cor:copy_potential_independence}, and part (e) is
Propositions~\ref{prop:finite_component_stability},
\ref{prop:norm_vanishing_component_defects}, and~\ref{prop:sidewise_homotopy}. For part (f),
Proposition~\ref{prop:interface_ideal_coarse_invariance} identifies the two fixed-collar ideals.
For part (g), a finite tangential defect inside a fixed edge collar is finite rank, and every
bounded finite-propagation perturbation supported in a fixed collar belongs to the algebraic
interface ideal; the invariance follows from
Propositions~\ref{prop:interface_supported_invariance} and
\ref{prop:sidewise_data_interface_invariance}. Part (h) is
Proposition~\ref{prop:finite_component_stability} on the deleted components, followed by
Theorem~\ref{thm:split_component_toeplitz}, Theorem~\ref{thm:finite_corner_toeplitz}, or
Corollary~\ref{cor:k1_trivial_corner_toeplitz} on the tail. The
Toeplitz winding homomorphism, and each component coordinate winding in the split or finite-corner
split case, depends
only on the resulting \(K_1\)-class, so the winding values remain those of
Corollary~\ref{M-cor:periodic_straight_edge_winding}.
\end{proof}

\subsection{Proof of Result~\ref{thm:rough_eventual_toeplitz_equivalence}}\label{app:proof:thm:rough_eventual_toeplitz_equivalence}
\begin{proof}[Proof of \Cref{thm:rough_eventual_toeplitz_equivalence}]
The equality of ideals is Proposition~\ref{prop:interface_ideal_coarse_invariance}. Since the
sidewise differences belong to this common ideal, the quotient images of \(J_i\) and \(T_i^\pm\)
are the same. The proof of Proposition~\ref{prop:sidewise_data_interface_invariance} then gives the
displayed equality of quotient Fermi projections and of \(p_{J_i}\); applying the connecting map
gives equality of interface classes. The class for the straight tail is evaluated by
Corollary~\ref{M-cor:primitive_periodic_toeplitz}. A perturbation supported in a finite tangential
interval and a fixed normal collar acts on only finitely many lattice degrees of freedom, hence is
finite rank and in particular belongs to \(\mathcal I\). An eventually periodic collar differs from
its periodic tail exactly by such a transition defect.
\end{proof}

\subsection{Proof of Result~\ref{thm:tail_split_component_windings}}\label{app:proof:thm:tail_split_component_windings}
\begin{proof}[Proof of \Cref{thm:tail_split_component_windings}]
Deleting or changing finitely many initial components changes the sidewise data by an operator in
the interface ideal by Proposition~\ref{prop:finite_component_stability}; hence the full interface
class and the tail class have the same image under the quotient that discards those components. On
the tail, the split case is Theorem~\ref{thm:split_component_toeplitz}. The finite-corner case is
Theorem~\ref{thm:finite_corner_toeplitz}, and the \(K_1\)-trivial version is
Corollary~\ref{cor:k1_trivial_corner_toeplitz}. In all cases, applying the \(j\)-th coordinate
projection and the winding homomorphism gives
\(\det[t_j\ n_j]\Ch(P)\). The finite-cover formula is
Corollary~\ref{cor:covered_primitive_toeplitz}.
\end{proof}

\subsection[Proof of Result~\ref{cor:fully_verified_straight_periodic}]{Proof of Result~\ref{cor:fully_verified_straight_periodic}}\label{app:proof:cor:fully_verified_straight_periodic}
\begin{proof}[Proof of \Cref{cor:fully_verified_straight_periodic}]
The zero gap and Chern signs for (a)--(c) are Theorem~\ref{M-thm:periodic_triangular_bulk},
Theorem~\ref{M-thm:periodic_triangular_chern}, and
Corollary~\ref{M-cor:periodic_chern_stability}; finite covers and cell enlargements are
Corollary~\ref{cor:periodic_covers_sublattices}. The straight-edge winding is
Theorem~\ref{M-thm:periodic_toeplitz_pairing} in coordinate directions and
Corollary~\ref{M-cor:primitive_periodic_toeplitz} in arbitrary primitive directions. Finite split
sums are Theorem~\ref{thm:split_component_toeplitz}; finite-corner split sums are
Theorem~\ref{thm:finite_corner_toeplitz} and
Corollary~\ref{cor:k1_trivial_corner_toeplitz}; the collar-separation criterion is
Corollary~\ref{cor:finite_corner_geometric_criterion}. Bounded-distance and eventually periodic
collars are Theorem~\ref{thm:rough_eventual_toeplitz_equivalence}; finite-defect and tail-split
variants are Corollary~\ref{cor:finite_defect_tail_toeplitz} and
Theorem~\ref{thm:tail_split_component_windings}. The scalar-jam, componentwise-jam, and positive
copy-confinement invariance follows from Corollary~\ref{M-cor:periodic_straight_edge_winding} and
Corollary~\ref{cor:straight_toeplitz_robustness}. Norm-vanishing component defects are
Proposition~\ref{prop:norm_vanishing_component_defects}. The final sentence is the fixed-collar and
homotopy part of Corollary~\ref{cor:straight_toeplitz_robustness}.
\end{proof}

\subsection{Proof of Result~\ref{lem:hex_massless_closure}}\label{app:proof:lem:hex_massless_closure}
\begin{proof}[Proof of \Cref{lem:hex_massless_closure}]
At \(k=(0,0)\), \(a_1=a_2=0\) and \(a_3^{\hexagon}=1\). Thus
\[
H_{\hexagon}(0,0;0)=
\begin{pmatrix}
0&0\\
0&0
\end{pmatrix},
\]
so \(\det H_{\hexagon}(0,0;0)=0\) and zero belongs to the spectrum.
\end{proof}

\subsection{Proof of Result~\ref{lem:hex_bulk_gap}}\label{app:proof:lem:hex_bulk_gap}
\begin{proof}[Proof of \Cref{lem:hex_bulk_gap}]
The trace of \(H_{\hexagon}(k;m)\) is zero, and a direct \(2\times2\) determinant calculation gives
\[
\det H_{\hexagon}(k;m)
=-\sin^2 k_1-\sin^2 k_2-m^2(\cos k_1+\cos k_2-1)^2 .
\]
Set \(s=\sin^2 k_1+\sin^2 k_2\). If \(s\ge1/4\), then
\(|\det H_{\hexagon}(k;m)|\ge1/4\). If \(s<1/4\), then
\(|\sin k_j|<1/2\), hence \(|\cos k_j|>\sqrt3/2\), for \(j=1,2\). Checking the four sign choices
for \((\cos k_1,\cos k_2)\) gives
\[
|\cos k_1+\cos k_2-1|\ge\sqrt{3}-1 .
\]
Therefore
\[
|\det H_{\hexagon}(k;m)|
\ge
\min\left\{\frac14,\,(\sqrt{3}-1)^2m^2\right\}.
\]
Since the eigenvalues of a trace-zero \(2\times2\) Hermitian matrix are
\(\pm\sqrt{|\det H_{\hexagon}(k;m)|}\), the stated spectral gap follows.
\end{proof}

\subsection{Proof of Result~\ref{lem:hex_chern_positive}}\label{app:proof:lem:hex_chern_positive}
\begin{proof}[Proof of \Cref{lem:hex_chern_positive}]
For \(m>0\), the path replacing \(m\) by \((1-s)m+s\), \(0\le s\le1\), stays in the positive-mass
region and is gapped by \Cref{lem:hex_bulk_gap}. Hence the Chern number equals the value at
\(m=1\). Write
\[
d_{\hexagon}(k)=(\sin k_1,\sin k_2,\cos k_1+\cos k_2-1),\qquad
\widehat d_{\hexagon}=d_{\hexagon}/|d_{\hexagon}|.
\]
For a two-band Hamiltonian \(d\cdot\sigma\), the negative-energy Chern number is
\[
\Ch(E_-)=
-\frac1{4\pi}\int_{\mathbb T^2}
\widehat d\cdot(\partial_{k_1}\widehat d\times\partial_{k_2}\widehat d)\,dk_1\,dk_2
=-\deg(\widehat d),
\]
with the same convention used in the Bloch-bundle formula; see
\cite[Chapters~4--5]{AsbothOroszlanyPalyi2016} and the Bellissard--van Elst--Schulz-Baldes
formulation of the Chern character \cite{BellissardVanElstSchulzBaldes1994}.

The north pole \((0,0,1)\) is a regular value. Its preimages satisfy
\(\sin k_1=\sin k_2=0\) and \(\cos k_1+\cos k_2-1>0\). The only such point is \((0,0)\). The local
degree there is
\[
\operatorname{sgn}\det
\begin{pmatrix}
\cos k_1&0\\
0&\cos k_2
\end{pmatrix}_{(0,0)}
=+1 .
\]
Thus \(\deg(\widehat d_{\hexagon})=1\), and the negative-energy Chern number is \(-1\).
\end{proof}

\subsection{Proof of Result~\ref{lem:hex_chern_negative}}\label{app:proof:lem:hex_chern_negative}
\begin{proof}[Proof of \Cref{lem:hex_chern_negative}]
Let \(U=\begin{pmatrix}0&-i\\ i&0\end{pmatrix}\). Since \(a_1\) and \(a_2\) are odd under
\(k\mapsto-k\), while \(a_3^{\hexagon}\) is even,
\[
H_{\hexagon}(k;-m)=U\,\overline{H_{\hexagon}(-k;m)}\,U^* .
\]
The base involution \(k\mapsto-k\) preserves the orientation \(dk_1\wedge dk_2\). Unitary
conjugation preserves the Bloch line bundle, while complex conjugation changes the sign of its first
Chern class. Therefore the negative-energy Chern number at \(-m\) is the negative of the value at
\(m>0\). By \Cref{lem:hex_chern_positive}, it is \(+1\).
\end{proof}

\subsection{Proof of Result~\ref{lem:hex_toeplitz_winding}}\label{app:proof:lem:hex_toeplitz_winding}
\begin{proof}[Proof of \Cref{lem:hex_toeplitz_winding}]
The entries of \(H_{\hexagon}(k;m)\) are Laurent polynomials in \(e^{ik_1}\) and \(e^{ik_2}\), so the
symbol is finite range. By \Cref{lem:hex_bulk_gap}, it is gapped at zero for \(m\ne0\). The straight
half-space Toeplitz extension therefore applies. The coordinate statement is
\Cref{M-thm:periodic_toeplitz_pairing}, and the primitive-basis statement is
\Cref{M-cor:primitive_periodic_toeplitz}. Substituting
\Cref{lem:hex_chern_positive,lem:hex_chern_negative} gives the displayed signs.
\end{proof}

\subsection{Proof of Result~\ref{thm:hex_bulk_edge}}\label{app:proof:thm:hex_bulk_edge}
\begin{proof}[Proof of \Cref{thm:hex_bulk_edge}]
The gap estimate is \Cref{lem:hex_bulk_gap}. The Chern value is
\Cref{lem:hex_chern_positive} for \(m>0\) and \Cref{lem:hex_chern_negative} for \(m<0\), which is the
formula \(-\operatorname{sgn}(m)\). The Toeplitz formula is
\Cref{lem:hex_toeplitz_winding}.
\end{proof}

\subsection{Proof of Result~\ref{lem:square_massless_closure}}\label{app:proof:lem:square_massless_closure}
\begin{proof}[Proof of \Cref{lem:square_massless_closure}]
At \(k=(0,0)\), \(b_1=b_2=0\) and \(b_3^{\square}=3\). Hence
\[
H_{\square}(0,0;0)=
\begin{pmatrix}
0&0\\
0&0
\end{pmatrix},
\]
so \(\det H_{\square}(0,0;0)=0\).
\end{proof}

\subsection{Proof of Result~\ref{lem:square_bulk_gap}}\label{app:proof:lem:square_bulk_gap}
\begin{proof}[Proof of \Cref{lem:square_bulk_gap}]
Again the matrix is trace zero, and
\[
\det H_{\square}(k;m)
=-\sin^2 k_1-\sin^2 k_2-m^2(1+\cos k_1+\cos k_2)^2 .
\]
If \(s=\sin^2k_1+\sin^2k_2\ge1/4\), then
\(|\det H_{\square}(k;m)|\ge1/4\). If \(s<1/4\), then
\(|\cos k_j|>\sqrt3/2\), \(j=1,2\), and the four possible sign choices give
\[
|1+\cos k_1+\cos k_2|\ge\sqrt{3}-1 .
\]
This gives the lower bound by \(c_{\square}(m)\). The eigenvalues are
\(\pm\sqrt{|\det H_{\square}(k;m)|}\), so the distance from zero to the spectrum is at least
\(\sqrt{c_{\square}(m)}\).
\end{proof}

\subsection{Proof of Result~\ref{lem:square_chern_positive}}\label{app:proof:lem:square_chern_positive}
\begin{proof}[Proof of \Cref{lem:square_chern_positive}]
For \(m>0\), the positive-mass path to \(m=1\) stays gapped by \Cref{lem:square_bulk_gap}. At
\(m=1\), set
\[
d_{\square}(k)=(\sin k_1,\sin k_2,1+\cos k_1+\cos k_2),\qquad
\widehat d_{\square}=d_{\square}/|d_{\square}|.
\]
As in the two-band Chern formula cited in the proof of \Cref{lem:hex_chern_positive},
\(\Ch(E_-)=-\deg(\widehat d_{\square})\). The north-pole preimages are
\[
(0,0),\qquad(\pi,0),\qquad(0,\pi),
\]
since \(1+\cos k_1+\cos k_2>0\) at exactly those three points among the four points with
\(\sin k_1=\sin k_2=0\). The local signs are respectively
\[
\operatorname{sgn}(\cos k_1\cos k_2)=+1,\ -1,\ -1 .
\]
Thus \(\deg(\widehat d_{\square})=-1\), and the negative-energy Chern number is \(+1\).
\end{proof}

\subsection{Proof of Result~\ref{lem:square_toeplitz_winding}}\label{app:proof:lem:square_toeplitz_winding}
\begin{proof}[Proof of \Cref{lem:square_toeplitz_winding}]
The entries of \(H_{\square}(k;m)\) are Laurent polynomials, and \Cref{lem:square_bulk_gap} gives a
zero gap for \(m\ne0\). The straight-edge Toeplitz index formula in
\Cref{M-thm:periodic_toeplitz_pairing} and the primitive-coordinate version
\Cref{M-cor:primitive_periodic_toeplitz} give
\[
\operatorname{wind}\!\left(\partial_{\mathrm T,t,n}([P_{\square,m}])\right)
=
\det[t\ n]\Ch(P_{\square,m}).
\]
For \(m>0\), \Cref{lem:square_chern_positive} gives \(+1\). For \(m<0\), the same antiunitary
identity used in \Cref{lem:hex_chern_negative}, with \(b_3^{\square}(-k)=b_3^{\square}(k)\), changes
the Chern sign. This gives \(-1\).
\end{proof}

\subsection{Proof of Result~\ref{thm:square_bulk_edge}}\label{app:proof:thm:square_bulk_edge}
\begin{proof}[Proof of \Cref{thm:square_bulk_edge}]
\Cref{lem:square_bulk_gap} gives the zero gap. \Cref{lem:square_chern_positive} gives the Chern
number \(+1\) for \(m>0\), and the antiunitary mass-reversal argument in the proof of
\Cref{lem:square_toeplitz_winding} gives \(-1\) for \(m<0\). The Toeplitz winding formula is
\Cref{lem:square_toeplitz_winding}.
\end{proof}

\subsection{Proof of Result~\ref{lem:roe_ideal}}\label{app:proof:lem:roe_ideal}
\begin{proof}[Proof of \Cref{lem:roe_ideal}]
At the algebraic level, operators supported within a bounded neighborhood of $Y^{(R)}$ are stable
under addition and adjoint. They are also stable under multiplication on either side by bounded
finite-propagation operators in $C^*(X)$, because finite propagation enlarges the support
neighborhood by at most a bounded amount. Taking norm closure gives a closed two-sided ideal.
\end{proof}

\subsection{Proof of Result~\ref{prop:collar_radius_independence}}\label{app:proof:prop:collar_radius_independence}
\begin{proof}[Proof of \Cref{prop:collar_radius_independence}]
Assume $R'\ge R$. Then
\[
Y^{(R)}\subset Y^{(R')}\subset N_{R'}(Y^{(R)}).
\]
Thus an operator supported within a bounded neighborhood of $Y^{(R')}$ is also supported within a
bounded neighborhood of $Y^{(R)}$, and the converse inclusion is immediate from
$Y^{(R)}\subset Y^{(R')}$. The corresponding algebraic ideals therefore coincide, and taking norm
closures gives the equality. The case $R>R'$ is symmetric.
\end{proof}

\subsection{Proof of Result~\ref{prop:interface_ideal_coarse_invariance}}\label{app:proof:prop:interface_ideal_coarse_invariance}
\begin{proof}[Proof of \Cref{prop:interface_ideal_coarse_invariance}]
Let \(T\) be an algebraic Roe operator whose matrix support is contained in
\(N_s(Y_0)\times N_s(Y_0)\) for some \(s<\infty\). Since \(Y_0\subset N_R(Y_1)\), the same matrix
support is contained in \(N_{s+R}(Y_1)\times N_{s+R}(Y_1)\). Hence the algebraic interface ideal
for \(Y_0\) is contained in the algebraic interface ideal for \(Y_1\). The reverse inclusion is
identical, using \(Y_1\subset N_R(Y_0)\). Taking norm closures gives the equality.
\end{proof}

\subsection{Proof of Result~\ref{lem:S_uniform_norm}}\label{app:proof:lem:S_uniform_norm}
\begin{proof}[Proof of \Cref{lem:S_uniform_norm}]
The propagation bound is exactly Assumption~\ref{ass:S_uniform}(i). For the norm bound, every row
sum and every column sum of absolute values is at most $N_Sc_S$ by
Assumption~\ref{ass:S_uniform}(ii)--(iii). The Schur test therefore gives
\[
\opnorm{S_n}\le N_Sc_S.
\]
\end{proof}

\subsection{Proof of Result~\ref{prop:uniform_locality_norm}}\label{app:proof:prop:uniform_locality_norm}
\begin{proof}[Proof of \Cref{prop:uniform_locality_norm}]
By the Higson--Prodan construction, $B+B^\dagger$ maps a simplex barycenter to barycenters of
simplices that share an incidence relation. In the barycentric graph metric, every such pair is at
distance $1$. Hence $B+B^\dagger$ has propagation exactly $1$ on every $\ell^2(X_n)$. The operator
$S$ is handled by Lemma~\ref{lem:S_uniform_norm}. Since
$H_{\pm,n}=(B+B^\dagger)\pm S$, the propagation satisfies
\[
\mathrm{prop}(H_{\pm,n})\le \max\{1,R_S\},
\]
which is bounded uniformly in $n$. This proves~(a).

For (b), $\opnorm{B+B^\dagger}$ is bounded by twice the maximum simplex degree in the
barycentric graph. Under uniform regularity (Definition~\ref{def:uniform_regular}), the ball sizes
$|B_{X_n}(x,1)|$ are uniformly bounded, so the maximum degree is uniformly bounded. Similarly,
$\opnorm{S_n}$ is bounded uniformly in $n$ by Lemma~\ref{lem:S_uniform_norm}. Hence
\[
\opnorm{H_{\pm,n}}\le \opnorm{B+B^\dagger}+\opnorm{S_n}\le B_{\mathrm{bulk}}
\]
for a constant $B_{\mathrm{bulk}}$ independent of $n$.
\end{proof}

\subsection{Proof of Result~\ref{thm:hp_closed_input}}\label{app:proof:thm:hp_closed_input}
\begin{proof}[Proof of \Cref{thm:hp_closed_input}]
Assumption~\ref{ass:hp_admissibility} includes uniform regularity of $(\K_n)_n$, and its part~(a)
gives Assumption~\ref{ass:S_uniform}. Its part~(b) gives the uniform norm bound and the uniform
spectral gap in Assumption~\ref{ass:uniform_bulk}; the propagation bound follows from
Proposition~\ref{prop:uniform_locality_norm}. The identification of \(P_\pm\) with the
Higson--Prodan Fermi projection and the Chern computation are precisely part~(c) of the imported
closed-bulk input.
\end{proof}

\subsection{Proof of Result~\ref{prop:bulk_class}}\label{app:proof:prop:bulk_class}
\begin{proof}[Proof of \Cref{prop:bulk_class}]
Boundedness and self-adjointness are immediate from the direct-sum construction and
Assumption~\ref{ass:uniform_bulk}(b). The propagation bound in
Assumption~\ref{ass:uniform_bulk}(a) passes to the direct sum, so $H_\pm^\infty$ has propagation
at most $R_{\mathrm{loc}}$ on $\ell^2(X)$.

It remains to verify local compactness. Let $F\subset X$ be finite and write $\chi_F$ for
multiplication by its characteristic function. Because $F$ is finite and $H_\pm^\infty$ has
propagation at most $R_{\mathrm{loc}}$, every vector in the range of $\chi_F H_\pm^\infty$ is
supported inside the finite set $N_{R_{\mathrm{loc}}}(F)$. Hence $\chi_F H_\pm^\infty$ has
finite-dimensional range and is therefore finite rank. The same argument applied to the adjoint
shows that $H_\pm^\infty\chi_F$ is finite rank as well. Thus $H_\pm^\infty$ is locally compact and
belongs to the algebraic Roe algebra; therefore $H_\pm^\infty\in C^*(X)$ after taking norm closure.

By Assumption~\ref{ass:uniform_bulk}(b),
\[
\spec_{\mathrm{op}}(H_\pm^\infty)\subset [-B_{\mathrm{bulk}},B_{\mathrm{bulk}}],
\]
and by Assumption~\ref{ass:uniform_bulk}(c),
\[
\spec_{\mathrm{op}}(H_\pm^\infty)\subset (-\infty,-g_{\mathrm{bulk}}]\cup[g_{\mathrm{bulk}},\infty).
\]
Choose a compactly supported continuous function $f\in C_c(\mathbb R)\subset C_0(\mathbb R)$ such
that
\[
f(t)=1\quad\text{for }t\in[-B_{\mathrm{bulk}},-g_{\mathrm{bulk}}/2],
\qquad
f(t)=0\quad\text{for }t\in[g_{\mathrm{bulk}}/2,B_{\mathrm{bulk}}],
\]
and, for example, $f(t)=0$ for $|t|\ge B_{\mathrm{bulk}}+1$. Then $f$ agrees with
$\chi_{(-\infty,0)}$ on $\spec_{\mathrm{op}}(H_\pm^\infty)$ (which avoids the interval
$(-g_{\mathrm{bulk}},g_{\mathrm{bulk}})$), so the operator functional calculus on
$\ell^2(X)$ gives
\[
f(H_\pm^\infty)=\chi_{(-\infty,0)}(H_\pm^\infty)=P_\pm.
\]
Since $H_\pm^\infty\in C^*(X)$ and $f\in C_0(\mathbb R)$, continuous functional calculus in
$C^*(X)$ gives $f(H_\pm^\infty)\in C^*(X)$.
Hence $P_\pm\in C^*(X)$, and so $P_\pm$ defines a class
\[
[P_\pm]\in K_0(C^*(X)).
\]
\end{proof}

\subsection{Proof of Result~\ref{prop:bulk_gap_homotopy}}\label{app:proof:prop:bulk_gap_homotopy}
\begin{proof}[Proof of \Cref{prop:bulk_gap_homotopy}]
Choose one function $f\in C_c(\mathbb R)$ such that
\[
f(t)=1\quad\text{for }t\in[-B,-g/2],
\qquad
f(t)=0\quad\text{for }t\in[g/2,B].
\]
Then $P_s=f(H_s)$ for every $s$. Continuous functional calculus is norm-continuous for a fixed
continuous function on a norm-continuous path of self-adjoint elements, so $s\mapsto P_s$ is a
norm-continuous path of projections in $C^*(X)$. A norm-continuous path of projections represents a
constant $K_0$-class.
\end{proof}

\subsection{Proof of Result~\ref{prop:quantitative_projection_stability}}\label{app:proof:prop:quantitative_projection_stability}
\begin{proof}[Proof of \Cref{prop:quantitative_projection_stability}]
For $z\in\Gamma$,
\[
\opnorm{(H-z)^{-1}}\le \delta_\Gamma^{-1}.
\]
Since
\[
\opnorm{V(H-z)^{-1}}\le \frac{\opnorm V}{\delta_\Gamma}<1,
\]
the operator $H+V-z$ is invertible and
\[
\opnorm{(H+V-z)^{-1}}\le \frac1{\delta_\Gamma-\opnorm V}.
\]
The resolvent identity gives
\[
(H+V-z)^{-1}-(H-z)^{-1}=-(H+V-z)^{-1}V(H-z)^{-1}.
\]
Integrating this identity around $\Gamma$ gives the stated projection estimate.
\end{proof}

\subsection{Proof of Result~\ref{cor:bulk_small_perturbation}}\label{app:proof:cor:bulk_small_perturbation}\label{app:proof:cor:positive_interface_small_perturbation}
\begin{proof}[Proof of \Cref{cor:bulk_small_perturbation}]
For $t\in[0,1]$, set $H_t:=H+tV$. Weyl's spectral-stability estimate gives
\[
\spec_{\mathrm{op}}(H_t)\cap\bigl(-(g-t\opnorm V),\,g-t\opnorm V\bigr)=\emptyset.
\]
Since $g-\opnorm V>0$, the path $H_t$ is uniformly gapped at zero. Proposition~\ref{prop:bulk_gap_homotopy}
therefore applies, with spectral bound $B+\opnorm V$ and gap $g-\opnorm V$, and gives equality of
the endpoint Fermi-projection classes. In the interface case, the path
\[
s\longmapsto \chi_{(-\infty,0)}(H_\pm^\infty+sV)
\]
is a norm-continuous path of projections joining $P_\pm$ to $P_V$. Applying
Proposition~\ref{prop:positive_boundary_homotopy} gives the stated equality.
\end{proof}

\subsection{Proof of Result~\ref{prop:collar_locality_finite}}\label{app:proof:prop:collar_locality_finite}
\begin{proof}[Proof of \Cref{prop:collar_locality_finite}]
The difference $H_{\pm,n}-H_{\pm,n}^{\mathrm{split}}$ is exactly the sum of the two off-diagonal
blocks
\[
P_{\mathrm{orig},n}H_{\pm,n}P_{\mathrm{copy},n}
+
P_{\mathrm{copy},n}H_{\pm,n}P_{\mathrm{orig},n}.
\]
Thus it suffices to prove that every nonzero matrix element between the original side and the strict
copy side is supported inside $Y_n$.

A vertex of the barycentric graph is the barycenter of a simplex of $D(\K_n)$, and an edge joins
two vertices precisely when the corresponding simplices are incident by inclusion. Let
\[
U_{n,\mathrm{orig}}:=X_{n,\mathrm{orig}}\setminus \Gamma_n,
\qquad
U_{n,\mathrm{copy}}:=X_{n,\mathrm{copy}}\setminus \Gamma_n.
\]
We claim that no graph edge joins a vertex of $U_{n,\mathrm{orig}}$ directly to a vertex of
$U_{n,\mathrm{copy}}$. An edge in the barycentric graph corresponds to a pair of simplices
$\sigma\subset\tau$ or $\tau\subset\sigma$. If one of these simplices lies strictly on the original
side and the other lies strictly on the copy side, then they cannot be nested: a simplex containing
vertices from both sides necessarily meets the glued boundary, and a boundary simplex itself already
belongs to $\Gamma_n$. Hence every adjacency from the original side to the copy side must pass
through a simplex meeting the glued boundary.

As a consequence, every graph path from $X_{n,\mathrm{orig}}$ to $X_{n,\mathrm{copy}}$ must pass through
$\Gamma_n$. Now let
\[
x\in X_{n,\mathrm{orig}},
\qquad
y\in X_{n,\mathrm{copy}},
\qquad
\langle x,H_{\pm,n}y\rangle\neq0.
\]
Since $H_{\pm,n}$ has propagation at most $R_{\mathrm{loc}}$, we have
\[
d_{X_n}(x,y)\le R_{\mathrm{loc}}.
\]
Choose a geodesic path from $x$ to $y$ in the barycentric graph. By the previous paragraph, this
path meets $\Gamma_n$ at some vertex $z$. Because the path has total length at most
$R_{\mathrm{loc}}$, we have
\[
d_{X_n}(x,z)\le R_{\mathrm{loc}},
\qquad
d_{X_n}(y,z)\le R_{\mathrm{loc}}.
\]
Hence both $x$ and $y$ belong to $N_{R_{\mathrm{loc}}}(\Gamma_n)=Y_n$. Therefore every
original--copy matrix coefficient of $H_{\pm,n}$ vanishes after compressing to $X_n\setminus Y_n$,
which is exactly \eqref{eq:finite_collar}.

For \eqref{eq:finite_collar_jam}, observe that the jam term is already diagonal with respect to the
decomposition $\Hil_{\mathrm{orig},n}\oplus\Hil_{\mathrm{copy},n}$. Consequently,
\[
H_{\pm,M,n}-\bigl(H_{\pm,n}^{\mathrm{split}}+MP_{\mathrm{copy},n}\bigr)
=
H_{\pm,n}-H_{\pm,n}^{\mathrm{split}},
\]
and \eqref{eq:finite_collar_jam} follows from \eqref{eq:finite_collar}.
\end{proof}

\subsection{Proof of Result~\ref{cor:collar_locality_coarse}}\label{app:proof:cor:collar_locality_coarse}
\begin{proof}[Proof of \Cref{cor:collar_locality_coarse}]
By Proposition~\ref{prop:collar_locality_finite}, each summand of either difference is supported
inside $Y_n\times Y_n$. In both differences the jam term is absent or cancels exactly, so their
propagations are uniformly bounded by $R_{\mathrm{loc}}$ and their norms are uniformly bounded by
the uniform bound on $(H_{\pm,n})_n$. Hence the direct sums are bounded operators on $\ell^2(X)$
with uniformly finite propagation and support contained in a bounded neighborhood of $Y$.

Because $X$ has bounded geometry and we work with scalar coefficients, such operators are locally
compact: if $F\subset X$ is finite, then finite propagation implies that both $T\chi_F$ and
$\chi_FT$ have range contained in $\ell^2$ of a finite neighborhood of $F$, hence are finite rank.
Therefore the two direct sums belong to the algebraic interface ideal, and thus to its norm closure
$C^*(Y\subset X)$.
\end{proof}

\subsection{Proof of Result~\ref{prop:trivial_jam}}\label{app:proof:prop:trivial_jam}
\begin{proof}[Proof of \Cref{prop:trivial_jam}]
The operator $\widetilde T^-_{\pm,M}$ is a direct sum of bounded local self-adjoint operators, so it
is bounded, self-adjoint, and local. On the original-side sector it equals the identity, hence has
spectral value $1$.

On the copy-side sector it equals $H_{\pm,\mathrm{copy}}^\infty+MP_{\mathrm{copy}}^\infty$. For each $n$
every spectral value $\lambda$ of $H_{\pm,n}^{\mathrm{copy}}+M\Id$ satisfies
\[
\lambda\ge M+\lambda_{\min}(H_{\pm,n}^{\mathrm{copy}}).
\]
By the definition of $B_{\mathrm{copy},-}$, the right-hand side is strictly positive when
$M>B_{\mathrm{copy},-}$. Thus the whole operator spectrum of $\widetilde T^-_{\pm,M}$ is contained in
$\{1\}\cup [M-B_{\mathrm{copy},-},\infty)$, which gives the stated spectral-distance bound. In
particular, the spectrum is contained in $(0,\infty)$, so the negative operator spectral projection
vanishes.
\end{proof}

\subsection{Proof of Result~\ref{prop:coarse_excision}}\label{app:proof:prop:coarse_excision}
\begin{proof}[Proof of \Cref{prop:coarse_excision}]
Fix $r>0$. By the definition of coarse disjoint union, there exists $N_r$ such that
\[
\dist(X_n,X_m)>2r
\qquad\text{whenever }n\neq m\text{ and }\max\{n,m\}>N_r.
\]
Let
\[
x\in N_r(X^+)\cap N_r(X^-).
\]
Choose $x_+\in X^+$ and $x_-\in X^-$ with
\[
d(x,x_+)\le r,
\qquad
d(x,x_-)\le r.
\]
Then $d(x_+,x_-)\le2r$. Suppose first that either $x_+$ or $x_-$ lies in a component $X_n$ with
$n>N_r$. The separation condition forces both points to lie in that same component. If either
point lies in $Y_n$, then $x\in N_r(Y)$. Otherwise
\[
x_+\in X_{n,\mathrm{orig}}\setminus Y_n,
\qquad
x_-\in X_{n,\mathrm{copy}}\setminus Y_n.
\]
By the same separation argument used in Proposition~\ref{prop:collar_locality_finite}, every
geodesic in the barycentric graph from $x_+$ to $x_-$ meets $\Gamma_n\subset Y_n$. Therefore
$d(x_+,Y_n)\le 2r$, and consequently
\[
d(x,Y)\le d(x,x_+)+d(x_+,Y_n)\le 3r.
\]
Thus every such $x$ lies in $N_{3r}(Y)$.

It remains to treat the case in which neither $x_+$ nor $x_-$ lies in a component of index
larger than $N_r$. Let
\[
F_r:=\bigsqcup_{n=1}^{N_r} X_n.
\]
This is a finite set. Because each $\K_n$ has nonempty boundary, every $Y_n$ is nonempty, so
\[
c_r:=\max_{x\in N_r(F_r)} d(x,Y)
\]
is finite. Therefore
\[
N_r(X^+)\cap N_r(X^-)\subset N_{s(r)}(Y),
\qquad
s(r):=\max\{3r,c_r\},
\]
which is exactly coarse excision.
\end{proof}

\subsection{Proof of Result~\ref{prop:quotient_split}}\label{app:proof:prop:quotient_split}
\begin{proof}[Proof of \Cref{prop:quotient_split}]
For $Q_YT$ and $TQ_Y$, finite propagation implies
\[
Q_YT=\chi_Y\,T\,\chi_{N_R(Y)},
\qquad
TQ_Y=\chi_{N_R(Y)}\,T\,\chi_Y,
\]
so both are supported in $N_R(Y)\times N_R(Y)$ and hence belong to $\mathcal I$.

Now consider $Q_+TQ_-$. If
\[
\langle x,Q_+TQ_-\,y\rangle\neq0,
\]
then $x\in Z^+$, $y\in Z^-$, and $d(x,y)\le R$. Therefore
\[
x\in N_R(X^-),\qquad y\in N_R(X^+).
\]
Since also $x\in X^+$ and $y\in X^-$, Proposition~\ref{prop:coarse_excision} gives
\[
x,y\in N_{s(R)}(Y).
\]
Hence $Q_+TQ_-$ is supported in $N_{s(R)}(Y)\times N_{s(R)}(Y)$, so $Q_+TQ_-\in\mathcal I$.
The same argument gives $Q_-TQ_+\in\mathcal I$.

Because $\mathcal I$ is closed and algebraic Roe operators are dense in $\mathcal A$, the same
containments hold for every $T\in\mathcal A$.

We also record that multiplication by $Q_+$, $Q_-$, or $Q_Y$ preserves $\mathcal I$: algebraic
Roe operators supported near $Y$ remain supported near $Y$ after such left or right multiplication,
and the claim follows by norm closure.

We now prove part~(ii). Since $Q_Y\in\mathcal I$ by part~(i), with $T=\Id$, we have
\[
e_++e_-=\pi(Q_++Q_-)=\pi(\Id-Q_Y)=\Id_{\mathcal A/\mathcal I}.
\]
Orthogonality is immediate from $Q_+Q_-=0$. The centrality identity proved below shows that
$e_\pm$ are central projections.

We next prove part~(iii). If $\pi(a)=\pi(a')$, then $a-a'\in\mathcal I$, hence
\[
Q_\pm(a-a')\in\mathcal I.
\]
Therefore $E_\pm$ are well-defined.

Idempotence and orthogonality are immediate:
\[
E_\pm^2(\pi(a))=\pi(Q_\pm Q_\pm a)=E_\pm(\pi(a)),
\qquad
E_+E_-(\pi(a))=\pi(Q_+Q_-a)=0,
\]
and similarly for $E_-E_+$. Also,
\[
\pi(a)=\pi\bigl((Q_++Q_-+Q_Y)a\bigr)=E_+(\pi(a))+E_-(\pi(a)),
\]
because $Q_Ya\in\mathcal I$. Thus $E_+ + E_- = \Id_{\mathcal A/\mathcal I}$.

To prove multiplicativity for $E_+$, let $a,b\in\mathcal A$. Then
\[
E_+(\pi(a))E_+(\pi(b))
=
\pi(Q_+a)\,\pi(Q_+b)
=
\pi(Q_+aQ_+b).
\]
Now
\[
Q_+ab-Q_+aQ_+b
=
Q_+a(Q_-+Q_Y)b.
\]
The term $Q_+aQ_-b$ lies in $\mathcal I$ because $Q_+aQ_-\in\mathcal I$ and $\mathcal I$ is an
ideal. The term $Q_+aQ_Yb$ also lies in $\mathcal I$ because $aQ_Y\in\mathcal I$. Hence
\[
\pi(Q_+aQ_+b)=\pi(Q_+ab)=E_+(\pi(ab)).
\]
So $E_+$ is multiplicative. The same proof applies to $E_-$. The $*$-property is immediate:
\[
E_+(\pi(a))^*=\pi((Q_+a)^*)=\pi(a^*Q_+)=E_+(\pi(a^*)),
\]
because
\[
a^*Q_+-Q_+a^*
=
(Q_+a-aQ_+)^*
\in\mathcal I,
\]
as shown next.

To prove the centrality identity, let $a,b\in\mathcal A$. Then
\[
E_+(\pi(a))\,\pi(b)-\pi(a)\,E_+(\pi(b))
=
\pi\bigl((Q_+a-aQ_+)b\bigr).
\]
Now
\[
Q_+a-aQ_+
=
Q_+a(Q_-+Q_Y)-(Q_-+Q_Y)aQ_+.
\]
Each term on the right belongs to $\mathcal I$: the $Q_+aQ_-$ and $Q_-aQ_+$ terms are in
$\mathcal I$ by part~(i), and the terms involving $Q_Y$ are in $\mathcal I$ because $aQ_Y$ and
$Q_Ya$ are in $\mathcal I$ and $\mathcal I$ is stable under multiplication by $Q_+$. Hence
$Q_+a-aQ_+\in\mathcal I$, so
\[
E_+(\pi(a))\,\pi(b)=\pi(a)\,E_+(\pi(b)).
\]
Thus $E_+$ is central in the sense stated. The same holds for $E_-$.

Finally, every $b\in\mathcal A/\mathcal I$ decomposes as
\[
b=E_+(b)+E_-(b),
\]
and orthogonality $E_+E_-=E_-E_+=0$ implies uniqueness of this decomposition. Since
$E_\pm(\mathcal A/\mathcal I)$ are closed two-sided ideals, this proves the direct-sum decomposition.
\end{proof}

\subsection{Proof of Result~\ref{prop:quotient_sidewise}}\label{app:proof:prop:quotient_sidewise}
\begin{proof}[Proof of \Cref{prop:quotient_sidewise}]
The operators \(J,T^+,T^-\) are locally compact: for every finite \(F\subset X\), the operators
\(\chi_FT\) and \(T\chi_F\) have finite-dimensional range, with \(T\) equal to \(J\), \(T^+\), or
\(T^-\). Hence
\[
J,T^+,T^-\in\mathcal A=C^*(X),
\]
so $\pi(J)$ and $\pi(T^\pm)$ are well-defined.

Since $0$ lies in an operator spectral gap of $T^\pm$, choose compactly supported continuous
functions $h_\pm$ that agree with $\chi_{(-\infty,0)}$ on $\spec_{\mathrm{op}}(T^\pm)$. Continuous
functional calculus in $\mathcal A$ gives
\[
P^\pm=h_\pm(T^\pm)\in\mathcal A.
\]
Thus $\pi(P^\pm)$ is well-defined.

Set
\[
D:=J-Q_+T^+Q_+-Q_-T^-Q_-.
\]
We prove $D\in\mathcal I$.

Expand $J$ using $Q_++Q_-+Q_Y=\Id$ on $\ell^2(X)$. Every term in this expansion containing
$Q_Y$ lies in $\mathcal I$ by Proposition~\ref{prop:quotient_split}(i). The same is true for the
cross terms $Q_+JQ_-$ and $Q_-JQ_+$.

It remains to treat the diagonal difference terms
\[
Q_+(J-T^+)Q_+,
\qquad
Q_-(J-T^-)Q_-.
\]
Let
\[
S_+:=\mathrm{prop}(J-T^+).
\]
Suppose
\[
\langle x,Q_+(J-T^+)Q_+\,y\rangle\neq0.
\]
Then $x,y\in Z^+$ and $d(x,y)\le S_+$. If $x\notin N_{R_{\mathrm{ag}}+S_+}(Y)$, then
\[
d(x,Y)>R_{\mathrm{ag}}+S_+,
\]
so
\[
d(y,Y)\ge d(x,Y)-d(x,y)>R_{\mathrm{ag}}.
\]
Hence both $x$ and $y$ lie in $X^+\setminus N_{R_{\mathrm{ag}}}(Y)$, and the assumed sidewise
agreement forces the matrix coefficient to vanish, a contradiction. Thus every nonzero matrix
coefficient of $Q_+(J-T^+)Q_+$ has
\[
x\in N_{R_{\mathrm{ag}}+S_+}(Y).
\]
By symmetry the same holds for $y$. Therefore
\[
Q_+(J-T^+)Q_+\in\mathcal I.
\]
The same argument gives
\[
Q_-(J-T^-)Q_-\in\mathcal I.
\]
Hence $D\in\mathcal I$.

Applying $\pi$ gives
\[
\pi(J)=\pi(Q_+T^+Q_+)+\pi(Q_-T^-Q_-).
\]
By Proposition~\ref{prop:quotient_split},
\[
\pi(Q_+T^+Q_+)=E_+(\pi(T^+)),
\qquad
\pi(Q_-T^-Q_-)=E_-(\pi(T^-)),
\]
which proves (i).

Now fix $0<\gamma<\min\{\gamma_+,\gamma_-\}$, choose
$B\ge\max\{\opnorm{T^+},\opnorm{T^-}\}$, and let $h$ be as in part~(ii).
Because
\[
\mathcal A/\mathcal I
=
E_+(\mathcal A/\mathcal I)\oplus E_-(\mathcal A/\mathcal I),
\]
continuous functional calculus is componentwise on this direct sum. Therefore part~(i) implies
\[
h(\pi(J))
=
E_+\!\bigl(h(\pi(T^+))\bigr)
+
E_-\!\bigl(h(\pi(T^-))\bigr).
\]
Since $\pi$ is a $*$-homomorphism, it commutes with continuous functional calculus:
\[
h(\pi(T^\pm))=\pi(h(T^\pm)).
\]
On the operator spectrum of $T^\pm$, the function $h$ agrees with $\chi_{(-\infty,0)}$ because
$\spec_{\mathrm{op}}(T^\pm)\subset[-B,B]$ and
$\spec_{\mathrm{op}}(T^\pm)\cap(-\gamma,\gamma)=\emptyset$. Hence the operator functional
calculus on $\ell^2(X)$ gives
\[
h(T^\pm)=\chi_{(-\infty,0)}(T^\pm)=P^\pm.
\]
Consequently,
\[
h(\pi(J))
=
E_+(\pi(P^+))+E_-(\pi(P^-)),
\]
which proves (ii).

For part~(iii), each $\pi(P^\pm)$ is a projection, $E_\pm$ are $*$-homomorphisms, and their
ranges are orthogonal ideals. Therefore
\[
p_J^2
=
E_+(\pi(P^+))^2+E_-(\pi(P^-))^2
=
p_J,
\qquad
p_J^*=p_J.
\]
So $p_J$ is a projection. The identity $p_J=h(\pi(J))=\pi(h(J))$ follows from part~(ii) and the
continuous functional calculus. Since the right-hand side is the same for every admissible cutoff
$h$, the projection $p_J$ is independent of the choice of $\gamma$ and $h$.
\end{proof}

\subsection{Proof of Result~\ref{prop:boundary_projection}}\label{app:proof:prop:boundary_projection}
\begin{proof}[Proof of \Cref{prop:boundary_projection}]
Because $\widetilde p$ is a self-adjoint lift of the projection $p$,
\[
\pi\!\bigl(\exp(2\pi i\,\widetilde p)\bigr)=\exp(2\pi i\,p)=\Id.
\]
Hence
\[
\exp(2\pi i\,\widetilde p)-\Id \in M_n(\mathcal I),
\]
so $\exp(2\pi i\,\widetilde p)$ is a unitary in $\Id+M_n(\mathcal I)$ and therefore defines a class
in $K_1(\mathcal I)$.

If $\widetilde p_0$ and $\widetilde p_1$ are two self-adjoint lifts of the same projection $p$, then
\[
\widetilde p_t:=(1-t)\widetilde p_0+t\widetilde p_1
\qquad (0\le t\le 1)
\]
is a norm-continuous path of self-adjoint lifts of $p$. Therefore
\[
t\longmapsto \exp(2\pi i\,\widetilde p_t)
\]
is a homotopy of unitaries in $\Id+M_n(\mathcal I)$, so the $K_1$-class is independent of the
chosen lift.

For the last statement, because $\mathcal A$ is unital, $h(a)\in M_n(\mathcal A)$ for every
continuous $h$. Functional calculus in $\mathcal A/\mathcal I$ gives
\[
\pi(h(a))=h(\pi(a)).
\]
Since $h$ agrees with $\chi_{(-\infty,0)}$ on $\spec(\pi(a))$, the element $h(\pi(a))$ is a
projection. Applying the first part with lift $h(a)$ proves
\[
\partial_Y\!\bigl(\pi(h(a))\bigr)
=
\bigl[\exp(2\pi i\,h(a))\bigr].
\]
\end{proof}

\subsection{Proof of Result~\ref{prop:positive_boundary_homotopy}}\label{app:proof:prop:positive_boundary_homotopy}
\begin{proof}[Proof of \Cref{prop:positive_boundary_homotopy}]
The map
\[
s\longmapsto E_+(\pi(P_s))
\]
is a norm-continuous path of projections in $\mathcal A/\mathcal I$, because $\pi$ and $E_+$ are
bounded $*$-homomorphisms. Hence the $K_0(\mathcal A/\mathcal I)$-class of
$E_+(\pi(P_s))$ is constant. Applying the connecting map
\[
\partial_Y:K_0(\mathcal A/\mathcal I)\longrightarrow K_1(\mathcal I)
\]
gives the claim.
\end{proof}

\subsection{Proof of Result~\ref{prop:sidewise_data_interface_invariance}}\label{app:proof:prop:interface_supported_invariance}\label{app:proof:prop:sidewise_data_interface_invariance}
\begin{proof}[Proof of \Cref{prop:sidewise_data_interface_invariance}]
Since $V,V^\pm\in\mathcal I$, the quotient images satisfy
\[
\pi(J')=\pi(J),
\qquad
\pi(T^{\prime\pm})=\pi(T^\pm).
\]
Let
\[
P^\pm:=\chi_{(-\infty,0)}(T^\pm),
\qquad
P^{\prime\pm}:=\chi_{(-\infty,0)}(T^{\prime\pm}).
\]
Choose continuous compactly supported cutoffs $h^\pm$ and $h^{\prime\pm}$ agreeing with
$\chi_{(-\infty,0)}$ on the operator spectra of $T^\pm$ and $T^{\prime\pm}$, respectively. Since
\[
\spec(\pi(T^\pm))=\spec(\pi(T^{\prime\pm}))
\]
and this quotient spectrum is contained in both operator spectra, both cutoffs give the same
functional calculus element in the quotient:
\[
\pi(P^{\prime\pm})
=
h^{\prime\pm}(\pi(T^{\prime\pm}))
=
h^{\prime\pm}(\pi(T^\pm))
=
h^\pm(\pi(T^\pm))
=
\pi(P^\pm).
\]
Therefore
\[
p_{J'}
=
E_+(\pi(P^{\prime +}))+E_-(\pi(P^{\prime -}))
=
E_+(\pi(P^+))+E_-(\pi(P^-))
=p_J.
\]
Applying $\partial_Y$ gives equality of the interface classes.
\end{proof}

\subsection{Proof of Result~\ref{prop:finite_component_stability}}\label{app:proof:prop:finite_component_stability}
\begin{proof}[Proof of \Cref{prop:finite_component_stability}]
Let
\[
F:=\bigsqcup_{n=1}^N X_n.
\]
This is finite. Since each $Y_n$ is nonempty,
\[
R_F:=\max_{x\in F}\dist(x,Y)
\]
is finite, so $F\subset N_{R_F}(Y)$. An operator supported on $F\times F$ is therefore supported
inside \(N_{R_F}(Y)\times N_{R_F}(Y)\). Since $V$ is bounded and finite-propagation, it
belongs to the algebraic interface ideal, hence to $C^*(Y\subset X)$. Applying this observation to
each finite-component change in the sidewise data, the final statement follows from
Proposition~\ref{prop:sidewise_data_interface_invariance}.
\end{proof}

\subsection{Proof of Result~\ref{prop:norm_vanishing_component_defects}}\label{app:proof:prop:norm_vanishing_component_defects}
\begin{proof}[Proof of \Cref{prop:norm_vanishing_component_defects}]
Let \(V^{(N)}:=\bigoplus_{n\le N}V_n\). This operator is supported on the finite set
\(\bigsqcup_{n\le N}X_n\). Since each \(Y_n\) is nonempty, the proof of
Proposition~\ref{prop:finite_component_stability} shows that \(V^{(N)}\in C^*(Y\subset X)\). Also
\[
\|V-V^{(N)}\|=\sup_{n>N}\|V_n\|\longrightarrow0.
\]
The interface ideal is norm closed, so \(V\in C^*(Y\subset X)\). The final statement is
Proposition~\ref{prop:sidewise_data_interface_invariance}.
\end{proof}

\subsection{Proof of Result~\ref{prop:sidewise_homotopy}}\label{app:proof:prop:sidewise_homotopy}
\begin{proof}[Proof of \Cref{prop:sidewise_homotopy}]
Let
\[
B:=\sup_{s\in[0,1]}\max\{\|T_s^+\|,\|T_s^-\|\}<\infty.
\]
Choose one compactly supported continuous cutoff $h$ that equals $1$ on
$[-B,-\gamma/2]$ and $0$ on $[\gamma/2,B]$. Then $h$ agrees with
$\chi_{(-\infty,0)}$ on the operator spectra of all $T_s^\pm$, and
\[
P_s^\pm:=h(T_s^\pm)
\]
are norm-continuous paths of projections in $\mathcal A$. Hence
\[
p_s:=E_+(\pi(P_s^+))+E_-(\pi(P_s^-))
\]
is a norm-continuous path of projections in $\mathcal A/\mathcal I$, so its $K_0$-class is
constant. Applying the boundary map
\[
\partial_Y:K_0(\mathcal A/\mathcal I)\to K_1(\mathcal I)
\]
shows that the interface class is independent of $s$.
\end{proof}

\subsection{Proof of Result~\ref{cor:copy_potential_independence}}\label{app:proof:cor:copy_potential_independence}
\begin{proof}[Proof of \Cref{cor:copy_potential_independence}]
The positive-side comparison is fixed and has the uniform bulk gap from Assumption~\ref{ass:uniform_bulk}. On the copy sector,
\[
H_{\pm,\mathrm{copy}}^\infty+V_s\ge (M_0-B_{\mathrm{copy},-})\Id,
\]
while on the original sector $T_s^-$ equals $\Id$. Hence the negative-side comparison has a
uniform positive gap at zero. The collar-locality argument is unchanged: away from the glued
boundary, $J_s$ agrees with $T_s^+$ on the positive side and with $T_s^-$ on the negative side.
Thus $(J_s;T_s^+,T_s^-)$ is a norm-continuous family of sidewise class-A interfaces with a uniform
comparison gap, and Proposition~\ref{prop:sidewise_homotopy} applies.
\end{proof}

\subsection{Proof of Result~\ref{thm:general_copy_interface}}\label{app:proof:thm:general_copy_interface}
\begin{proof}[Proof of \Cref{thm:general_copy_interface}]
The operators are bounded, self-adjoint, finite-propagation Roe operators. The positive comparison
$T^+=H_\pm^\infty$ has the bulk gap and Fermi projection $P_\pm$ by
Proposition~\ref{prop:bulk_class}. On the original sector, $T_V^-$ is the identity. On the copy
sector,
\[
H_{\pm,\mathrm{copy}}^\infty+V\ge (M_0-B_{\mathrm{copy},-})\Id>0,
\]
so
\[
\dist\bigl(0,\spec_{\mathrm{op}}(T_V^-)\bigr)\ge \min\{1,M_0-B_{\mathrm{copy},-}\}.
\]
Thus $T_V^-$ is strictly positive and its negative Fermi projection is zero.

It remains to check sidewise agreement. On the positive side outside $Y$ we are strictly in the
original sector, so $V$ vanishes and $J_V=T^+$. On the negative side outside $Y$, the original
sector projection vanishes, and Proposition~\ref{prop:collar_locality_finite} kills the
original--copy cross block. Thus $J_V=T_V^-$ there. Therefore Proposition~\ref{prop:quotient_sidewise}
applies, and since the negative-side Fermi projection is zero it gives
\[
p_{J_V}=E_+(\pi(P_\pm)).
\]
The formula for the interface class is Definition~\ref{def:interface_class}.
\end{proof}

\subsection{Proof of Result~\ref{thm:bulkboundary_main}}\label{app:proof:thm:bulkboundary_main}
\begin{proof}[Proof of \Cref{thm:bulkboundary_main}]
The operators $J_{\pm,M}^\infty$, $T^+=H_\pm^\infty$, and $T^-=\widetilde T^-_{\pm,M}$ are bounded
self-adjoint local operators on the scalar bounded-geometry space $X$, hence belong to
$\mathcal A=C^*(X)$.

By Proposition~\ref{prop:bulk_class},
\[
T^+=H_\pm^\infty
\]
is bounded, self-adjoint, local, has an operator gap at $0$, and has operator Fermi projection
\[
P^+=P_\pm\in C^*(X).
\]
By Proposition~\ref{prop:trivial_jam},
\[
T^-=\widetilde T^-_{\pm,M}
\]
is bounded, self-adjoint, local, and has
\[
\dist\bigl(0,\spec_{\mathrm{op}}(T^-)\bigr)\ge \min\{1,M-B_{\mathrm{copy},-}\}
\]
when $M>B_{\mathrm{copy},-}$. Therefore
\[
P^-:=\chi_{(-\infty,0)}(T^-)=0.
\]

It remains to verify the sidewise agreement conditions. On the positive side outside $Y$ we are
strictly on the original side, so the jam term vanishes and
\[
\chi_{X^+\setminus Y}\bigl(J_{\pm,M}^\infty-T^+\bigr)\chi_{X^+\setminus Y}=0.
\]
On the negative side outside $Y$ we are strictly on the copy side. There $P_{\mathrm{orig}}^\infty$
vanishes, and Proposition~\ref{prop:collar_locality_finite} shows that all original--copy matrix
coefficients vanish there. Hence
\[
\chi_{X^-\setminus Y}\bigl(J_{\pm,M}^\infty-T^-\bigr)\chi_{X^-\setminus Y}=0.
\]
Thus $(J_{\pm,M}^\infty;T^+,T^-)$ is a sidewise class-A interface along $Y$.

Now Proposition~\ref{prop:quotient_sidewise} applies. Since $P^-=0$, it yields
\[
p_{J_{\pm,M}^\infty}
=
E_+\!\bigl(\pi(P^+)\bigr)+E_-\!\bigl(\pi(P^-)\bigr)
=
E_+\!\bigl(\pi(P_\pm)\bigr),
\]
which is \eqref{eq:quotient_fermi_projection}. The equivalent cutoff statement follows from the
same proposition.

Finally, Definition~\ref{def:interface_class} gives
\[
\Indint(J_{\pm,M}^\infty;T^+,T^-)
=
\partial_Y\!\bigl(p_{J_{\pm,M}^\infty}\bigr)
=
\partial_Y\!\bigl(E_+(\pi(P_\pm))\bigr),
\]
which is \eqref{eq:main_interface_class}. The statement about the jammed side is exactly the
vanishing of $P^-$ and hence of the negative-side quotient contribution.
\end{proof}

\subsection{Proof of Result~\ref{cor:jam_independence}}\label{app:proof:cor:jam_independence}
\begin{proof}[Proof of \Cref{cor:jam_independence}]
The right-hand side depends only on the bulk projection $P_\pm$ and the positive quotient
endomorphism $E_+$, not on $M$. Apply Theorem~\ref{thm:bulkboundary_main}.
\end{proof}

\subsection{Proof of Result~\ref{thm:bulk_chern}}\label{app:proof:thm:bulk_chern}
\begin{proof}[Proof of \Cref{thm:bulk_chern}]
Each $H_{\pm,n}=(B_n+B_n^\dagger)\pm S_n$ is the closed-case Higson--Prodan Hamiltonian on $D(\K_n)$,
so the direct-sum family $H_\pm^\infty$ is the Higson--Prodan family for the refinement sequence of
closed doubles. The bulk projection $P_\pm$ constructed in Proposition~\ref{prop:bulk_class} is
therefore the Higson--Prodan Fermi projection, and the Chern-number computation is the imported
content of the Higson--Prodan construction and its Higson--Roe analytic-surgery input
\cite{HigsonProdanUniversalPRL,HigsonProdanUniversalArxiv,HigsonRoeI,HigsonRoeII,HigsonRoeIII}.
\end{proof}

\subsection{Proof of Result~\ref{cor:interface_from_chern}}\label{app:proof:cor:interface_from_chern}
\begin{proof}[Proof of \Cref{cor:interface_from_chern}]
The identity is Theorem~\ref{thm:bulkboundary_main} under the additional Chern-number input
recorded in Theorem~\ref{thm:bulk_chern}. The independence of the large scalar jam scale is
Corollary~\ref{cor:jam_independence}. This gives exactly the displayed formula and the stated
identification of the bulk projection.
\end{proof}

\subsection{Proof of Result~\ref{cor:hp_admissible_interface}}\label{app:proof:cor:hp_admissible_interface}
\begin{proof}[Proof of \Cref{cor:hp_admissible_interface}]
Theorem~\ref{thm:hp_closed_input} gives Assumptions~\ref{ass:S_uniform}
and~\ref{ass:uniform_bulk}. The interface-class formula is then
Theorem~\ref{thm:bulkboundary_main}, and the Chern-number statement is the Higson--Prodan bulk
theorem as recorded in Theorem~\ref{thm:bulk_chern}. For the original/copy partition, the winding
statement is Theorem~\ref{M-thm:pathb_roe_winding}.
\end{proof}

\subsection{Proof of Result~\ref{prop:hp_em_criterion}}\label{app:proof:prop:hp_em_criterion}
\begin{proof}[Proof of \Cref{prop:hp_em_criterion}]
The closed-bulk input gives the positive-side model and its Chern pairing.
Proposition~\ref{prop:coarse_excision} proves that the original/copy decomposition gives a coarse
partition with interface \(Y\). Corollary~\ref{cor:collar_locality_coarse} gives sidewise agreement
away from \(Y\): the positive side agrees with the closed Higson--Prodan model and the copy side
agrees with the jammed comparison model. Proposition~\ref{prop:trivial_jam} shows that this
comparison model is trivial on the negative side once \(M>B_{\mathrm{copy},-}\). These are the
geometric and sidewise pieces of Definition~\ref{def:hp_em_admissible}; the Ewert--Meyer/Roe
identification is Theorem~\ref{M-thm:pathb_roe_winding}. None of these steps uses a boundary gap.
\end{proof}

\subsection{Proof of Result~\ref{cor:concrete_hp_em_classes}}\label{app:proof:cor:concrete_hp_em_classes}
\begin{proof}[Proof of \Cref{cor:concrete_hp_em_classes}]
The double of a full-boundary triangulated oriented surface is a closed oriented triangulated
surface by Proposition~\ref{M-prop:double_closed}. The listed refinement schemes have uniformly
bounded incidence under the stated hypotheses, and doubling preserves that bounded geometry. The
HP-covered condition supplies the controlled cap-product operator, the closed-bulk gap, and the
Chern pairing \(\Ch([P_\pm])=\pm1\). Proposition~\ref{prop:coarse_excision} supplies the
original/copy coarse partition, Proposition~\ref{prop:hp_em_criterion} records the Ewert--Meyer
convention, and Theorem~\ref{M-thm:pathb_roe_winding} gives the \(K_1\)-class identity.
\end{proof}

\subsection{Proof of Result~\ref{cor:tail_admissibility}}\label{app:proof:cor:tail_admissibility}
\begin{proof}[Proof of \Cref{cor:tail_admissibility}]
Changing or deleting finitely many components gives an operator supported on a finite subset of
the coarse disjoint union. Since each interface component is nonempty, such an operator belongs to
the interface ideal by Proposition~\ref{prop:finite_component_stability}. The sidewise interface
\(K_1\)-class is therefore unchanged by Proposition~\ref{prop:sidewise_data_interface_invariance}.
In the straight periodic Toeplitz case (\Cref{M-cor:toeplitz_winding}) the integer winding,
as a \(K_1\)-homomorphism, has the same value on the unchanged \(K_1\)-class.
\end{proof}

\subsection{Proof of Result~\ref{thm:full_hp_bec}}\label{app:proof:thm:full_hp_bec}
\begin{proof}[Proof of \Cref{thm:full_hp_bec}]
The Higson--Prodan closed-bulk theorem gives Assumptions~\ref{ass:S_uniform} and
\ref{ass:uniform_bulk} through Theorem~\ref{thm:hp_closed_input}. It also identifies \(P_\pm\) as
the Higson--Prodan Fermi projection with \(\Ch([P_\pm])=\pm1\). Theorem~\ref{thm:bulkboundary_main}
gives the sidewise interface realization and the formula
\[
\Indint(J_{\pm,M}^\infty;T^+,T^-)
=
\partial_Y(E_+(\pi(P_\pm))).
\]
The negative comparison operator has no negative-energy Fermi projection by
Proposition~\ref{prop:trivial_jam}, so the negative side contributes Chern number \(0\). The
\(K_1\)-class identity is therefore unconditional modulo HP. In the straight periodic
Toeplitz case (\Cref{M-cor:toeplitz_winding}), the Toeplitz integer pairing applied to this
class yields \(\pm1\) internally. For general partitions the integer pairing is recorded in
\Cref{M-cor:conditional_winding}, conditional on the existence of an integer-valued
Chern-jump pairing on \(K_1(\mathcal I)\). Independence of \(M\)
is Corollary~\ref{cor:jam_independence}.
\end{proof}

\subsection{Proof of Result~\ref{cor:relabel_orientation}}\label{app:proof:cor:relabel_orientation}
\begin{proof}[Proof of \Cref{cor:relabel_orientation}]
An orientation-preserving boundary-compatible simplicial isomorphism induces a unitary
identification of the simplex Hilbert spaces. It conjugates the boundary operator, cap-product
operator, original/copy projections, collar ideals, and quotient maps, so the corresponding
\(K\)-classes and pairings are unchanged. Reversing orientation changes the sign of the fundamental
and relative fundamental classes. Hence the Poincar\'e-duality operator \(S\) changes sign, up to
the same unitary relabeling, and \(H_+=(B+B^\dagger)+S\) is identified with the old \(H_-\). The
Higson--Prodan Chern number and the compatible interface winding therefore change sign.
\end{proof}

\subsection{Proof of Result~\ref{cor:interface_winding_robustness}}\label{app:proof:cor:interface_winding_robustness}
\begin{proof}[Proof of \Cref{cor:interface_winding_robustness}]
The scalar jam-scale statement is Corollary~\ref{cor:jam_independence}. Positive copy-side
confinements are covered by Theorem~\ref{thm:general_copy_interface} and
Corollary~\ref{cor:copy_potential_independence}. Interface-supported changes are
Proposition~\ref{prop:sidewise_data_interface_invariance}, finite-component changes are
Proposition~\ref{prop:finite_component_stability}, norm-vanishing component defects are
Proposition~\ref{prop:norm_vanishing_component_defects}, and norm-continuous sidewise homotopies are
Proposition~\ref{prop:sidewise_homotopy}. Each operation preserves the \(K_1\)-class
\(\partial_Y(E_+(\pi(P_\pm)))\). In the straight periodic Toeplitz case
(\Cref{M-cor:toeplitz_winding}) the integer winding, as a \(K_1\)-homomorphism, takes the
same value on the preserved \(K_1\)-class.
\end{proof}

\subsection{Appendix proofs for Section~\ref{M-sec:pathb_internalization}}\label{app:pathb_proofs}

\subsection{Proof of Result~\ref{lem:pathb_hp_operator_bounds}}\label{app:proof:lem:pathb_hp_operator_bounds}
\begin{proof}[Proof of \Cref{lem:pathb_hp_operator_bounds}]
The construction of \(S_K\) is the Higson--Prodan cap-product construction
\cite{HigsonProdanUniversalPRL,HigsonProdanUniversalArxiv}. The boundary matrix \(B\) has entries
only between incident simplices, so \(B+B^\dagger\) is self-adjoint and has propagation \(1\) in the
barycentric graph. The operator \(S_K\) is Hermitian by construction and is supported on the
uniformly bounded barycentric neighborhood prescribed by the controlled cap-product formula.
Therefore \(H_\pm(K)\) is self-adjoint and finite propagation.

For a bounded-geometry family, the number of nonzero entries in each row of \(B\), \(B^\dagger\),
and \(S_K\) is uniformly bounded, and the matrix coefficients belong to a fixed finite set of
incidence signs and cap-product coefficients. The Schur test gives a uniform norm bound for each
summand, hence for \(H_\pm(K)\).
\end{proof}

\subsection{Proof of Result~\ref{M-lem:pathb_hp_gap}}\label{app:proof:lem:pathb_hp_gap}
\begin{proof}[Proof (citation)]
This is the closed universal-model theorem of Higson--Prodan,
\emph{Universal Chern Model on Arbitrary Triangulations}, Phys. Rev. Lett. \textbf{136} (2026),
096602, and arXiv:2510.20862 \cite{HigsonProdanUniversalPRL,HigsonProdanUniversalArxiv}. Their
theorem is stated for the same universal Hamiltonians built from the simplicial boundary operator
and the controlled Poincar\'e-duality operator on bounded-geometry closed triangulations. The closed
doubles \(D(\K_n)\) are closed oriented triangulations with the bounded-geometry
control required in that theorem, so its uniform clean Fermi-gap conclusion applies. No part of the
Higson--Prodan gap theorem is reproved here.
\end{proof}

\subsection{Proof of Result~\ref{M-lem:pathb_hp_chern}}\label{app:proof:lem:pathb_hp_chern}
\begin{proof}[Proof (citation)]
The same Higson--Prodan closed universal-model theorem identifies the Fermi projection of
\(H_\pm=(B+B^\dagger)\pm S\) with the nontrivial closed-bulk class and computes its Chern pairing as
\(\pm1\), with the sign fixed by the two choices \(H_+\), \(H_-\) and the orientation convention
\cite{HigsonProdanUniversalPRL,HigsonProdanUniversalArxiv}. The headline theorem uses exactly that
closed-bulk class on the closed doubles. Thus \(P_\pm\) has \(\Ch([P_\pm])=\pm1\). The proof is purely by
citation to the published closed theorem.
\end{proof}

\subsection{Proof of Result~\ref{M-lem:pathb_roe_setup}}\label{app:proof:lem:pathb_roe_setup}
\begin{proof}[Proof of \Cref{M-lem:pathb_roe_setup}]
The identity operator on \(\ell^2(X)\) has propagation \(0\). On a discrete bounded-geometry scalar
Roe space, finite-support compressions of the identity are finite-rank, so the identity is locally
compact and \(C^*(X)\) is unital.

By Lemmas~\ref{lem:pathb_hp_operator_bounds} and~\ref{M-lem:pathb_hp_gap}, the direct sum
\(H_\pm^\infty=\bigoplus_n H_\pm(D(\K_n))\) is a bounded self-adjoint finite-propagation element of
\(C^*(X)\) with a uniform spectral gap about \(0\). Choose \(f\in C_c(\mathbb R)\) that equals
\(1\) on the negative spectral window and \(0\) on the positive spectral window. Functional calculus
in the unital Roe algebra gives
\[
P_\pm=f(H_\pm^\infty)=\chi_{(-\infty,0)}(H_\pm^\infty)\in C^*(X).
\]
This is the same argument as Proposition~\ref{prop:bulk_class}, now with the gap supplied by
Lemma~\ref{M-lem:pathb_hp_gap}.
\end{proof}

\subsection{Proof of Result~\ref{lem:pathb_combes_thomas}}\label{app:proof:lem:pathb_combes_thomas}
\begin{proof}[Proof of \Cref{lem:pathb_combes_thomas}]
Fix a component \(X_n\) and a basepoint \(w\in X_n\). Let \(R\) be a uniform propagation bound for
\(H_{\pm,n}\), let \(B_0\) be a uniform norm bound, and let
\[
W_a\delta_x=e^{a d(x,w)}\delta_x
\]
for \(a>0\). The conjugated operator \(H_a=W_aH_{\pm,n}W_a^{-1}\) has the same matrix support as
\(H_{\pm,n}\), and each matrix coefficient changes by a factor between \(e^{-aR}\) and \(e^{aR}\).
The Schur-test constants from bounded geometry therefore give
\[
\|H_a-H_{\pm,n}\|\le C_0(e^{aR}-1)
\]
with \(C_0\) independent of \(n\). Choose \(a>0\) so that this norm is at most \(g_{\mathrm{HP}}/4\).

Let \(\Gamma\) be a fixed contour enclosing the negative spectral window and staying at distance
at least \(g_{\mathrm{HP}}/2\) from \(\spec(H_{\pm,n})\). The previous estimate and the resolvent
identity imply that, for \(z\in\Gamma\),
\[
\|(H_a-z)^{-1}\|\le 4/g_{\mathrm{HP}}
\]
after decreasing \(a\) if necessary. Since
\[
(H_{\pm,n}-z)^{-1}
=W_a^{-1}(H_a-z)^{-1}W_a,
\]
we obtain, for all \(v,w\in X_n\),
\[
\left|\langle\delta_v,(H_{\pm,n}-z)^{-1}\delta_w\rangle\right|
\le (4/g_{\mathrm{HP}})e^{-a d(v,w)}.
\]
The Riesz formula
\[
P_{\pm,n}=\frac{1}{2\pi i}\int_\Gamma (z-H_{\pm,n})^{-1}\,dz
\]
then gives the stated bound with \(c=a\) and \(C=|\Gamma|(4/g_{\mathrm{HP}})/(2\pi)\).
\end{proof}

\subsection{Proof of Result~\ref{M-lem:pathb_exact_sequence}}\label{app:proof:lem:pathb_exact_sequence}
\begin{proof}[Proof of \Cref{M-lem:pathb_exact_sequence}]
The ideal property is Lemma~\ref{lem:roe_ideal}. Hence the inclusion
\(\mathcal I\hookrightarrow \mathcal A\) and quotient map \(\pi:\mathcal A\to\mathcal A/\mathcal I\)
form a short exact sequence of \(C^*\)-algebras. The six-term exact sequence in \(K\)-theory for a
short exact sequence of \(C^*\)-algebras gives the connecting map
\[
\partial_Y:K_0(\mathcal A/\mathcal I)\to K_1(\mathcal I)
\]
with the sign convention fixed in Proposition~\ref{prop:boundary_projection}; see, for example,
Blackadar's account of \(C^*\)-algebra \(K\)-theory \cite{BlackadarKTheory}. The coarse geometric
meaning of the same extension is Roe's interface ideal construction \cite{RoeLectures}.
\end{proof}

\subsection{Proof of Result~\ref{lem:pathb_boundary_map}}\label{app:proof:lem:pathb_boundary_map}
\begin{proof}[Proof of \Cref{lem:pathb_boundary_map}]
By Definition~\ref{M-def:pathb_interface_projection}, \(h(J_{\pm,M}^\infty)\) is a self-adjoint lift
of the quotient projection \(p_J\). The exponential formula for the even-to-odd connecting map gives
\[
\partial_Y([p_J])
=
\left[\exp(2\pi i\,h(J_{\pm,M}^\infty))\right].
\]
Because \(\pi(h(J_{\pm,M}^\infty))=p_J\), applying \(\pi\) to the exponential gives
\(\exp(2\pi i p_J)=\Id\); hence the exponential differs from the identity by an element of
\(\mathcal I\). If \(h_0\) and \(h_1\) are two admissible cutoffs, the straight-line path
\((1-t)h_0(J_{\pm,M}^\infty)+t h_1(J_{\pm,M}^\infty)\) consists of self-adjoint lifts of the same
projection \(p_J\), because their quotient images are both \(p_J\). Exponentiating gives a homotopy
in \(\Id+\mathcal I\), so the \(K_1\)-class is independent of the cutoff.
\end{proof}

\subsection{Computational details for Theorem~\ref{M-thm:pathb_roe_winding}}\label{app:proof:thm:pathb_roe_winding}
The body proof of Theorem~\ref{M-thm:pathb_roe_winding} is given immediately after the theorem
statement in Section~\ref{M-sec:pathb_internalization}; see the four-step
\emph{Body proof of Theorem~\ref{M-thm:pathb_roe_winding}}. This appendix subsection retains the
explicit computational verifications referenced from the four body steps. It is intended for the
reader who wishes to inspect (i) the cutoff-functional-calculus identity used to lift the quotient
projection to a self-adjoint element of \(\mathcal A\), (ii) the explicit quotient identification of
the comparison-operator pair at finite jam scale, and (iii) the homotopy verification underlying
the admissible-cutoff invariance.
\begin{proof}[Computational verifications for Theorem~\ref{M-thm:pathb_roe_winding}]
\textit{Inputs.}
Lemma~\ref{M-lem:pathb_hp_gap} gives the Higson--Prodan closed-doubled-bulk gap, and
Lemma~\ref{M-lem:pathb_hp_chern} gives \(\Ch([P_\pm])=\pm1\). Lemma~\ref{M-lem:pathb_roe_setup} places
\(P_\pm\) in \(C^*(X)\). These are the inputs to the body proof's Step (3) and
Definition~\ref{M-def:pathb_interface_projection}, recorded here for ease of cross-reference.

\smallskip\textit{Quotient identification of the comparison-operator pair.}
The finite-\(M\) sidewise realization is the one already constructed in
Theorem~\ref{thm:bulkboundary_main}. Its proof uses only collar locality, the quotient splitting,
and the positivity of the auxiliary copy-side model. Specifically, for \(M>B_{\mathrm{copy},-}\),
the negative comparison operator \(\widetilde T^-_{\pm,M}\) is strictly positive by
Proposition~\ref{prop:trivial_jam}, so its negative Fermi projection is \(0\). On the positive side,
the comparison operator is \(T^+=H_\pm^\infty\) with Fermi projection \(P_\pm\). The quotient-side
functional calculus of Proposition~\ref{prop:quotient_sidewise} therefore gives
\[
p_J=E_+(\pi(P_\pm))+E_-(\pi(0))=E_+(\pi(P_\pm)).
\]
This is the explicit form of the quotient identification used in the body proof's Step (4).

\smallskip\textit{Cutoff-functional-calculus identity.}
The same proposition gives \(p_J=\pi(h(J_{\pm,M}^\infty))\) for every admissible cutoff \(h\).
Lemma~\ref{lem:pathb_boundary_map} then identifies the Roe boundary map with the cutoff unitary:
\[
\partial_Y([p_J])
=
[\,\exp(2\pi i\,h(J_{\pm,M}^\infty))\,].
\]
This is exactly the interface class \(\Indint(J_{\pm,M}^\infty;T^+,T^-)\) of
Definition~\ref{def:interface_class}, and is the cutoff-unitary representative used in the body
proof's Step (1).

\smallskip\textit{Ewert--Meyer compatibility for the original/copy partition.}
The original/copy decomposition is coarsely excisive by Proposition~\ref{prop:coarse_excision}, so
it is represented by the same Mayer--Vietoris extension
\[
0\to \mathcal I\to\mathcal A\to\mathcal A/\mathcal I\to0
\]
used to define \(\partial_Y\); see also the top horizontal row of
Figure~\ref{M-fig:pathb_diagram}. Ewert--Meyer's coarse bulk--interface construction pairs the
boundary class of a quotient bulk projection with the interface cocycle for this partition
\cite{EwertMeyer2019}. In the present doubled-jammed model, the quotient splitting
\[
\mathcal A/\mathcal I=E_+(\mathcal A/\mathcal I)\oplus E_-(\mathcal A/\mathcal I)
\]
of Proposition~\ref{prop:quotient_split} identifies the paired class as the Chern jump between the
positive and negative comparison projections. The positive projection is \(P_\pm\), and the
negative projection is \(0\). Therefore
\[
\langle \operatorname{EM}_Y,[p_J]\rangle
=\Ch([P_\pm])-\Ch([0])
=\pm1.
\]
This is the detailed Ewert--Meyer/Roe identification used in the body proof's Step (2); the
abstract naturality argument given in the body uses only that the same six-term exact sequence
underlies both boundary maps with the orientation choice fixed in
Convention~\ref{M-conv:pathb_sign}.

\smallskip\textit{Cutoff homotopy and \(M\)-invariance.}
For two admissible cutoffs \(h_0,h_1\) and the straight-line interpolation
\((1-t)h_0(J_{\pm,M}^\infty)+t h_1(J_{\pm,M}^\infty)\), both endpoints are self-adjoint lifts of
\(p_J\). Their quotient images coincide, so the difference is a self-adjoint element of
\(\mathcal I\); the homotopy
\(t\mapsto\exp(2\pi i((1-t)h_0+th_1)(J_{\pm,M}^\infty))\) takes values in \(\Id+\mathcal I\) and
gives a continuous path of unitaries representing the same \(K_1\)-class. Independence from
\(M\) along admissible jam scales is Corollary~\ref{cor:jam_independence}; this completes the
\(K_1\)-limit step of the body proof's Step (4).
\end{proof}

\subsection{Proof of Result~\ref{M-cor:toeplitz_winding}}\label{app:proof:cor:toeplitz_winding}
\begin{proof}[Proof of \Cref{M-cor:toeplitz_winding}]
Suppose \(Y\subset X\) is the straight periodic interface of the periodic triangular class with
primitive tangential/positive-normal basis \((t,n)\). By Theorem~\ref{M-thm:pathb_roe_winding}
the interface class is \(\partial_Y(E_+(\pi(P_\pm)))\in K_1(\mathcal I)\). Compressing the
\(n\)-direction to a half-space realizes \(\mathcal I\) as the kernel of the canonical Toeplitz
extension recalled in Theorem~\ref{M-thm:periodic_toeplitz_pairing}, whose connecting map
\(\partial_{\mathrm T,t,n}:K_0(C(\mathbb T^2))\to K_1(C(S^1))\cong\mathbb Z\) satisfies
\(\operatorname{wind}\!\left(\partial_{\mathrm T,t,n}([P])\right)=\varepsilon(t,n)\,\Ch(P)\). For
the original/copy partition the Fermi projection \(P_+\) of the positive-side closed-bulk
operator \(H_\pm^\triangle\) is gapped by Theorem~\ref{M-thm:periodic_triangular_bulk}, and the
trivial-jam statement (Proposition~\ref{prop:trivial_jam}) gives \([P_-]=[0]\). The Chern signs
\(\Ch(E_-(H_\pm^\triangle))=\mp1\) of Theorem~\ref{M-thm:periodic_triangular_chern} therefore
determine the integer
\[
\operatorname{wind}\bigl(\partial_Y(E_+(\pi(P_\pm)))\bigr)
=\varepsilon(t,n)\,\Ch(E_-(H_\pm^\triangle))
=\mp\,\varepsilon(t,n)=\pm1.
\]
Every step is internal to this paper: the Toeplitz identification of \(K_1(\mathcal I)\) with
\(\mathbb Z\) and the integer pairing are recalled in Theorem~\ref{M-thm:periodic_toeplitz_pairing}.
\end{proof}

\subsection{Proof of Result~\ref{M-cor:conditional_winding}}\label{app:proof:cor:conditional_winding}
\begin{proof}[Proof of \Cref{M-cor:conditional_winding}]
Assume an integer-valued Chern-jump-compatible pairing
\(\langle\,\cdot\,,\,\cdot\,\rangle:K_0(\mathcal A/\mathcal I)\times K_1(\mathcal I)\to\mathbb Z\)
exists, in the sense made precise in the statement. By
Theorem~\ref{M-thm:pathb_roe_winding} the class \(\partial_Y([p_J])\) coincides with the
\(K_1\)-class \([\exp(2\pi i\,h(J_{\pm,M}^\infty))]\). The hypothesized pairing then evaluates
\(\langle [p_J],\partial_Y([p_J])\rangle=\Ch([P_+])-\Ch([P_-])\). The trivial-jam statement
(Proposition~\ref{prop:trivial_jam}) gives \([P_-]=[0]\), and the closed-bulk Chern part of
Higson--Prodan \cite{HigsonProdanUniversalPRL,HigsonProdanUniversalArxiv} gives
\(\Ch([P_+])=\pm1\). Hence the pairing yields \(\pm1\), conditional on the existence of the
integer-valued Chern-jump pairing on the general partition. In the periodic triangular Toeplitz
case the pairing exists internally (Corollary~\ref{M-cor:toeplitz_winding}); for general
partitions the construction of the integer pairing is not carried out in this paper and is
flagged in the corollary statement.
\end{proof}

\subsection{Proof of Result~\ref{M-cor:pathb_finite_m_winding}}\label{app:proof:cor:pathb_finite_m_winding}
\begin{proof}[Proof of \Cref{M-cor:pathb_finite_m_winding}]
For \(M>B_{\mathrm{copy},-}\) the quotient projection \(p_J=E_+(\pi(P_\pm))\) is independent of
\(M\) by Proposition~\ref{prop:quotient_sidewise}: the jam scale only enters through the
copy-side block, whose negative spectrum is empty above the trivial-jam threshold and whose
quotient image is trivial. The connecting class \(\partial_Y([p_J])\in K_1(\mathcal I)\)
depends only on the quotient projection, hence is independent of \(M\). Cutoff-homotopy
independence is the statement of Corollary~\ref{cor:jam_independence}.
\end{proof}

\subsection{Proof of Result~\ref{M-lem:pathb_collar_stability}}\label{app:proof:lem:pathb_collar_stability}
\begin{proof}[Proof of \Cref{M-lem:pathb_collar_stability}]
Because \(V\in\mathcal I=C^*(Y\subset X)\), the quotient image of the perturbed operator is the same
as the quotient image of the original operator. Thus the quotient projection \(p_J\) is unchanged.
The perturbation is not required to be small in norm; the only spectral requirement is that the
comparison gaps used to choose the sidewise Fermi cutoffs remain open. With those gaps open, the
same cutoff-functional-calculus argument applies and the \(K_1\)-class
\(\partial_Y([p_J])\) is unchanged. In the straight periodic Toeplitz case the integer
winding (\Cref{M-cor:toeplitz_winding}) is a \(K_1\)-homomorphism of this class, hence is
unchanged as well.
\end{proof}

\subsection{Proof of Result~\ref{M-lem:pathb_bg_stability}}\label{app:proof:lem:pathb_bg_stability}
\begin{proof}[Proof of \Cref{M-lem:pathb_bg_stability}]
Let \(H_s\), \(0\le s\le1\), be the deformed closed doubled Hamiltonian path. The preserved uniform
gap gives one continuous cutoff \(f\) with \(P_s=f(H_s)\). Hence \(s\mapsto P_s\) is a
norm-continuous path of projections in \(C^*(X)\). The quotient projections
\[
E_+(\pi(P_s))
\]
form a norm-continuous path in \(\mathcal A/\mathcal I\), so their \(K_0\)-class is constant. Applying
\(\partial_Y\) gives a constant interface \(K_1\)-class. Any \(K_1\)-homomorphism (in
particular the Toeplitz integer pairing in the straight periodic case,
\Cref{M-cor:toeplitz_winding}) takes the same value at the two endpoints.
\end{proof}

\subsection{Proof of Result~\ref{M-lem:pathb_refinement_stability}}\label{app:proof:lem:pathb_refinement_stability}
\begin{proof}[Proof of \Cref{M-lem:pathb_refinement_stability}]
The Higson--Prodan gap bound is uniform along the refinement family, and
\(M_n>B_{\mathrm{copy},-}+1\) gives a uniform positive gap for the auxiliary negative comparison
operators. Passing to a cofinal subsequence changes the coarse disjoint union only by deleting or
coarsely reindexing components. Deleting finitely many components produces an operator supported on
a finite subset of \(X\), hence an element of \(C^*(Y\subset X)\) by
Proposition~\ref{prop:finite_component_stability}. Cofinal reindexing gives the same quotient
projection and therefore the same boundary class. Since \(M_n\to\infty\), the finite-\(M_n\)
correction terms vanish in the Schur-complement estimates, so the limiting interface class is
eventually constant.
\end{proof}

\subsection{Proof of Result~\ref{M-lem:pathb_periodic_triangular}}\label{app:proof:lem:pathb_periodic_triangular}
\begin{proof}[Proof of \Cref{M-lem:pathb_periodic_triangular}]
For the periodic triangular model, Theorem~\ref{M-thm:periodic_triangular_bulk} gives the closed
zero-gap input and Theorem~\ref{M-thm:periodic_triangular_chern} gives the Chern signs directly.
Theorem~\ref{M-thm:pathb_roe_winding} uses the same positive closed-bulk Fermi projection and the
same trivial negative comparison projection. When the interface is primitive and straight,
Theorem~\ref{M-thm:periodic_toeplitz_pairing} identifies the Ewert--Meyer interface cocycle with the
Toeplitz winding. Thus Theorem~\ref{M-thm:pathb_roe_winding} and the internal periodic calculation
evaluate the same \(K_1\)-class and give the same sign.
\end{proof}

\subsection{Proof of Result~\ref{M-lem:pathb_finite_defect}}\label{app:proof:lem:pathb_finite_defect}
\begin{proof}[Proof of \Cref{M-lem:pathb_finite_defect}]
Fixed-collar defects are elements of the interface ideal, so
Lemma~\ref{M-lem:pathb_collar_stability} applies. Defects supported in finitely many coarse
components are also in the interface ideal by Proposition~\ref{prop:finite_component_stability}.
Therefore the quotient projection \(p_J\) and the \(K_1\)-class
\(\partial_Y(E_+(\pi(P_\pm)))\) agree with the finite-defect robustness statements; in the
straight periodic Toeplitz case the integer winding also agrees via
\Cref{M-cor:toeplitz_winding}.
\end{proof}

\subsection{Proof of Result~\ref{M-lem:pathb_hex_square}}\label{app:proof:lem:pathb_hex_square}
\begin{proof}[Proof of \Cref{M-lem:pathb_hex_square}]
A gapped homotopy between a hexagonal or square appendix symbol and the corresponding closed-bulk
class gives a norm-continuous path of Fermi projections. Lemma~\ref{M-lem:pathb_bg_stability} keeps
the \(K_1\)-class of Theorem~\ref{M-thm:pathb_roe_winding} fixed along that path. The endpoint
Chern sign is the one computed in Lemmas~\ref{lem:hex_chern_positive}, \ref{lem:hex_chern_negative},
and \ref{lem:square_chern_positive}, with the relevant orientation convention. In the
straight periodic Toeplitz case the integer winding agrees with the appendix benchmark via
\Cref{M-cor:toeplitz_winding}.
\end{proof}

\subsection{Proof of Result~\ref{M-lem:pathb_tail_annular}}\label{app:proof:lem:pathb_tail_annular}
\begin{proof}[Proof of \Cref{M-lem:pathb_tail_annular}]
The tail-split and annular constructions decompose the interface ideal into the same component
quotient coordinates used by the Toeplitz maps in Appendix~\ref{app:robustness_tail_annular}.
The interface class of Theorem~\ref{M-thm:pathb_roe_winding} is \(\partial_Y(E_+(\pi(P_\pm)))\). Projecting this class to a component quotient
commutes with the boundary map because both operations are induced by quotient homomorphisms of the
same Roe extension. Each component therefore receives exactly the Toeplitz \(K_1\)-coordinate used
in Theorem~\ref{thm:tail_split_component_windings}.
\end{proof}

\section{Numerical experiments and reproducibility}\label{sec:numerics}
%====================================================================

The numerical experiments record the finite matrices used to test the mechanisms in the paper:
doubled bulk gaps on sampled meshes, the intrinsic Poincar\'e--Lefschetz/compression comparison,
boundary localization after jamming, finite-size flux insertion, chirality reversal, collar support,
and the \(M^{-1}\)/\(M^{-2}\) decoupling laws. The HP theorem or the periodic Bloch calculation
supplies the closed-bulk gap in the theorems; the numerics do not define separate
\(K\)-theoretic indices for individual boundary components. The localization, transport, flux,
disorder, and refinement-scaling figures are included to document finite-matrix behavior. They are
checks of the mechanisms proved above, not inputs to the proofs.

\medskip\noindent\textbf{Pipeline and examples.}
The companion code generates boundary triangulations from a subdivided icosphere, verifies the
doubling identities, assembles \(B,S,H_\pm(D(\K))\), forms \(H_{\pm,M}\), and evaluates the
localization, transport, and decoupling quantities below. The reported computations correspond to
the linked repository snapshot
\texttt{1c888c900cdbbfff078f2b55b445ee01f2b4bd17}, accessed 16 May 2026.

\begin{table}[H]
\centering
\small
\caption{Mesh sizes for the three boundary triangulations \(\K\) and their doubles \(D(\K)\).
The counts agree with the identities in Theorem~\ref{M-thm:counts}.}
\label{tab:meshstats}
\IfFileExists{mesh_stats_table.tex}{\begin{tabular}{@{}lrrrrrrrrrrrr@{}}
\toprule
ex & $\#\partial$ & $|\partial\K_0|$ & $|\partial\K_1|$ & $|\K_0|$ & $|\K_1|$ & $|\K_2|$ & $\dim\Hil(\K)$ & $|D_0|$ & $|D_1|$ & $|D_2|$ & $\dim\Hil(D)$ & $g$ \\
\midrule
ex1 & 1 & 22 & 22 & 73 & 194 & 122 & 389 & 124 & 366 & 244 & 734 & $0.5666$ \\
ex2 & 2 & 88 & 88 & 272 & 728 & 456 & 1456 & 456 & 1368 & 912 & 2736 & $0.3380$ \\
ex3 & 3 & 78 & 78 & 329 & 912 & 582 & 1823 & 580 & 1746 & 1164 & 3490 & $0.3272$ \\
\bottomrule
\end{tabular}
}{%
\emph{(File \texttt{mesh\_stats\_table.tex} not found. Run the companion code to generate it.)}}
\end{table}

\begin{table}[H]
\centering
\small
\caption{Selected numerical outputs: the measured doubled bulk half-gap \(g\) for
\(H_+(D(\K))\), the number of jammed eigenvalues selected by \(|E|<g\) at \(M=50\), and fitted
drift slopes when reported.}
\label{tab:numsummary}
\IfFileExists{numerics_summary.tex}{\begin{tabular}{@{}lrrrr@{}}
\toprule
ex & $g$ & \# in-gap ($|E|<g$) & $v_{+}$ & $v_{-}$ \\
\midrule
ex1 & $0.5666$ & 5 & -0.890 & 0.893 \\
ex2 & $0.3380$ & 14 & -- & -- \\
ex3 & $0.3272$ & 12 & -- & -- \\
\bottomrule
\end{tabular}
}{%
\emph{(File \texttt{numerics\_summary.tex} not found. Run the companion code to generate it.)}}
\end{table}

\begin{table}[H]
\centering
\small
\caption{Finite-\(M\) checks. The measured slopes are fitted from the log-log data generated by
the companion script.}
\label{tab:conditionchecks}
\begin{tabular}{@{}lll@{}}
\toprule
Quantity & Theoretical scale & Observed check \\
\midrule
Compressed resolvent error & \(M^{-1}\) & slope \(-0.943\) \\
First-order resolvent error & \(M^{-2}\) & slope \(-1.942\) \\
Mean squared copy leakage & \(M^{-2}\) & slope \(-1.965\) \\
First-order eigenvalue localization & \(M^{-2}\) & slope \(-1.901\) \\
Contour margin \(\delta_H(\Gamma)\) & positive & positive on sampled \(M\)-range \\
Window separation \(g_M(E_{\max})\) & positive for large \(M\) & linear growth after threshold \\
Original--copy coupling support & bounded collar & all nonzero entries inside measured collar \\
\bottomrule
\end{tabular}
\end{table}

\medskip\noindent\textbf{Meshes, doubled gaps, and jammed spectra.}
\Cref{fig:geom} shows the boundary triangulations and \Cref{fig:bulk} shows the sorted spectra of
\(H_+(D(\K))\). In the displayed finite matrices all three doubled models have a positive measured
gap at zero. \Cref{fig:jammed} then shows the spectra after adding \(M\Proj_{\mathrm{copy}}\) with
\(M=50\). The high-energy cluster is copy-supported; the low-energy levels inside \(|E|<g\) are the
finite-size interface candidates.

\begin{figure}[H]
\centering
\begin{subfigure}{0.32\linewidth}
  \centering
  \safeincludegraphics[width=\linewidth]{fig_ex1_mesh.png}
  \caption{\ExOneName.}
\end{subfigure}\hfill
\begin{subfigure}{0.32\linewidth}
  \centering
  \safeincludegraphics[width=\linewidth]{fig_ex2_mesh.png}
  \caption{\ExTwoName.}
\end{subfigure}\hfill
\begin{subfigure}{0.32\linewidth}
  \centering
  \safeincludegraphics[width=\linewidth]{fig_ex3_mesh.png}
  \caption{\ExThreeName.}
\end{subfigure}
\caption{Boundary triangulations used in the experiments. Boundary edges are highlighted.}
\label{fig:geom}
\end{figure}

\begin{figure}[H]
\centering
\begin{subfigure}{0.49\linewidth}
  \centering
  \safeincludegraphics[width=\linewidth]{fig_ex1_bulk_spectrum.png}
  \caption{Example 1.}
\end{subfigure}\hfill
\begin{subfigure}{0.49\linewidth}
  \centering
  \safeincludegraphics[width=\linewidth]{fig_ex2_bulk_spectrum.png}
  \caption{Example 2.}
\end{subfigure}

\vspace{0.5em}
\begin{subfigure}{0.49\linewidth}
  \centering
  \safeincludegraphics[width=\linewidth]{fig_ex3_bulk_spectrum.png}
  \caption{Example 3.}
\end{subfigure}
\caption{Sorted spectra of \(H_+(D(\K))\). Dashed lines mark the measured half-gap
\(\pm g\), where \(g=\min|\spec(H_+(D(\K)))|\).}
\label{fig:bulk}
\end{figure}

\begin{figure}[H]
\centering
\begin{subfigure}{0.49\linewidth}
  \centering
  \safeincludegraphics[width=\linewidth]{fig_ex1_jammed_spectrum.png}
  \caption{Example 1.}
\end{subfigure}\hfill
\begin{subfigure}{0.49\linewidth}
  \centering
  \safeincludegraphics[width=\linewidth]{fig_ex2_jammed_spectrum.png}
  \caption{Example 2.}
\end{subfigure}

\vspace{0.5em}
\begin{subfigure}{0.49\linewidth}
  \centering
  \safeincludegraphics[width=\linewidth]{fig_ex3_jammed_spectrum.png}
  \caption{Example 3.}
\end{subfigure}
\caption{Spectra of \(H_{+,M}\) at \(M=50\). Dashed lines indicate the measured doubled bulk
gap scale inherited from \(H_+(D(\K))\).}
\label{fig:jammed}
\end{figure}

\medskip\noindent\textbf{Boundary localization.}
\Cref{fig:loc} plots the boundary weight \(w_{\partial}(\psi)=\|\Proj_{\partial}\psi\|^2\) for
Example 1. \Cref{fig:component_localization} repeats this check component by component for the
two- and three-boundary examples. The in-gap states have enhanced boundary weight, while the
high-energy copy cluster has negligible original-boundary weight.

\begin{figure}[H]
\centering
\safeincludegraphics[width=0.72\linewidth]{fig_ex1_boundary_localization.png}
\caption{Example 1: boundary weight versus eigenvalue. States selected by \(|E|<g\) have
enhanced boundary weight.}
\label{fig:loc}
\end{figure}

\begin{figure}[H]
\centering
\safeincludegraphics[width=0.86\linewidth]{fig_component_localization_ex23.png}
\caption{Examples 2 and 3: component-resolved boundary weights. The dashed vertical lines mark
the finite-size gap window \(|E|<g\).}
\label{fig:component_localization}
\end{figure}

\medskip\noindent\textbf{Intrinsic boundary comparison.}
\Cref{fig:pl_boundary_comparison,tab:pl_boundary_comparison} compare the intrinsic
Poincar\'e--Lefschetz boundary Hamiltonian \(H^{\mathrm{PL}}_+(\K)\), the doubled original-side
compression \(H^\partial_+(\K)\), and the finite-\(M\) jammed model. This is the numerical check
most directly tied to Theorem~\ref{M-thm:relative_boundary_comparison}: the discrepancy
\(R_{\partial,+}=H^\partial_+(\K)-H^{\mathrm{PL}}_+(\K)\) is visible but confined to the boundary
collar, and the jammed low-energy spectrum tracks the compressed boundary model.

\begin{figure}[H]
\centering
\safeincludegraphics[width=0.92\linewidth]{fig_pl_boundary_comparison.png}
\caption{Example 1: central spectra of \(H^{\mathrm{PL}}_+(\K)\), \(H^\partial_+(\K)\), and
\(H_{+,M}\), together with the measured support distribution of the collar correction
\(R_{\partial,+}\). The correction is not norm-small, but it is boundary-collar supported.}
\label{fig:pl_boundary_comparison}
\end{figure}

\begin{table}[H]
\centering
\small
\caption{Example 1: numerical size and support of
\(R_{\partial,+}=H^\partial_+(\K)-H^{\mathrm{PL}}_+(\K)\).}
\label{tab:pl_boundary_comparison}
\IfFileExists{pl_boundary_comparison_table.tex}{\begin{tabular}{@{}lrrrrrr@{}}
\toprule
example & $\dim\Hil(\K)$ & $\operatorname{rank}R$ & support dim. & support radius & $\|R\|$ & $\|R\|_F$ \\
\midrule
ex1 & 389 & 32 & 44 & 2 & $1.0000$ & $3.7417$ \\
\bottomrule
\end{tabular}
}{%
\emph{(File \texttt{pl\_boundary\_comparison\_table.tex} not found. Run the companion code to generate it.)}}
\end{table}

\medskip\noindent\textbf{Boundary transport and chirality.}
We evolve a wavepacket projected to the finite in-gap subspace and order the boundary vertices
cyclically along each component. Diagonal ridges in \Cref{fig:edgeprop} show one-way boundary
motion. \Cref{fig:chirality} and \Cref{fig:component_drifts} show that the signed drift reverses
under \(H_{+,M}\leftrightarrow H_{-,M}\). These are finite-matrix transport checks; the
global \(K_1\)-class remains the one in Theorem~\ref{M-thm:pathb_roe_winding}.

\begin{figure}[H]
\centering
\begin{minipage}{0.49\linewidth}
  \centering
  \safeincludegraphics[width=\linewidth]{fig_ex1_edge_propagation.png}
  \small Example 1.
\end{minipage}\hfill
\begin{minipage}{0.49\linewidth}
  \centering
  \safeincludegraphics[width=\linewidth]{fig_ex2_edge_propagation.png}
  \small Example 2.
\end{minipage}

\vspace{0.5em}
\begin{minipage}{0.78\linewidth}
  \centering
  \safeincludegraphics[width=\linewidth]{fig_ex3_edge_propagation.png}
  \small Example 3.
\end{minipage}
\caption{Boundary transport heatmaps for wavepackets projected to the finite-size in-gap subspace.}
\label{fig:edgeprop}
\end{figure}

\begin{figure}[H]
\centering
\safeincludegraphics[width=0.68\linewidth]{fig_ex1_chirality_flip.png}
\caption{Example 1: opposite signed drift for \(H_{+,M}\) and \(H_{-,M}\).}
\label{fig:chirality}
\end{figure}

\begin{figure}[H]
\centering
\safeincludegraphics[width=0.86\linewidth]{fig_multiboundary_component_drifts.png}
\caption{Component-resolved signed drift slopes for the two- and three-boundary examples. The
sign reverses under \(H_{+,M}\leftrightarrow H_{-,M}\) on every sampled boundary component.}
\label{fig:component_drifts}
\end{figure}

\medskip\noindent\textbf{Flux insertion.}
For the annular example we also insert a \(U(1)\) twist across a fixed angular cut and recompute
near-zero eigenvalues for \(\theta\in[0,2\pi]\). \Cref{fig:flux_spectral_flow} is a finite-size
spectral-flow check for the same interface class appearing in Theorem~\ref{M-thm:pathb_roe_winding}:
edge branches move through the bulk gap under flux, with the sign-changing events summarized in
\Cref{tab:flux_spectral_flow}. This figure is evidence for the interface mechanism, not a
replacement for the coarse bulk--interface pairing theorem.

\begin{figure}[H]
\centering
\safeincludegraphics[width=0.82\linewidth]{fig_flux_spectral_flow_ex2.png}
\caption{Example 2: finite-size flux-insertion spectral flow for \(H_{+,M}\). The curves are the
computed near-zero eigenvalues as the twist angle runs from \(0\) to \(2\pi\).}
\label{fig:flux_spectral_flow}
\end{figure}

\begin{table}[H]
\centering
\small
\caption{Example 2: flux-insertion check. Sign changes count crossings between consecutive
sampled flux values among the near-zero eigenvalues computed by sparse shift-invert solves.}
\label{tab:flux_spectral_flow}
\IfFileExists{flux_spectral_flow_table.tex}{\begin{tabular}{@{}lrrrr@{}}
\toprule
example & flux samples & eigenvalues/sample & sign changes & $\min |E|$ \\
\midrule
ex2 & 25 & 18 & 4 & $0.0011$ \\
\bottomrule
\end{tabular}
}{%
\emph{(File \texttt{flux\_spectral\_flow\_table.tex} not found. Run the companion code to generate it.)}}
\end{table}

\medskip\noindent\textbf{Finite-\(M\) decoupling checks.}
The next figures test the quantitative estimates in Section~\ref{M-sec:decoupling}. The compressed
resolvent error follows the predicted \(M^{-1}\) behavior from Theorem~\ref{M-thm:resolvent}. The
first-order resolvent correction, squared copy leakage, and first-order eigenvalue localization
follow the predicted \(M^{-2}\) behavior from Theorems~\ref{M-thm:first_order_eff},
\ref{M-thm:eigvec_leak}, and~\ref{M-thm:first_order_spectral_localization}.

\begin{figure}[H]
\centering
\begin{subfigure}{0.49\linewidth}
  \centering
  \safeincludegraphics[width=\linewidth]{fig_resolvent_loglog.png}
  \caption{Compressed resolvent.}
\end{subfigure}\hfill
\begin{subfigure}{0.49\linewidth}
  \centering
  \safeincludegraphics[width=\linewidth]{fig_first_order_resolvent_loglog.png}
  \caption{First-order resolvent.}
\end{subfigure}

\vspace{0.5em}
\begin{subfigure}{0.49\linewidth}
  \centering
  \safeincludegraphics[width=\linewidth]{fig_leakage_loglog.png}
  \caption{Copy leakage.}
\end{subfigure}\hfill
\begin{subfigure}{0.49\linewidth}
  \centering
  \safeincludegraphics[width=\linewidth]{fig_first_order_localization_loglog.png}
  \caption{First-order eigenvalue localization.}
\end{subfigure}
\caption{Log-log finite-\(M\) checks. The fitted slopes are reported in
\Cref{tab:conditionchecks}.}
\label{fig:finiteM_loglog}
\end{figure}

\begin{figure}[H]
\centering
\begin{subfigure}{0.49\linewidth}
  \centering
  \safeincludegraphics[width=\linewidth]{fig_resolvent_convergence.png}
  \caption{Compressed resolvent.}
\end{subfigure}\hfill
\begin{subfigure}{0.49\linewidth}
  \centering
  \safeincludegraphics[width=\linewidth]{fig_ex1_leakage_vs_M.png}
  \caption{Copy leakage.}
\end{subfigure}

\vspace{0.5em}
\begin{subfigure}{0.49\linewidth}
  \centering
  \safeincludegraphics[width=\linewidth]{fig_ex1_eig_localization_vs_M.png}
  \caption{Eigenvalue localization.}
\end{subfigure}
\caption{Example 1 finite-\(M\) decoupling quantities with explicit comparison bounds.}
\label{fig:finiteM_ex1}
\end{figure}

\begin{table}[H]
\centering
\small
\caption{Example 1: numerical eigenvalue-localization summary for several \(M\) values.}
\label{tab:eigloc}
\IfFileExists{eig_localization_summary.tex}{\begin{tabular}{@{}rrrrr@{}}
\toprule
$M$ & $1/M$ & \# in-gap & $\max_{|E|<g}\mathrm{dist}(E,\mathrm{spec}(A))$ & bound \\
\midrule
10.0 & 0.1000 & 4 & $0.1027$ & $0.6591$ \\
20.0 & 0.0500 & 4 & $0.0817$ & $0.2568$ \\
40.0 & 0.0250 & 5 & $0.0409$ & $0.1156$ \\
50.0 & 0.0200 & 5 & $0.0328$ & $0.0907$ \\
100.0 & 0.0100 & 5 & $0.0164$ & $0.0437$ \\
\bottomrule
\end{tabular}
}{%
\emph{(File \texttt{eig\_localization\_summary.tex} not found. Run the companion code to generate it.)}}
\end{table}

\medskip\noindent\textbf{Collar support, contour separation, refinement, and robustness.}
The collar and contour checks in \Cref{fig:collar_contour} test hypotheses used by the contour
projection and interface arguments. The refinement plot in \Cref{fig:refine} records the measured
bulk half-gap across the sampled subdivision levels; the theorem-level gap input comes from HP or
from the periodic proof. \Cref{fig:pl_collar_scaling} records a separate refinement check for the
relative-boundary comparison: the support fraction of \(R_{\partial,+}\) decreases with
subdivision and remains in a radius-two boundary collar in the sampled family. \Cref{fig:disorder}
adds weak onsite disorder and shows that the qualitative boundary transport ridge persists.
\Cref{fig:copy_confinement} tests the robustness of the low-energy boundary sector under a
nonconstant positive copy-side confining potential.

\begin{figure}[H]
\centering
\begin{subfigure}{0.32\linewidth}
  \centering
  \safeincludegraphics[width=\linewidth]{fig_collar_support_histogram.png}
  \caption{Collar support.}
\end{subfigure}\hfill
\begin{subfigure}{0.32\linewidth}
  \centering
  \safeincludegraphics[width=\linewidth]{fig_contour_margin_lower_bound.png}
  \caption{Contour margin.}
\end{subfigure}\hfill
\begin{subfigure}{0.32\linewidth}
  \centering
  \safeincludegraphics[width=\linewidth]{fig_window_separation_thresholds.png}
  \caption{Window separation.}
\end{subfigure}
\caption{Numerical checks of collar support and separation quantities used in the finite-\(M\)
projection estimates.}
\label{fig:collar_contour}
\end{figure}

\begin{figure}[H]
\centering
\safeincludegraphics[width=0.82\linewidth]{fig_refinement_scaling.png}
\caption{Measured bulk half-gap \(g_{\mathrm{gap}}\) versus subdivision level for the three
example families. This finite-mesh check documents the sampled matrices; theorem-level gap input
comes from HP or from the periodic proof.}
\label{fig:refine}
\end{figure}

\begin{figure}[H]
\centering
\safeincludegraphics[width=0.82\linewidth]{fig_pl_collar_scaling.png}
\caption{Example 1 refinement check for the relative-boundary comparison. The measured
support fraction of the collar correction \(R_{\partial,+}\) decreases with refinement, consistent
with the boundary-rank interpretation in Lemma~\ref{M-lem:collar_rank_counts}.}
\label{fig:pl_collar_scaling}
\end{figure}

\begin{table}[H]
\centering
\small
\caption{Example 1: subdivision scaling of the support of
\(R_{\partial,+}=H^\partial_+(\K)-H^{\mathrm{PL}}_+(\K)\).}
\label{tab:pl_collar_scaling}
\IfFileExists{pl_collar_scaling_table.tex}{\begin{tabular}{@{}rrrrrrr@{}}
\toprule
subdiv. & bary rounds & $\dim\Hil(\K)$ & boundary dim. & support dim. & radius & support/dim \\
\midrule
1 & 0 & 77 & 20 & 20 & 2 & $0.2597$ \\
2 & 0 & 389 & 44 & 44 & 2 & $0.1131$ \\
3 & 0 & 1533 & 88 & 114 & 2 & $0.0744$ \\
4 & 0 & 6337 & 180 & 144 & 2 & $0.0227$ \\
\bottomrule
\end{tabular}
}{%
\emph{(File \texttt{pl\_collar\_scaling\_table.tex} not found. Run the companion code to generate it.)}}
\end{table}

\begin{figure}[H]
\centering
\safeincludegraphics[width=0.82\linewidth]{fig_ex1_edge_propagation_disorder.png}
\caption{Example 1 with weak onsite disorder: the qualitative one-way boundary transport remains
visible in the finite in-gap subspace.}
\label{fig:disorder}
\end{figure}

\begin{figure}[H]
\centering
\safeincludegraphics[width=0.88\linewidth]{fig_copy_confinement_robustness.png}
\caption{Example 1: scalar jamming compared with a nonconstant positive copy-side confining
potential. The in-gap count and mean boundary weight are unchanged at the displayed precision,
while the copy weight remains small.}
\label{fig:copy_confinement}
\end{figure}

\begin{table}[H]
\centering
\small
\caption{Example 1: robustness under nonconstant positive copy-side confinement. The
nonconstant potential has entries in the displayed range on the copy sector.}
\label{tab:copy_confinement}
\IfFileExists{copy_confinement_robustness_table.tex}{\begin{tabular}{@{}lrrrr@{}}
\toprule
copy potential & \# in-gap & mean boundary weight & mean copy weight & potential range \\
\midrule
scalar $M$ & 5 & $0.6721$ & $5.982\times 10^{-4}$ & $[50.0,50.0]$ \
nonconstant $V$ & 5 & $0.6722$ & $5.942\times 10^{-4}$ & $[32.7,67.1]$ \
\bottomrule
\end{tabular}
}{%
\emph{(File \texttt{copy\_confinement\_robustness\_table.tex} not found. Run the companion code to generate it.)}}
\end{table}

\subsection{Reproducibility appendix}\label{app:reproducibility_appendix}
The repository includes the rendered figures and summary tables. The companion code is available in
the linked
\href{https://github.com/yspennstate/Doubling-and-Jamming-A-Universal-Chern-Hamiltonian-on-Triangulated-Surfaces-with-Boundary/tree/main}{project repository};
the computations reported here correspond to main-branch snapshot
\texttt{1c888c900cdbbfff078f2b55b445ee01f2b4bd17}, accessed 16 May 2026. The repository URL,
snapshot hash, and included rendered artifacts are the availability record for this manuscript.

The script fixes the random seeds used for disorder and nonconstant copy-side confinement. It also
exports the finite matrices, the barycentric adjacency graph, and the boundary-collar data used
above. The dependency versions used for the local regeneration check were Python~3.14.3,
numpy~2.4.3, scipy~1.17.1, and matplotlib~3.10.8. Matplotlib artifacts are PNG files.

\subsection{Minimal reproduction script}\label{app:minimal_reproduction_script}
After installing the packages listed in \texttt{requirements.txt}, the key periodic triangular
certificate tables, including
\Cref{tab:periodic_triangular_certificate,tab:periodic_triangular_extension_certificate}, are
regenerated with this command:
\begin{verbatim}
python periodic_triangular_bulk_certificate.py --check --write-table
\end{verbatim}
The finite-matrix figure panels and data-backed tables in Appendix~\ref{sec:numerics} are
regenerated with this command:
\begin{verbatim}
python run_companion_numerics_with_robustness.py
\end{verbatim}
The static theorem-map and notation tables are generated by this \LaTeX{} source.

The local source tree expected in the GitHub repository contains the following numerical artifacts:
\begin{verbatim}
periodic_triangular_bulk_certificate.py
periodic_triangular_certificate_table.tex
periodic_triangular_extension_certificate_table.tex
run_companion_numerics_with_robustness.py
requirements.txt
mesh_stats_table.tex
numerics_summary.tex
pl_boundary_comparison_table.tex
pl_collar_scaling_table.tex
flux_spectral_flow_table.tex
eig_localization_summary.tex
copy_confinement_robustness_table.tex
fig_ex1_mesh.png
fig_ex2_mesh.png
fig_ex3_mesh.png
fig_ex1_bulk_spectrum.png
fig_ex2_bulk_spectrum.png
fig_ex3_bulk_spectrum.png
fig_ex1_jammed_spectrum.png
fig_ex2_jammed_spectrum.png
fig_ex3_jammed_spectrum.png
fig_ex1_boundary_localization.png
fig_component_localization_ex23.png
fig_pl_boundary_comparison.png
fig_ex1_edge_propagation.png
fig_ex2_edge_propagation.png
fig_ex3_edge_propagation.png
fig_ex1_chirality_flip.png
fig_multiboundary_component_drifts.png
fig_flux_spectral_flow_ex2.png
fig_resolvent_loglog.png
fig_first_order_resolvent_loglog.png
fig_leakage_loglog.png
fig_first_order_localization_loglog.png
fig_resolvent_convergence.png
fig_ex1_leakage_vs_M.png
fig_ex1_eig_localization_vs_M.png
fig_collar_support_histogram.png
fig_contour_margin_lower_bound.png
fig_window_separation_thresholds.png
fig_refinement_scaling.png
fig_pl_collar_scaling.png
fig_ex1_edge_propagation_disorder.png
fig_copy_confinement_robustness.png
\end{verbatim}



\section*{Acknowledgments}
The research presented in this paper was supported by the European Research Council (ERC) under the European Union's Horizon Europe research and innovation programme (grant agreement No.~101041711), by the Simons Foundation (as part of the Collaboration on the Mathematical and Scientific Foundations of Deep Learning), by the Israel Science Foundation (grant number 2258/19), by the Israel Science Foundation (ISF Grant~4101/25), and by the U.S. National Science Foundation (NSF Grant~OISE-2401227).

\section*{Declaration of generative AI and AI-assisted technologies in the manuscript preparation process}
During the preparation of this work the author used generative AI tools to accelerate drafting
and revision. All mathematical claims, proofs, and bibliographic information were subsequently
reviewed and edited by the author, who takes full responsibility for the content of the manuscript.

\FloatBarrier

\begin{thebibliography}{99}

\bibitem{AsbothOroszlanyPalyi2016}
J.~K.~Asboth, L.~Oroszlany, and A.~Palyi,
\emph{A Short Course on Topological Insulators: Band Structure and Edge States in One and Two
Dimensions},
Lecture Notes in Physics \textbf{919}, Springer, 2016.

\bibitem{BellissardVanElstSchulzBaldes1994}
J.~Bellissard, A.~van Elst, and H.~Schulz-Baldes,
\emph{The noncommutative geometry of the quantum Hall effect},
J. Math. Phys. \textbf{35} (1994), 5373--5451.

\bibitem{Haldane1988}
F.~D.~M.~Haldane,
\emph{Model for a Quantum Hall Effect without Landau Levels: Condensed-Matter Realization of the Parity Anomaly},
Phys. Rev. Lett. \textbf{61} (1988), 2015--2018.

\bibitem{SchnyderRyuFurusakiLudwig}
A.~P.~Schnyder, S.~Ryu, A.~Furusaki, and A.~W.~W.~Ludwig,
\emph{Classification of topological insulators and superconductors in three spatial dimensions},
Phys. Rev. B \textbf{78} (2008), 195125.

\bibitem{KitaevPeriodic}
A.~Kitaev,
\emph{Periodic table for topological insulators and superconductors},
AIP Conf. Proc. \textbf{1134} (2009), 22--30.

\bibitem{QiZhangRMP}
X.-L.~Qi and S.-C.~Zhang,
\emph{Topological insulators and superconductors},
Rev. Mod. Phys. \textbf{83} (2011), 1057--1110.

\bibitem{FukuiHatsugaiSuzuki}
T.~Fukui, Y.~Hatsugai, and H.~Suzuki,
\emph{Chern numbers in discretized Brillouin zone: efficient method of computing (spin) Hall conductances},
J. Phys. Soc. Jpn. \textbf{74} (2005), 1674--1677.

\bibitem{HigsonProdanUniversalPRL}
N.~Higson and E.~Prodan,
\emph{Universal Chern Model on Arbitrary Triangulations},
Phys. Rev. Lett. \textbf{136} (2026), no.~9, 096602,
doi:10.1103/66x5-dvqq.

\bibitem{HigsonProdanUniversalArxiv}
N.~Higson and E.~Prodan,
\emph{A Universal Chern Model on Arbitrary Triangulations},
arXiv:2510.20862, doi:10.48550/arXiv.2510.20862.

\bibitem{HigsonRoeI}
N.~Higson and J.~Roe,
\emph{Mapping Surgery to Analysis I: Analytic Signatures},
K-Theory \textbf{33} (2004), 277--299.

\bibitem{HigsonRoeII}
N.~Higson and J.~Roe,
\emph{Mapping Surgery to Analysis II: Geometric Signatures},
K-Theory \textbf{33} (2004), 301--324.

\bibitem{HigsonRoeIII}
N.~Higson and J.~Roe,
\emph{Mapping Surgery to Analysis III: Exact Sequences},
K-Theory \textbf{33} (2004), 325--346.

\bibitem{WallSurgery}
C.~T.~C.~Wall,
\emph{Surgery on Compact Manifolds},
AMS, 1999.

\bibitem{MunkresElements}
J.~R.~Munkres,
\emph{Elements of Algebraic Topology},
Addison-Wesley, 1984.

\bibitem{Hatcher}
A.~Hatcher,
\emph{Algebraic Topology},
Cambridge University Press, 2002.

\bibitem{KatoPerturbation}
T.~Kato,
\emph{Perturbation Theory for Linear Operators},
Classics in Mathematics, Springer, 1995.

\bibitem{BhatiaMatrixAnalysis}
R.~Bhatia,
\emph{Matrix Analysis},
Graduate Texts in Mathematics, vol.~169, Springer, 1997.

\bibitem{AtiyahKTheory}
M.~F.~Atiyah,
\emph{K-Theory},
W.~A.~Benjamin, 1967.

\bibitem{BlackadarKTheory}
B.~Blackadar,
\emph{K-Theory for Operator Algebras},
2nd ed., Mathematical Sciences Research Institute Publications, vol.~5,
Cambridge University Press, 1998.

\bibitem{BDF1977}
L.~G.~Brown, R.~G.~Douglas, and P.~A.~Fillmore,
\emph{Extensions of $C^*$-algebras and $K$-homology},
Ann. of Math. (2) \textbf{105} (1977), no.~2, 265--324.

\bibitem{ProdanSchulzBaldesBook}
E.~Prodan and H.~Schulz-Baldes,
\emph{Bulk and Boundary Invariants for Complex Topological Insulators: From K-Theory to Physics},
Springer, 2016.

\bibitem{Mitchell2018}
N.~P.~Mitchell, L.~M.~Nash, D.~Hexner, A.~Turner, and W.~T.~M.~Irvine,
\emph{Amorphous topological insulators constructed from random point sets},
Nat. Phys. \textbf{14} (2018), 380--385.

\bibitem{BourneProdan2018}
C.~Bourne and E.~Prodan,
\emph{Non-commutative Chern numbers for generic aperiodic discrete systems},
J. Phys. A: Math. Theor. \textbf{51} (2018), 235202.

\bibitem{BourneMesland2019}
C.~Bourne and B.~Mesland,
\emph{Index theory and topological phases of aperiodic lattices},
Ann. Henri Poincar\'e \textbf{20} (2019), 1969--2038.

\bibitem{ElbauGraf2002}
P.~Elbau and G.~M.~Graf,
\emph{Equality of Bulk and Edge Hall Conductance Revisited},
Comm. Math. Phys. \textbf{229} (2002), 415--432.

\bibitem{ElgartGrafSchenker2005}
A.~Elgart, G.~M.~Graf, and J.~H.~Schenker,
\emph{Equality of the Bulk and Edge Hall Conductances in a Mobility Gap},
Comm. Math. Phys. \textbf{259} (2005), 185--221.

\bibitem{GrafPorta2013}
G.~M.~Graf and M.~Porta,
\emph{Bulk-Edge Correspondence for Two-Dimensional Topological Insulators},
Comm. Math. Phys. \textbf{324} (2013), 851--895.

\bibitem{KubotaControlled}
Y.~Kubota,
\emph{Controlled topological phases and bulk-edge correspondence},
Comm. Math. Phys. \textbf{349} (2017), 493--525.

\bibitem{EwertMeyer2019}
E.~E.~Ewert and R.~Meyer,
\emph{Coarse geometry and topological phases},
Comm. Math. Phys. \textbf{366} (2019), 1069--1098,
doi:10.1007/s00220-019-03303-z.

\bibitem{BourneKellendonkRennie}
C.~Bourne, J.~Kellendonk, and A.~Rennie,
\emph{The $K$-theoretic bulk-edge correspondence for topological insulators},
Ann. Henri Poincar\'e \textbf{18} (2017), 1833--1866.

\bibitem{AlldridgeMaxZirnbauer2020}
A.~Alldridge, C.~Max, and M.~R.~Zirnbauer,
\emph{Bulk-Boundary Correspondence for Disordered Free-Fermion Topological Phases},
Comm. Math. Phys. \textbf{377} (2020), 1761--1821.

\bibitem{LudewigThiangEdge}
M.~Ludewig and G.~C.~Thiang,
\emph{Cobordism invariance of topological edge-following states},
Adv. Theor. Math. Phys. \textbf{26} (2022), no.~3, 673--710.

\bibitem{DrouotZhuEdgeSpectrum}
A.~Drouot and X.~Zhu,
\emph{Topological edge spectrum along curved interfaces},
Int. Math. Res. Not. IMRN \textbf{2024} (2024), no.~22, 13870--13889.

\bibitem{DrouotZhuBEC}
A.~Drouot and X.~Zhu,
\emph{The bulk-edge correspondence for curved interfaces},
arXiv:2408.07950.

\bibitem{RoeLectures}
J.~Roe,
\emph{Lectures on Coarse Geometry},
University Lecture Series \textbf{31}, American Mathematical Society, 2003.

\end{thebibliography}


\end{document}
